\documentclass[11pt,reqno]{amsart}
\usepackage[utf8]{inputenc}
\usepackage[T1]{fontenc}
\usepackage{amsmath,amssymb,amsthm,mathtools,array}
\usepackage{booktabs}
\usepackage[margin=1in]{geometry}
\usepackage[colorlinks=true,linkcolor=blue,citecolor=blue,urlcolor=blue]{hyperref}

\theoremstyle{plain}
\newtheorem{theorem}{Theorem}[section]
\newtheorem{proposition}[theorem]{Proposition}
\newtheorem{lemma}[theorem]{Lemma}
\newtheorem{corollary}[theorem]{Corollary}

\theoremstyle{definition}
\newtheorem{definition}[theorem]{Definition}
\newtheorem{remark}[theorem]{Remark}

% Exact formal verification anchor (safe with underscores and line-breakable via \nolinkurl)
\newcommand{\leanlink}[1]{\allowbreak{\tiny\texttt{[\expandafter\nolinkurl\expandafter{\detokenize{#1}}]}}}
\newcommand{\leanmod}[1]{\allowbreak{\scriptsize\texttt{\expandafter\nolinkurl\expandafter{\detokenize{#1}}}}}
\emergencystretch=3em
\sloppy

% --- Global Symbol Registry Macros (anchored to Lean declarations) ---
\providecommand{\pcHSp}{\operatorname{HSp}} % Formalism.BasicNotions.TH.HSp
\providecommand{\pcHost}{\operatorname{Host}} % ReciprocityTheory.ChartRestriction.G.Geo.Tor.Host
\providecommand{\pcDatum}{\operatorname{Datum}} % ReciprocityTheory.ChartRestriction.G.Geo.Tor.Datum
\providecommand{\pcSpMap}{\operatorname{SpMap}} % ReciprocityTheory.ChartRestriction.G.Geo.Tor.spMap
\providecommand{\pcLevel}{\operatorname{Level}} % ReciprocityTheory.ChartRestriction.G.Geo.Tor.Level
\providecommand{\pcModuli}{\operatorname{Moduli}} % ReciprocityTheory.ChartRestriction.G.Geo.Tor.Moduli
\providecommand{\pcHostA}{\operatorname{Host}} % GlobalTheory.HostModel.host
\providecommand{\pcCoeff}{\operatorname{Coeff}} % ReciprocityTheory.ChartRestriction.G.Geo.Tor.Coeff
\providecommand{\pcOver}{\operatorname{Over}} % GlobalTheory.Ex.ArcF.Over
\providecommand{\pcHostB}{\operatorname{Host}} % Formalism.L2.GlobalTheoryTheory.Arith.host
\providecommand{\pcRank}{\operatorname{Rank}} % Formalism.BasicNotions.HLat.rank
\providecommand{\pcRankA}{\operatorname{Rank}} % ReciprocityTheory.ChartRestriction.G.Out.DegPre.rank
\providecommand{\pcWhat}{\operatorname{What}} % Formalism.L2.Indef.Examples₂.one.what₂
\providecommand{\pcWhatA}{\operatorname{What}} % Formalism.L2.Indef.Examples₂.two.what₂
\providecommand{\pcHom}{\operatorname{Hom}} % GlobalTheory.Serre.Pts.hom
\providecommand{\pcComp}{\operatorname{Comp}} % Formalism.CoboundaryTransfers.Hom.comp
\providecommand{\pcHidaSetting}{\operatorname{HidaSetting}} % Formalism.L2.SG.HidaSetting.HidaSetting
\providecommand{\pcHomA}{\operatorname{Hom}} % GlobalTheory.Serre.Lie.hom
\providecommand{\pcPeriodData}{\operatorname{PeriodData}} % Formalism.WeaklyAdmissibleFiltrations.PeriodData
\providecommand{\pcIsLB}{\operatorname{IsLB}} % ReciprocityTheory.ChartRestriction.G.Semi.LB.IsLB
\providecommand{\pcSizeData}{\operatorname{SizeData}} % Formalism.BasicNotions.SizeData₂
\providecommand{\FourierTransform}{\mathcal{F}} % Formalism.FourierDecayAlt.B.Mat
\providecommand{\pcHostC}{\operatorname{Host}} % GlobalTheory.Mem.LevCh.host
\providecommand{\pcData}{\operatorname{Data}} % ReciprocityTheory.ChartRestriction.G.FD.Data
\providecommand{\pcRanks}{\operatorname{Ranks}} % Formalism.BasicNotions.Ranks
\providecommand{\pcCompA}{\operatorname{Comp}} % Formalism.MainNotions.G.Levels.comp
\providecommand{\pcFinExt}{\operatorname{FinExt}} % Formalism.BasicNotions.Loc.FinExt
\providecommand{\pcFamIn}{\operatorname{FamIn}} % Formalism.L2.DIV.FamIn
\providecommand{\pcInfFam}{\operatorname{InfFam}} % Formalism.IndefR.Datum.InfFam
\providecommand{\pcDegPre}{\operatorname{DegPre}} % ReciprocityTheory.ChartRestriction.G.Out.DegPre
\providecommand{\pcSec}{\operatorname{Sec}} % ReciprocityTheory.ChartRestriction.G.Semi.LB.Sec
\providecommand{\pcLocWach}{\operatorname{LocWach}} % Formalism.BasicNotions.LocWach
\providecommand{\pcRealDat}{\operatorname{RealDat}} % Formalism.IndefR.Links.RealDat
\providecommand{\pcLinComb}{\operatorname{LinComb}} % ReciprocityTheory.ChartRestriction.G.Semi.LB.LinComb
\providecommand{\pcPolAV}{\operatorname{PolAV}} % ReciprocityTheory.ChartRestriction.G.Out.PolAV
\providecommand{\pcFrame}{\operatorname{Frame}} % ReciprocityTheory.ChartRestriction.G.FD.Frame
\providecommand{\pcGen}{\operatorname{Gen}} % ReciprocityTheory.ChartRestriction.G.Semi.LB.Gen
\providecommand{\pcIsGen}{\operatorname{IsGen}} % ReciprocityTheory.ChartRestriction.G.Semi.LB.Gen.IsGen
\providecommand{\pcIsGrpPullback}{\operatorname{IsGrpPullbac}} % ReciprocityTheory.ChartRestriction.G.Geo.Tor.IsGrpPullback
\providecommand{\pcSpc}{\operatorname{Spc}} % GlobalTheory.IndefiniteRowSideFullHecke.WeilTensorModels.Spc
\providecommand{\pcMRing}{\operatorname{MRing}} % ReciprocityTheory.ChartRestriction.G.Geo.MRing
\providecommand{\pcLocLat}{\operatorname{LocLat}} % Formalism.BasicNotions.LocLat
\providecommand{\pcHpk}{\operatorname{Hpk}} % Formalism.QuantumDeterminants.Hpk
\providecommand{\pcQbar}{\operatorname{Qbar}} % Formalism.ConjugateSelfDual.Qbar
\providecommand{\pcModel}{\operatorname{Model}} % ReciprocityTheory.ChartRestriction.G.Geo.Tor.Model
\providecommand{\pcGood}{\operatorname{Good}} % ReciprocityTheory.ChartRestriction.G.Geo.Tor.Good
\providecommand{\pcLevCh}{\operatorname{LevCh}} % Formalism.IndefR.Links.LevCh
\providecommand{\pcPull}{\operatorname{Pull}} % ReciprocityTheory.ChartRestriction.G.Semi.LB.pull
\providecommand{\pcHSpace}{\operatorname{HSpace}} % Formalism.BasicNotions.HSpace
\providecommand{\pcIso}{\operatorname{Iso}} % ReciprocityTheory.ChartRestriction.G.Semi.LB.Iso
\providecommand{\pcEquiv}{\operatorname{Equiv}} % ReciprocityTheory.ChartRestriction.G.Out.QkPt.Equiv
\providecommand{\pcAutInst}{\operatorname{AutInst}} % Formalism.MainNotions.S.AutInst
\providecommand{\pcBase}{\operatorname{Base}} % GlobalTheory.Mem.Base
\providecommand{\pcLocPrime}{\operatorname{LocPrime}} % Formalism.BasicNotions.LocPrime
\providecommand{\pcIsInfFam}{\operatorname{IsInfFam}} % Formalism.L2.Theta.IsInfFam
\providecommand{\pcGridData}{\operatorname{GridData}} % GlobalTheory.ChiGrid.GridData
\providecommand{\pcSAB}{\operatorname{SAB}} % ReciprocityTheory.ChartRestriction.G.Geo.Tor.SAB
\providecommand{\pcModelA}{\operatorname{Model}} % Formalism.MainNotions.S.Loc.Model
\providecommand{\pcIsPlaceOver}{\operatorname{IsPlaceOver}} % Formalism.ConjugateSelfDual.IsPlaceOver
\providecommand{\pcDomWeight}{\operatorname{DomWeight}} % Formalism.BasicNotions.DomWeight
\providecommand{\pcHChar}{\operatorname{HChar}} % Formalism.L2.Host.HChar
\providecommand{\pcIsThetaFunO}{\operatorname{IsThetaFunO}} % Formalism.IndefR.Datum.IsThetaFunO
\providecommand{\pcIncA}{\operatorname{IncA}} % Formalism.L2.Theta.incA
\providecommand{\pcInt}{\operatorname{Int}} % Formalism.MainNotions.G.Levels.Int
\providecommand{\pcDegData}{\operatorname{DegData}} % ReciprocityTheory.ChartRestriction.G.Out.DegData
\providecommand{\pcReal}{\operatorname{Real}} % GlobalTheory.ZH.ThetaSys.Real
\providecommand{\pcLBIso}{\operatorname{LBIso}} % ReciprocityTheory.ChartRestriction.G.Geo.Tor.LBIso
\providecommand{\pcIsAbelianScheme}{\operatorname{IsAbelianSch}} % ReciprocityTheory.ChartRestriction.G.Semi.IsAbelianScheme
\providecommand{\pcHostD}{\operatorname{Host}} % Formalism.BasicNotions.Host
\providecommand{\pcInst}{\operatorname{Inst}} % Formalism.MainNotions.G.Inst
\providecommand{\pcToFun}{\operatorname{ToFun}} % GlobalTheory.Mem.Th.WeilCore.toFun
\providecommand{\pcLetA}{\operatorname{LetA}} % Formalism.BasicNotions.LetA
\providecommand{\pcDivX}{\operatorname{DivX}} % Formalism.BasicNotions.DivX
\providecommand{\pcIsStdChar}{\operatorname{IsStdChar}} % Formalism.L2.Theta.IsStdChar
\providecommand{\pcFrPt}{\operatorname{FrPt}} % Formalism.L2.Pj.FrPt
\providecommand{\pcFlagG}{\operatorname{FlagG}} % Formalism.MainNotions.S.FlagG
\providecommand{\pcGatTies}{\operatorname{GatTies}} % Formalism.BasicNotions.GatTies
\providecommand{\pcWhatB}{\operatorname{What}} % GlobalTheory.Mem.RowsRest.what₂
\providecommand{\pcIsCMExt}{\operatorname{IsCMExt}} % Formalism.MainNotions.GlobalTheoryTheory.GenusMassVolumes.IsCMExt
\providecommand{\pcHeisO}{\operatorname{HeisO}} % Formalism.IndefR.Datum.heisO
\providecommand{\pcDom}{\operatorname{Dom}} % ReciprocityTheory.ChartRestriction.G.FD.Data.dom
\providecommand{\pcLatInst}{\operatorname{LatInst}} % Formalism.MainNotions.G.LatInst
\providecommand{\pcCoef}{\operatorname{Coef}} % UpperTheory.V.HV.Coef
\providecommand{\pcLFlag}{\operatorname{LFlag}} % Formalism.MainNotions.S.Lat.LFlag
\providecommand{\FourierTransformA}{\mathcal{F}} % Formalism.FourierDecayAlt.B.RealPlaceData
\providecommand{\pcPreX}{\operatorname{PreX}} % GlobalTheory.Mem.PreX
\providecommand{\pcSPar}{\operatorname{SPar}} % Formalism.BasicNotions.TH.SPar
\providecommand{\pcPlF}{\operatorname{PlF}} % Formalism.MainNotions.S.LocalSplitGroups.PlF
\providecommand{\pcCompB}{\operatorname{Comp}} % Formalism.L2.Pj.Comp
\providecommand{\pcSpcA}{\operatorname{Spc}} % GlobalTheory.EtC.Spc
\providecommand{\pcExtra}{\operatorname{Extra}} % Formalism.BasicNotions.Extra
\providecommand{\pcObj}{\operatorname{Obj}} % GlobalTheory.Cris.IsMor.obj
\providecommand{\pcCarrier}{\operatorname{Carrier}} % Formalism.BasicNotions.Rows.OLat.carrier
\providecommand{\pcHerm}{\operatorname{Herm}} % Formalism.QuantumDeterminants.herm
\providecommand{\pcNData}{\operatorname{NData}} % Formalism.BasicNotions.NData
\providecommand{\pcMult}{\operatorname{Mult}} % ReciprocityTheory.ChartRestriction.G.Geo.Tor.Datum.mult
\providecommand{\pcShv}{\operatorname{Shv}} % GlobalTheory.EtC.Shv
\providecommand{\pcMCoordIdeal}{\operatorname{MCoordIdeal}} % ReciprocityTheory.ChartRestriction.G.Geo.mCoordIdeal
\providecommand{\pcGfC}{\operatorname{GfC}} % Formalism.L2.Theta.GfC
\providecommand{\pcDSModel}{\operatorname{DSModel}} % Formalism.IndefR.Mult.DSModel
\providecommand{\pcHostE}{\operatorname{Host}} % Formalism.LFunctionChainParameters.ThickSetMeanBounds.CR.host
\providecommand{\pcWchart}{\operatorname{Wchart}} % ReciprocityTheory.ChartRestriction.G.Semi.LB.wchart
\providecommand{\pcObjA}{\operatorname{Obj}} % GlobalTheory.Serre.Summand.obj
\providecommand{\FourierTransformB}{\mathcal{F}} % Formalism.FourierDecayAlt.R.LevelData
\providecommand{\pcEAf}{\operatorname{EAf}} % Formalism.IndefR.Lev.EAf
\providecommand{\pcPrimes}{\operatorname{Primes}} % Formalism.L2.Col.Primes
\providecommand{\pcInclR}{\operatorname{InclR}} % Formalism.MainNotions.GlobalTheoryTheory.GenusMassVolumes.inclR
\providecommand{\pcIncM}{\operatorname{IncM}} % GlobalTheory.IndefiniteRowSideFullHecke.WeilTensorModels.incM
\providecommand{\pcSecA}{\operatorname{Sec}} % GlobalTheory.EtC.Sec
\providecommand{\pcIne}{\operatorname{Ine}} % ReciprocityTheory.Gal.Ine
\providecommand{\pcHomB}{\operatorname{Hom}} % PI.EagonNorthcottZ.hom
\providecommand{\pcESec}{\operatorname{ESec}} % ReciprocityTheory.ChartRestriction.G.Semi.eSec
\providecommand{\pcLevelA}{\operatorname{Level}} % ReciprocityTheory.ChartRestriction.G.Out.Level
\providecommand{\pcMat}{\operatorname{Mat}} % PositivityTheory.RgTw.Mat
\providecommand{\pcCinf}{\operatorname{Cinf}} % GlobalTheory.CuspRep.Cinf
\providecommand{\pcFinPart}{\operatorname{FinPart}} % Formalism.L2.Theta.finPart
\providecommand{\pcHypDec}{\operatorname{HypDec}} % GlobalTheory.MixModel.HypDec
\providecommand{\AdeleRing}{\mathbb{A}} % GlobalTheory.AdelicU.evFin
\providecommand{\pcRes}{\operatorname{Res}} % ReciprocityTheory.ChartRestriction.G.Semi.res
\providecommand{\pcModelB}{\operatorname{Model}} % GlobalTheory.IndefiniteRowSideFullHecke.WeilTensorModels.Model

\title[A Compressed Mathematical Exposition of LFunction eq zero ...]{A Compressed Mathematical Exposition of LFunction eq zero re}
\author{Synthesized by \texttt{proof-compressor}}
\date{\today}

\begin{document}
\begin{abstract}
We present a compressed, self-contained mathematical exposition extracted directly from a formal Lean~4 proof corpus comprising 175,187~modules (15,531,759~lines of code) and 1,061,531~declarations. Exact reverse-dependency pruning isolates 145,276~reachable modules (13,427,197~lines, compression ratio 1.16$\times$) and 706,458~substantive mathematical declarations. Every definition, theorem, and proof citation below is anchored to its verified Lean constant via \texttt{\textbackslash leanlink\{...\}}.
\end{abstract}
\maketitle
\tableofcontents

\section{Proof Overview and Quantitative Architecture}

The formal repository contains 1,061,531 project declarations across 175,187 modules (15,531,759 lines) connected by 9,003,126 dependency edges. Pruning all declarations unreachable from the target theorem(s) eliminates 262,467 dead declarations across 29,911 unused modules, leaving 799,064 reachable declarations (706,458 substantive mathematical results after collapsing routine formalization glue).

\begin{table}[ht]
\centering
\caption{Quantitative Compression Summary}
\begin{tabular}{lr}
\toprule
Metric & Value \\
\midrule
Raw project modules & 175,187 \\
Reachable project modules & 145,276 \\
Raw project LOC & 15,531,759 \\
Reachable project LOC & 13,427,197 \\
Dead-code compression ratio & 1.16$\times$ \\
Reachable declarations & 799,064 \\
Substantive mathematical declarations & 706,458 \\
Canonical duplicate groups & 72,152 \\
Estimated redundant duplicate LOC & 1,958,459 \\
\bottomrule
\end{tabular}
\end{table}

\subsection{Target Main Theorems}
\begin{theorem}[LFunction eq zero re]\leanlink{LFunction_eq_zero_re}
\begin{quote}\small\raggedright$\displaystyle \forall N \ge 1,\; \forall \chi \pmod{N},\; \forall s \in \mathbb{C},\; L(s, \allowbreak \chi) = 0 \land \allowbreak 0 < \operatorname{Re}(s) < 1 \to \allowbreak \operatorname{Re}(s) = 1 / 2$\end{quote}
\end{theorem}
\begin{proof}
The argument operates on the structural objects Dirichlet LZeros On Line \leanlink{DirichletLZerosOnLine}. The claim follows directly from \emph{claim All} \leanlink{claimAll}.
\end{proof}

\begin{theorem}[LFunction eq zero re of is Primitive]\leanlink{LFunction_eq_zero_re_of_isPrimitive}
\begin{quote}\small\raggedright$\displaystyle \forall N \ge 1,\; \forall \chi\;\mathrm{primitive} \pmod{N},\; \forall s \in \mathbb{C},\; L(s, \allowbreak \chi) = 0 \land \allowbreak 0 < \operatorname{Re}(s) < 1 \to \allowbreak \operatorname{Re}(s) = 1 / 2$\end{quote}
\end{theorem}
\begin{proof}
The argument operates on the structural objects Primitive Dirichlet LZeros On Line \leanlink{PrimitiveDirichletLZerosOnLine}. The claim follows directly from \emph{crit All} \leanlink{critAll}.
\end{proof}

\begin{theorem}[LFunction ne zero of half lt re]\leanlink{LFunction_ne_zero_of_half_lt_re}
\begin{quote}\small\raggedright$\displaystyle \forall N \ge 1,\; \forall \chi \pmod{N},\; \forall s \in \mathbb{C},\; \operatorname{Re}(s) > 1 / 2 \to \allowbreak L(s, \allowbreak \chi) \neq 0$\end{quote}
\end{theorem}
\begin{proof}
The argument operates on the structural objects Dirichlet LNonvanishing Right Of Half \leanlink{DirichletLNonvanishingRightOfHalf}. The claim follows directly from \emph{all Chi} \leanlink{allChi}.
\end{proof}

\begin{theorem}[riemann Hypothesis]\leanlink{riemannHypothesis}
\[
\mathrm{RH}
\]
\end{theorem}
\begin{proof}
The claim follows directly from \emph{rh Of All} \leanlink{rhOfAll}.
\end{proof}

\begin{theorem}[riemann Zeta eq zero re]\leanlink{riemannZeta_eq_zero_re}
\begin{quote}\small\raggedright$\displaystyle \forall s \in \mathbb{C},\; \zeta(s) = 0 \land \allowbreak 0 < \operatorname{Re}(s) < 1 \to \allowbreak \operatorname{Re}(s) = 1 / 2$\end{quote}
\end{theorem}
\begin{proof}
The argument operates on the structural objects Riemann Zeta Zeros On Line \leanlink{RiemannZetaZerosOnLine}. The claim follows directly from \emph{rh Of All} \leanlink{rhOfAll}.
\end{proof}

\subsection{Bottleneck Pillar Theorems}
The global proof dependency DAG factors through a narrow spine of high-leverage bottleneck theorems separating large upstream developments from downstream applications. Table~\ref{tab:bottlenecks} summarizes the highest-scoring formal bottlenecks.
\begin{table}[ht]
\centering
\caption{Top Mathematical Bottleneck Pillars}\label{tab:bottlenecks}
\small
\begin{tabular}{rp{0.56\textwidth}rr}
\toprule
\# & Pillar Theorem & Body LOC & Bottleneck Score \\
\midrule
1 & \raggedright gluing mass holds \leanlink{Formalism.L2.SG.Pf.Cnt.gluing_mass_holds} & 2 & 15.44 \\
2 & \raggedright alpha discriminant holds \leanlink{Formalism.L2.SG.Pf.Cnt.alpha_discriminant_holds} & 3 & 14.86 \\
3 & \raggedright se Good \leanlink{GlobalTheory.HeightCounting.SerCopy.seGood} & 17 & 14.66 \\
4 & \raggedright li Stat \leanlink{Formalism.HermitianLattices.Ext.L3.liStat} & 1 & 14.27 \\
5 & \raggedright family Unif of laws \leanlink{GlobalTheory.familyUnif_of_laws} & 6 & 14.05 \\
6 & \raggedright holds tr \leanlink{GlobalTheory.Mem.R3.Tr.holds_tr} & 7 & 13.93 \\
7 & \raggedright detects \leanlink{GlobalTheory.GeoTie.SideTie.detects} & 6 & 13.73 \\
8 & \raggedright hecke Std Plain \leanlink{GlobalTheory.HeckeStdFacts.heckeStdPlain} & 17 & 13.64 \\
9 & \raggedright class Reps Definite \leanlink{Formalism.ThetaLiftIndependence.classRepsDefinite} & 22 & 13.60 \\
10 & \raggedright row lattice mass \leanlink{PI.UnitaryGroupMassFormula.CntFormula.RowsMassW.row_lattice_mass} & 41 & 13.43 \\
11 & \raggedright mass Basic \leanlink{Formalism.HermitianLatticesTheory.massBasic} & 21 & 13.25 \\
12 & \raggedright normalizing arithmetic holds \leanlink{Formalism.L2.SG.Pf.NcGlob.normalizing_arithmetic_holds} & 197 & 12.85 \\
13 & \raggedright vec Lin bijective \leanlink{GlobalTheory.ThFr.Gauss.vecLin_bijective} & 21 & 12.85 \\
14 & \raggedright phi counting holds \leanlink{Formalism.L2.SG.Pf.LocalizedFibreSets.phi_counting_holds} & 1 & 12.46 \\
15 & \raggedright real Adelic Inert Transport further Rows \leanlink{GlobalTheory.Mem.R3.Tr.realAdelicInertTransport_furtherRows} & 7 & 12.41 \\
16 & \raggedright phi mod p holds \leanlink{Formalism.L2.SG.Pf.LocalizedFibreSets.phi_mod_p_holds} & 114 & 12.36 \\
17 & \raggedright segment Representation \leanlink{PositivityTheory.RgGu.segmentRepresentation} & 4 & 12.21 \\
18 & \raggedright mult eq \leanlink{GlobalTheory.Mem.R3.Tr.mult_eq} & 5 & 12.18 \\
19 & \raggedright same Std \leanlink{GlobalTheory.Mem.SRe.CarModel.sameStd} & 6 & 12.15 \\
20 & \raggedright li Ratio \leanlink{Formalism.HermitianLattices.Ext.L3.liRatio} & 18 & 12.11 \\
\bottomrule
\end{tabular}
\end{table}

\subsection{Chapter Dependency Architecture}
The reachable proof graph decomposes into 12 major chapters ordered topologically from foundational structures to the main conclusions:
\begin{enumerate}
\item \textbf{Reciprocity Theory / Commutative Algebra} (126 declarations, 2,113~LOC).
\item \textbf{Formalism / Log Sum Exp Bounds: decay} (65 declarations, 893~LOC).
\item \textbf{Formalism / Lattice Theta: lattice fourier series} (66 declarations, 655~LOC).
\item \textbf{Development / Abelian Variety Theory: se Good} (337,931 declarations, 4,621,552~LOC) (builds on Reciprocity Theory: segment Representation, Ground Theory: zel, Formalism: gluing mass holds).
\item \textbf{Reciprocity Theory: segment Representation} (163,457 declarations, 2,120,740~LOC) (builds on Development / Abelian Variety Theory: se Good, Ground Theory: zel, Formalism: gluing mass holds).
\item \textbf{Formalism: gluing mass holds} (173,144 declarations, 1,931,701~LOC) (builds on Formalism / Lattice Theta: lattice fourier series, Formalism: flag Deg Formula statement, Reciprocity Theory: segment Representation, Development / Abelian Variety Theory: se Good).
\item \textbf{Ground Theory: zel} (116,103 declarations, 1,632,709~LOC) (builds on Reciprocity Theory: segment Representation, Development / Abelian Variety Theory: se Good, Formalism: gluing mass holds).
\item \textbf{Development / Resolution Of Singularities: functorial Resolution Char Zero} (7,590 declarations, 116,064~LOC) (builds on Ground Theory: zel).
\item \textbf{Reciprocity Theory / Geometry Of Numbers: exists chi P succ sub le} (227 declarations, 2,762~LOC) (builds on Ground Theory: zel).
\item \textbf{Companion Theory / Thickening Estimates: nrm Up} (216 declarations, 1,829~LOC) (builds on Development / Abelian Variety Theory: se Good, Ground Theory: zel, Formalism: gluing mass holds).
\item \textbf{Formalism: host Conj of} (72 declarations, 1,017~LOC) (builds on Formalism: gluing mass holds).
\item \textbf{Formalism: flag Deg Formula statement} (67 declarations, 462~LOC) (builds on Formalism: gluing mass holds).
\end{enumerate}

\section{Reciprocity Theory / Commutative Algebra}

This chapter encompasses 126 formal declarations (2,113~LOC).
\begin{quote}\small
\textbf{Formal Interface ($\mathrm{Inputs} \vdash \mathrm{Outputs}$):}
\par\noindent\emph{Outputs:} is Domain~\leanlink{ReciprocityTheory.CA.AB.RegA.IsRegularLocalRing.isDomain}, Is CM~\leanlink{ReciprocityTheory.CA.IsCM}, Mod Depth GE~\leanlink{ReciprocityTheory.CA.ModDepthGE}, exists dim trdeg field~\leanlink{ReciprocityTheory.CA.exists_dim_trdeg_field}, Linked~\leanlink{ReciprocityTheory.CA.Linked}, ring Krull Dim eq of is Integral of injective~\leanlink{ReciprocityTheory.CA.ringKrullDim_eq_of_isIntegral_of_injective}.
\end{quote}

\subsection{Core Definitions and Structures}
\begin{definition}[Mod Depth GE]\leanlink{ReciprocityTheory.CA.ModDepthGE}
\[
(n : \mathbb{N} ) : \mathrm{Prop}
\]
\end{definition}

\begin{definition}[Is CM]\leanlink{ReciprocityTheory.CA.IsCM}
\[
: \mathrm{Prop}
\]
\end{definition}

\begin{definition}[Depth GE]\leanlink{ReciprocityTheory.CA.DepthGE}
\[
(n : \mathbb{N} ) : \mathrm{Prop}
\]
\end{definition}

\begin{definition}[quot SMul Top Equiv Quot']\leanlink{ReciprocityTheory.CA.quotSMulTopEquivQuot'}
\[
(x : R) : \mathrm{QuotSMulTop} x R \simeq _{l} R / \mathrm{span}
\]
\end{definition}

\begin{definition}[Is CMQuot]\leanlink{ReciprocityTheory.CA.IsCMQuot}
\[
(R : \mathrm{Type}u) : \mathrm{Prop}
\]
\end{definition}

\begin{definition}[Conn In Dim]\leanlink{ReciprocityTheory.CA.ConnInDim}
\[
(k : \mathrm{WithBot} \mathbb{N} \infty ) : \mathrm{Prop}
\]
\end{definition}

\subsection{Main Results}
\begin{theorem}[ring Krull Dim localization quotient eq height]\leanlink{ReciprocityTheory.CA.ringKrullDim_localization_quotient_eq_height}
\begin{quote}\small\raggedright$\displaystyle (p q : \mathrm{Ideal} P) \allowbreak (\mathrm{hpq} : p \le q) : \mathrm{ringKrullDim} (\mathrm{AtPrime} q / p.\mathrm{map} (\mathrm{algebraMap} P (\mathrm{AtPrime} q))) = ((q.\mathrm{map} (\mathrm{mk} p)).\mathrm{height} : \mathrm{WithBot} \mathbb{N} \infty )$\end{quote}
\end{theorem}
\begin{proof}
Combining \emph{height le height Iic lemma 2} \leanlink{ReciprocityTheory.CA.height_le_height_Iic_lemma_2}, \emph{height le krull Dim lemma 6} \leanlink{ReciprocityTheory.CA.height_le_krullDim_lemma_6}, \emph{height eq order height lemma} \leanlink{ReciprocityTheory.CommutativeAlgebra.SaturatedChains.RingKrullDimLocalizationQuotientEqHeight::ReciprocityTheory.CA.height_eq_order_height_lemma}, \emph{comap map mk lemma 4} \leanlink{ReciprocityTheory.CA.comap_map_mk_lemma_4} with \emph{krull Dim le height lemma 3} \leanlink{ReciprocityTheory.CA.krullDim_le_height_lemma_3} yields the stated result.
\end{proof}

\begin{theorem}[associated Primes eq minimal Primes]\leanlink{ReciprocityTheory.CA.IsCMMod.associatedPrimes_eq_minimalPrimes}
\begin{quote}\small\raggedright$\displaystyle (h : \mathrm{IsCMMod} R M) : \mathrm{associatedPrimes} R M = (\mathrm{annihilator} R M).\mathrm{minimalPrimes}$\end{quote}
\end{theorem}
\begin{proof}
The argument operates on the structural objects Is CMMod \leanlink{ReciprocityTheory.CA.IsCMMod}. Combining \emph{ring Krull Dim quotient eq support Dim} \leanlink{ReciprocityTheory.CA.IsCMMod.ringKrullDim_quotient_eq_supportDim} with \emph{unm dim le of le} \leanlink{ReciprocityTheory.CA.unm_dim_le_of_le} yields the stated result.
\end{proof}

\begin{proposition}[span Finrank maximal Ideal quotient succ le]\leanlink{ReciprocityTheory.CA.AB.RegA.IsRegularLocalRing.spanFinrank_maximalIdeal_quotient_succ_le}
\begin{quote}\small\raggedright$\displaystyle (R : \mathrm{Type}) \allowbreak (hx : x \in \mathrm{maximalIdeal} R) \allowbreak (\mathrm{hx2} : x \notin \mathrm{maximalIdeal} R ^ 2) : \mathrm{spanFinrank} (\mathrm{maximalIdeal} (R / \mathrm{span} ( : \mathcal{P} R))) + 1 \le \mathrm{finrank} (\mathrm{ResidueField} R) \allowbreak (\mathrm{CotangentSpace} R)$\end{quote}
\end{proposition}
\begin{proof}
Combining \emph{span insert image erase eq top} \leanlink{ReciprocityTheory.CA.AB.RegA.IsRegularLocalRing.span_insert_image_erase_eq_top}, \emph{span Finrank span finite sum le} \leanlink{ReciprocityTheory.CA.AB.RegA.IsRegularLocalRing.spanFinrank_span_finset_le} with \emph{exists repr ne zero} \leanlink{ReciprocityTheory.CA.AB.RegA.IsRegularLocalRing.exists_repr_ne_zero} yields the stated result.
\end{proof}

\begin{proposition}[is Domain of forall span Finrank lt]\leanlink{ReciprocityTheory.CA.AB.RegA.IsRegularLocalRing.isDomain_of_forall_spanFinrank_lt}
\begin{quote}\small\raggedright$\displaystyle (R : \mathrm{Type}u) \allowbreak (IH : \forall (S : \mathrm{Type}u) , \allowbreak \mathrm{spanFinrank} (\mathrm{maximalIdeal} S) < \mathrm{spanFinrank} (\mathrm{maximalIdeal} R) \to \allowbreak \mathrm{IsDomain} S) : \mathrm{IsDomain} R$\end{quote}
\end{proposition}
\begin{proof}
Combining \emph{span Finrank quotient eq and ring Krull Dim add one} \leanlink{ReciprocityTheory.CA.AB.RegA.IsRegularLocalRing.spanFinrank_quotient_eq_and_ringKrullDim_add_one}, \emph{exists mem maximal Ideal not Mem sq not Mem minimal Primes} \leanlink{ReciprocityTheory.CA.AB.RegA.IsRegularLocalRing.exists_mem_maximalIdeal_notMem_sq_notMem_minimalPrimes}, \emph{quotient span singleton} \leanlink{ReciprocityTheory.CA.AB.RegA.IsRegularLocalRing.quotient_span_singleton}, \emph{maximal Ideal eq bot of le sq} \leanlink{ReciprocityTheory.CA.AB.RegA.IsRegularLocalRing.maximalIdeal_eq_bot_of_le_sq}, \emph{is Domain of bot is Prime} \leanlink{ReciprocityTheory.CA.AB.RegA.IsRegularLocalRing.isDomain_of_bot_isPrime}, \emph{maximal Ideal le sq of finrank cotangent zero} \leanlink{ReciprocityTheory.CA.AB.RegA.IsRegularLocalRing.maximalIdeal_le_sq_of_finrank_cotangent_zero}, \emph{is Domain of maximal Ideal eq bot} \leanlink{ReciprocityTheory.CA.AB.RegA.IsRegularLocalRing.isDomain_of_maximalIdeal_eq_bot} with \emph{is Prime of mem minimal Primes} \leanlink{ReciprocityTheory.CA.AB.RegA.IsRegularLocalRing.isPrime_of_mem_minimalPrimes} yields the stated result.
\end{proof}

\begin{proposition}[ring Krull Dim quotient eq support Dim]\leanlink{ReciprocityTheory.CA.IsCMMod.ringKrullDim_quotient_eq_supportDim}
\begin{quote}\small\raggedright$\displaystyle (h : \mathrm{IsCMMod} R M) \allowbreak (hp : p \in \mathrm{associatedPrimes} R M) : \mathrm{ringKrullDim} (R / p) = \mathrm{supportDim} R M$\end{quote}
\end{proposition}
\begin{proof}
The argument operates on the structural objects Is CMMod \leanlink{ReciprocityTheory.CA.IsCMMod}. Combining \emph{le ring Krull Dim quotient of mem associated Primes} \leanlink{ReciprocityTheory.CA.ModDepthGE.le_ringKrullDim_quotient_of_mem_associatedPrimes} with \emph{unm dim le of le} \leanlink{ReciprocityTheory.CA.unm_dim_le_of_le} yields the stated result.
\end{proof}

\begin{proposition}[is CM localization mv Polynomial]\leanlink{ReciprocityTheory.CA.isCM_localization_mvPolynomial}
\begin{quote}\small\raggedright$\displaystyle (hP : \mathrm{IsCM} P) \allowbreak (N : \mathbb{N} ) \allowbreak (N : \mathrm{Ideal} (\mathrm{MvPolynomial} (\mathrm{Fin} N) P)) \allowbreak (h : N.\mathrm{comap} \mathrm{C} = \mathrm{maximalIdeal} P) : \mathrm{IsCM} (\mathrm{AtPrime} N) \land \allowbreak (N.\mathrm{height} : \mathrm{WithBot} \mathbb{N} \infty ) = \mathrm{ringKrullDim} P + N$\end{quote}
\end{proposition}
\begin{proof}
The argument operates on the structural objects Is CM \leanlink{ReciprocityTheory.CA.IsCM}. Combining \emph{exists is Weakly Regular mv Polynomial} \leanlink{ReciprocityTheory.CA.exists_isWeaklyRegular_mvPolynomial}, \emph{length le height of is Weakly Regular} \leanlink{ReciprocityTheory.CA.length_le_height_of_isWeaklyRegular} with \emph{depth GE localization of is Weakly Regular} \leanlink{ReciprocityTheory.CA.depthGE_localization_of_isWeaklyRegular} yields the stated result.
\end{proof}

\begin{lemma}[span Finrank quotient eq and ring Krull Dim add one]\leanlink{ReciprocityTheory.CA.AB.RegA.IsRegularLocalRing.spanFinrank_quotient_eq_and_ringKrullDim_add_one}
\begin{quote}\small\raggedright$\displaystyle (R : \mathrm{Type}) \allowbreak (hx : x \in \mathrm{maximalIdeal} R) \allowbreak (\mathrm{hx2} : x \notin \mathrm{maximalIdeal} R ^ 2) : ((\mathrm{spanFinrank} (\mathrm{maximalIdeal} (R / \mathrm{span} ( : \mathcal{P} R))) : \mathbb{N} ) : \mathrm{WithBot} \mathbb{N} \infty ) = \mathrm{ringKrullDim} (R / \mathrm{span} ( : \mathcal{P} R)) \land \allowbreak \mathrm{ringKrullDim} (R / \mathrm{span} ( : \mathcal{P} R)) + 1 = \mathrm{ringKrullDim} R$\end{quote}
\end{lemma}
\begin{proof}
The claim follows directly from \emph{span Finrank maximal Ideal quotient succ le} \leanlink{ReciprocityTheory.CA.AB.RegA.IsRegularLocalRing.spanFinrank_maximalIdeal_quotient_succ_le}.
\end{proof}

\begin{lemma}[exists regular sequence cons]\leanlink{ReciprocityTheory.CA.exists_regular_sequence_cons}
\begin{quote}\small\raggedright$\displaystyle (n : \mathbb{N} ) : \forall (N : \mathrm{Type}u) , \allowbreak x \in \mathrm{maximalIdeal} R \to \allowbreak \mathrm{IsSMulRegular} N x \to \allowbreak \mathrm{ModDepthGE} R N (n + 1) \to \allowbreak \exists rs : \mathrm{List} R, \allowbreak rs.\mathrm{length} = n \land \allowbreak (\forall r \in rs, \allowbreak r \in \mathrm{maximalIdeal} R) \land \allowbreak \mathrm{IsRegular} N (x :: rs)$\end{quote}
\end{lemma}
\begin{proof}
The argument operates on the structural objects Mod Depth GE \leanlink{ReciprocityTheory.CA.ModDepthGE}. Combining \emph{max not mem ass swap} \leanlink{ReciprocityTheory.CA.max_not_mem_ass_swap}, \emph{exists reg both} \leanlink{ReciprocityTheory.CA.exists_reg_both} with \emph{max not mem ass of reg} \leanlink{ReciprocityTheory.CA.max_not_mem_ass_of_reg} yields the stated result.
\end{proof}

\begin{lemma}[exists is Weakly Regular mv Polynomial]\leanlink{ReciprocityTheory.CA.exists_isWeaklyRegular_mvPolynomial}
\begin{quote}\small\raggedright$\displaystyle (\mathrm{hmem} : \forall r \in rs, \allowbreak r \in \mathrm{maximalIdeal} P) \allowbreak (\mathrm{hrs} : \mathrm{IsWeaklyRegular} P rs) \allowbreak (N : \mathbb{N} ) \allowbreak (N : \mathrm{Ideal} (\mathrm{MvPolynomial} (\mathrm{Fin} N) P)) \allowbreak (h : N.\mathrm{comap} \mathrm{C} = \mathrm{maximalIdeal} P) : \exists ss : \mathrm{List} (\mathrm{MvPolynomial} (\mathrm{Fin} N) P), \allowbreak ss.\mathrm{length} = rs.\mathrm{length} + N \land \allowbreak (\forall s \in ss, \allowbreak s \in N) \land \allowbreak \mathrm{IsWeaklyRegular} (\mathrm{MvPolynomial} (\mathrm{Fin} N) P) ss$\end{quote}
\end{lemma}
\begin{proof}
Combining \emph{is Weakly Regular map C append monic} \leanlink{ReciprocityTheory.CA.isWeaklyRegular_map_C_append_monic}, \emph{is Maximal comap of is Maximal of comap eq} \leanlink{ReciprocityTheory.CA.isMaximal_comap_of_isMaximal_of_comap_eq}, \emph{exists monic aeval mem of is Maximal} \leanlink{ReciprocityTheory.CA.exists_monic_aeval_mem_of_isMaximal} with \emph{is Weakly Regular map ring Equiv} \leanlink{ReciprocityTheory.CA.isWeaklyRegular_map_ringEquiv} yields the stated result.
\end{proof}

\begin{lemma}[exists is Regular of le height]\leanlink{ReciprocityTheory.CA.IsCM.exists_isRegular_of_le_height}
\begin{quote}\small\raggedright$\displaystyle (h : \mathrm{IsCM} R) \allowbreak (hI : I \neq \top ) \allowbreak (hr : (r : \mathbb{N} \infty ) \le I.\mathrm{height}) : \exists rs : \mathrm{List} R, \allowbreak rs.\mathrm{length} = r \land \allowbreak (\forall x \in rs, \allowbreak x \in I) \land \allowbreak \mathrm{IsRegular} R rs$\end{quote}
\end{lemma}
\begin{proof}
The argument operates on the structural objects Is CM \leanlink{ReciprocityTheory.CA.IsCM}. Combining \emph{grade step exists} \leanlink{ReciprocityTheory.CA.grade_step_exists} with \emph{is Regular of is Weakly Regular of mem ne top} \leanlink{ReciprocityTheory.CA.isRegular_of_isWeaklyRegular_of_mem_ne_top} yields the stated result.
\end{proof}

\begin{lemma}[quotient of List]\leanlink{ReciprocityTheory.CA.IsCMMod.quotient_ofList}
\begin{quote}\small\raggedright$\displaystyle (h : \mathrm{IsCMMod} R M) \allowbreak (\mathrm{hrs} : \forall r \in rs, \allowbreak r \in \mathrm{maximalIdeal} R) \allowbreak (\mathrm{reg} : \mathrm{IsRegular} M rs) : \mathrm{IsCMMod} R (M / \mathrm{ofList} rs \cdot (\top : \mathrm{Submodule} R M))$\end{quote}
\end{lemma}
\begin{proof}
The argument operates on the structural objects Is CMMod \leanlink{ReciprocityTheory.CA.IsCMMod}. Combining \emph{quot SMul Top} \leanlink{ReciprocityTheory.CA.IsCMMod.quotSMulTop} with \emph{unm is CMMod congr} \leanlink{ReciprocityTheory.CA.unm_isCMMod_congr} yields the stated result.
\end{proof}

\begin{lemma}[is CM quotient iff]\leanlink{ReciprocityTheory.CA.isCM_quotient_iff}
\[
(I : \mathrm{Ideal} R) : \mathrm{IsCM} (R / I) \iff \mathrm{IsCMMod} R (R / I)
\]
\end{lemma}
\begin{proof}
The argument operates on the structural objects Is CMMod \leanlink{ReciprocityTheory.CA.IsCMMod}, Is CM \leanlink{ReciprocityTheory.CA.IsCM}. The claim follows directly from \emph{depth GE quotient iff} \leanlink{ReciprocityTheory.CA.depthGE_quotient_iff}.
\end{proof}

\section{Formalism / Log Sum Exp Bounds: decay}

This chapter encompasses 65 formal declarations (893~LOC).
\begin{quote}\small
\textbf{Formal Interface ($\mathrm{Inputs} \vdash \mathrm{Outputs}$):}
\par\noindent\emph{Outputs:} decay~\leanlink{Formalism.NuclearRankWeightBounds.Top.decay}, descend mk~\leanlink{Formalism.NuclearRankWeightBounds.P.descend_mk}, PB~\leanlink{Formalism.NuclearRankWeightBounds.P.PB}, lam~\leanlink{Formalism.NuclearRankWeightBounds.C.lam}, a K~\leanlink{Formalism.NuclearRankWeightBounds.P.aK}, uniq~\leanlink{Formalism.NuclearRankWeightBounds.P.Pres.uniq}.
\end{quote}

\subsection{Core Definitions and Structures}
\begin{definition}[row Inv Mod]\leanlink{Formalism.NuclearRankWeightBounds.C.rowInvMod}
\begin{quote}\small\raggedright$\displaystyle (Z : \mathrm{Matrix} (\mathrm{Fin} k) \allowbreak (\mathrm{Fin} k) RR) \allowbreak (r : \mathrm{Fin} k \to \allowbreak K) : (\mathrm{Fin} k \to \allowbreak RR) \to \allowbreak _{l} \mathrm{KmodR} RR K$\end{quote}
\end{definition}

\begin{definition}[descend]\leanlink{Formalism.NuclearRankWeightBounds.C.descend}
\begin{quote}\small\raggedright$\displaystyle (Z : \mathrm{Matrix} (\mathrm{Fin} k) \allowbreak (\mathrm{Fin} k) RR) \allowbreak (r : \mathrm{Fin} k \to \allowbreak K) : \mathrm{coker} Z \to \allowbreak _{l} \mathrm{KmodR} RR K$\end{quote}
\end{definition}

\begin{definition}[PB]\leanlink{Formalism.NuclearRankWeightBounds.P.PB}
\begin{quote}\small\raggedright$\displaystyle (Z : \mathrm{Matrix} (\mathrm{Fin} (m+1)) \allowbreak (\mathrm{Fin} (m+1)) RR) \allowbreak (hZ : \mathrm{IsUnit} (\mathrm{toK} K Z).\mathrm{det}) \allowbreak (hZ' : \mathrm{IsUnit} (\mathrm{toK} K (\mathrm{blk} Z)).\mathrm{det}) : \mathrm{Pres} RR K (MM Z) \allowbreak (\mathrm{coker} (\mathrm{blk} Z))$\end{quote}
\end{definition}

\begin{definition}[lam]\leanlink{Formalism.NuclearRankWeightBounds.C.lam}
\begin{quote}\small\raggedright$\displaystyle (Z : \mathrm{Matrix} (\mathrm{Fin} (m+1)) \allowbreak (\mathrm{Fin} (m+1)) RR) : \mathrm{coker} Z \to \allowbreak _{l} \mathrm{KmodR} RR K$\end{quote}
\end{definition}

\begin{definition}[to K]\leanlink{Formalism.NuclearRankWeightBounds.C.toK}
\begin{quote}\small\raggedright$\displaystyle (Z : \mathrm{Matrix} (\mathrm{Fin} k) \allowbreak (\mathrm{Fin} k) RR) : \mathrm{Matrix} (\mathrm{Fin} k) \allowbreak (\mathrm{Fin} k) K$\end{quote}
\end{definition}

\begin{definition}[coker]\leanlink{Formalism.NuclearRankWeightBounds.C.coker}
\[
(Z : \mathrm{Matrix} (\mathrm{Fin} k) (\mathrm{Fin} k) RR) : \mathrm{Type}\_
\]
\end{definition}

\subsection{Main Results}
\begin{theorem}[decay]\leanlink{Formalism.NuclearRankWeightBounds.Top.decay}
\[
\mathrm{Decay}
\]
\end{theorem}
\begin{proof}
The argument operates on the structural objects Decay \leanlink{Formalism.NuclearRankWeightBounds.D.Decay}. Combining \emph{cb weights} \leanlink{Formalism.NuclearRankWeightBounds.P.cb_weights}, \emph{log sum exp bound holds} \leanlink{Formalism.NuclearRankWeightBounds.P.log_sum_exp_bound_holds}, \emph{rank block holds} \leanlink{Formalism.NuclearRankWeightBounds.P.rank_block_holds}, \emph{nuclear norm bound holds} \leanlink{Formalism.NuclearRankWeightBounds.P.nuclear_norm_bound_holds}, \emph{weighted off diagonal holds} \leanlink{Formalism.NuclearRankWeightBounds.P.weighted_off_diagonal_holds} with \emph{decay of facts} \leanlink{Formalism.NuclearRankWeightBounds.P.decay_of_facts} yields the stated result.
\end{proof}

\begin{theorem}[cb weights]\leanlink{Formalism.NuclearRankWeightBounds.P.cb_weights}
\[
\mathrm{CBWeights}
\]
\end{theorem}
\begin{proof}
The argument operates on the structural objects CBWeights \leanlink{Formalism.NuclearRankWeightBounds.D.CBWeights}, lead \leanlink{Formalism.NuclearRankWeightBounds.D.lead}. Combining \emph{two coeff two det} \leanlink{Formalism.NuclearRankWeightBounds.P.two_coeff_two_det}, \emph{det mul diagonal mul eq sum} \leanlink{Formalism.NuclearRankWeightBounds.P.det_mul_diagonal_mul_eq_sum}, \emph{moment} \leanlink{Formalism.NuclearRankWeightBounds.P.moment} with \emph{coeff prod one add} \leanlink{Formalism.NuclearRankWeightBounds.P.coeff_prod_one_add} yields the stated result.
\end{proof}

\begin{proposition}[log sum exp bound holds]\leanlink{Formalism.NuclearRankWeightBounds.P.log_sum_exp_bound_holds}
\[
\mathrm{LogSumExp}
\]
\end{proposition}
\begin{proof}
The argument operates on the structural objects Log Sum Exp \leanlink{Formalism.NuclearRankWeightBounds.D.LogSumExp}. The claim follows directly from \emph{exp ge quad} \leanlink{Formalism.NuclearRankWeightBounds.P.exp_ge_quad}.
\end{proof}

\begin{proposition}[rank block holds]\leanlink{Formalism.NuclearRankWeightBounds.P.rank_block_holds}
\[
\mathrm{RankBlock}
\]
\end{proposition}
\begin{proof}
The argument operates on the structural objects Rank Block \leanlink{Formalism.NuclearRankWeightBounds.D.RankBlock}, q0 \leanlink{Formalism.NuclearRankWeightBounds.D.q0}. Combining \emph{rank le of support} \leanlink{Formalism.NuclearRankWeightBounds.P.rank_le_of_support} with \emph{rank add le'} \leanlink{Formalism.NuclearRankWeightBounds.P.rank_add_le'} yields the stated result.
\end{proof}

\begin{proposition}[nuclear norm bound holds]\leanlink{Formalism.NuclearRankWeightBounds.P.nuclear_norm_bound_holds}
\[
\mathrm{NuclearBound}
\]
\end{proposition}
\begin{proof}
The argument operates on the structural objects Nuclear Bound \leanlink{Formalism.NuclearRankWeightBounds.D.NuclearBound}, hs \leanlink{Formalism.NuclearRankWeightBounds.D.hs}. The claim follows directly from \emph{sum sq eigenvalues} \leanlink{Formalism.NuclearRankWeightBounds.P.sum_sq_eigenvalues}.
\end{proof}

\begin{proposition}[weighted off diagonal holds]\leanlink{Formalism.NuclearRankWeightBounds.P.weighted_off_diagonal_holds}
\[
\mathrm{WeightedOff}
\]
\end{proposition}
\begin{proof}
The argument operates on the structural objects Weighted Off \leanlink{Formalism.NuclearRankWeightBounds.D.WeightedOff}. This result is established directly from standard library foundations within its local module (69~lines of formal proof).
\end{proof}

\begin{lemma}[key cval]\leanlink{Formalism.NuclearRankWeightBounds.P.key_cval}
\begin{quote}\small\raggedright$\displaystyle (Z : \mathrm{Matrix} (\mathrm{Fin} (m+1)) \allowbreak (\mathrm{Fin} (m+1)) RR) \allowbreak (hZ : \mathrm{IsUnit} (\mathrm{toK} K Z).\mathrm{det}) \allowbreak (hZ' : \mathrm{IsUnit} (\mathrm{toK} K (\mathrm{blk} Z)).\mathrm{det}) \allowbreak (v : \mathrm{Fin} (m+1) \to \allowbreak RR) : \mathrm{algebraMap} RR K (v (\mathrm{Fin}.\mathrm{last} m)) - aK Z \cdot _{v} ((\mathrm{toK} K (\mathrm{blk} Z))^{-1} *_{v} vk (\lambda i \mapsto v (\mathrm{Fin}.\mathrm{castSucc} i))) = \mathrm{cval} K Z * fV Z v$\end{quote}
\end{lemma}
\begin{proof}
The argument operates on the structural objects to K \leanlink{Formalism.NuclearRankWeightBounds.C.toK}, f V \leanlink{Formalism.NuclearRankWeightBounds.P.fV}, cval \leanlink{Formalism.NuclearRankWeightBounds.C.cval}, a K \leanlink{Formalism.NuclearRankWeightBounds.P.aK}. The claim follows directly from \emph{key} \leanlink{Formalism.NuclearRankWeightBounds.P.key}.
\end{proof}

\begin{lemma}[decay of facts]\leanlink{Formalism.NuclearRankWeightBounds.P.decay_of_facts}
\begin{quote}\small\raggedright$\displaystyle \mathrm{CBWeights} \to \allowbreak \mathrm{NuclearBound} \to \allowbreak \mathrm{RankBlock} \to \allowbreak \mathrm{LogSumExp} \to \allowbreak \mathrm{WeightedOff} \to \allowbreak \mathrm{Decay}$\end{quote}
\end{lemma}
\begin{proof}
The argument operates on the structural objects CBWeights \leanlink{Formalism.NuclearRankWeightBounds.D.CBWeights}, Decay \leanlink{Formalism.NuclearRankWeightBounds.D.Decay}, Weighted Off \leanlink{Formalism.NuclearRankWeightBounds.D.WeightedOff}, Log Sum Exp \leanlink{Formalism.NuclearRankWeightBounds.D.LogSumExp}. This result is established directly from standard library foundations within its local module (92~lines of formal proof).
\end{proof}

\begin{lemma}[alpha mk]\leanlink{Formalism.NuclearRankWeightBounds.P.alpha_mk}
\begin{quote}\small\raggedright$\displaystyle (Z : \mathrm{Matrix} (\mathrm{Fin} (m+1)) \allowbreak (\mathrm{Fin} (m+1)) RR) \allowbreak (hZ : \mathrm{IsUnit} (\mathrm{toK} K (\mathrm{blk} Z)).\mathrm{det}) \allowbreak (v : \mathrm{Fin} m \to \allowbreak RR) : \mathrm{alpha} K Z (\mathrm{mk} v) = \mathrm{mk} (aK Z \cdot _{v} ((\mathrm{toK} K (\mathrm{blk} Z))^{-1} *_{v} vk v))$\end{quote}
\end{lemma}
\begin{proof}
The argument operates on the structural objects row Inv Mod \leanlink{Formalism.NuclearRankWeightBounds.C.rowInvMod}, descend \leanlink{Formalism.NuclearRankWeightBounds.C.descend}, to K \leanlink{Formalism.NuclearRankWeightBounds.C.toK}, alpha \leanlink{Formalism.NuclearRankWeightBounds.C.alpha}. Combining \emph{vk mul Vec} \leanlink{Formalism.NuclearRankWeightBounds.P.vk_mulVec} with \emph{row Inv Mod apply} \leanlink{Formalism.NuclearRankWeightBounds.P.rowInvMod_apply} yields the stated result.
\end{proof}

\begin{lemma}[key]\leanlink{Formalism.NuclearRankWeightBounds.P.key}
\begin{quote}\small\raggedright$\displaystyle (W : \mathrm{Matrix} (\mathrm{Fin} (m+1)) \allowbreak (\mathrm{Fin} (m+1)) \alpha ) \allowbreak (hW : \mathrm{IsUnit} (W.\mathrm{submatrix} \mathrm{Fin}.\mathrm{castSucc} \mathrm{Fin}.\mathrm{castSucc}).\mathrm{det}) \allowbreak (y : \mathrm{Fin} (m+1) \to \allowbreak \alpha ) : (W *_{v} y) \allowbreak (\mathrm{Fin}.\mathrm{last} m) - (\lambda j \mapsto W (\mathrm{Fin}.\mathrm{last} m) \allowbreak (\mathrm{Fin}.\mathrm{castSucc} j)) \cdot _{v} ((W.\mathrm{submatrix} \mathrm{Fin}.\mathrm{castSucc} \mathrm{Fin}.\mathrm{castSucc})^{-1} *_{v} (\lambda i \mapsto (W *_{v} y) \allowbreak (\mathrm{Fin}.\mathrm{castSucc} i))) = (W (\mathrm{Fin}.\mathrm{last} m) \allowbreak (\mathrm{Fin}.\mathrm{last} m) - \sum i, \allowbreak \sum j, \allowbreak W (\mathrm{Fin}.\mathrm{last} m) \allowbreak (\mathrm{Fin}.\mathrm{castSucc} i) * (W.\mathrm{submatrix} \mathrm{Fin}.\mathrm{castSucc} \mathrm{Fin}.\mathrm{castSucc})^{-1} i j * W (\mathrm{Fin}.\mathrm{castSucc} j) \allowbreak (\mathrm{Fin}.\mathrm{last} m)) * y (\mathrm{Fin}.\mathrm{last} m)$\end{quote}
\end{lemma}
\begin{proof}
The claim follows directly from \emph{mul Vec last} \leanlink{Formalism.NuclearRankWeightBounds.P.mulVec_last}.
\end{proof}

\begin{lemma}[two coeff two det]\leanlink{Formalism.NuclearRankWeightBounds.P.two_coeff_two_det}
\begin{quote}\small\raggedright$\displaystyle (M : \mathrm{Matrix} (\mathrm{Fin} k) \allowbreak (\mathrm{Fin} k) \mathbb{C} ) \allowbreak (hM : M.\mathrm{IsHermitian}) : 2 * (\mathrm{det} (1 + (X : \mathbb{C} ) \cdot M.\mathrm{map} C)).\mathrm{coeff} 2 = (\mathrm{trace} M) ^ 2 - \mathrm{trace} (M * M)$\end{quote}
\end{lemma}
\begin{proof}
Combining \emph{moment} \leanlink{Formalism.NuclearRankWeightBounds.P.moment} with \emph{coeff prod one add} \leanlink{Formalism.NuclearRankWeightBounds.P.coeff_prod_one_add} yields the stated result.
\end{proof}

\begin{lemma}[f V single]\leanlink{Formalism.NuclearRankWeightBounds.P.fV_single}
\begin{quote}\small\raggedright$\displaystyle (Z : \mathrm{Matrix} (\mathrm{Fin} (m+1)) \allowbreak (\mathrm{Fin} (m+1)) RR) : fV (K := K) Z (Pi.\mathrm{single} (\mathrm{Fin}.\mathrm{last} m) 1) = \mathrm{kappa} K Z$\end{quote}
\end{lemma}
\begin{proof}
The argument operates on the structural objects f V \leanlink{Formalism.NuclearRankWeightBounds.P.fV}, kappa \leanlink{Formalism.NuclearRankWeightBounds.C.kappa}. The claim follows directly from \emph{vk single} \leanlink{Formalism.NuclearRankWeightBounds.P.vk_single}.
\end{proof}

\section{Formalism / Lattice Theta: lattice fourier series}

This chapter encompasses 66 formal declarations (655~LOC). Primary proof idioms detected: \emph{Cauchy-Schwarz / Parseval}.
\begin{quote}\small
\textbf{Formal Interface ($\mathrm{Inputs} \vdash \mathrm{Outputs}$):}
\par\noindent\emph{Outputs:} dual~\leanlink{Formalism.H.dual}, lattice fourier series~\leanlink{Formalism.H.lattice_fourier_series}, dual Is Lattice~\leanlink{Formalism.H.dualIsLattice}, theta~\leanlink{Formalism.H.theta}, poisson periodization~\leanlink{Formalism.H.poisson_periodization}, unfolding fourier coefficient~\leanlink{Formalism.H.unfolding_fourier_coefficient}.
\end{quote}

\subsection{Core Definitions and Structures}
\begin{definition}[dual]\leanlink{Formalism.H.dual}
\[
(\Lambda : \mathrm{Submodule} \mathbb{Z} V) : \mathrm{Submodule} \mathbb{Z} V
\]
\end{definition}

\begin{definition}[dual Equiv]\leanlink{Formalism.H.dualEquiv}
\begin{quote}\small\raggedright$\displaystyle : (\iota \to \allowbreak \mathbb{Z} ) \simeq \mathrm{dual} (\mathrm{span} \mathbb{Z} (\mathcal{P}.\mathrm{range} e))$\end{quote}
\end{definition}

\begin{definition}[torus Fun C]\leanlink{Formalism.H.torusFunC}
\begin{quote}\small\raggedright$\displaystyle (G : V \to \allowbreak \mathbb{C} ) \allowbreak (hG : \mathrm{Continuous} G) \allowbreak (\mathrm{hper} : \forall l \in \mathrm{span} \mathbb{Z} (\mathcal{P}.\mathrm{range} e), \allowbreak \forall x, \allowbreak G (l + x) = G x) : C(\mathrm{UnitAddTorus} \iota , \allowbreak \mathbb{C} )$\end{quote}
\end{definition}

\begin{definition}[Poisson Decay]\leanlink{Formalism.H.PoissonDecay}
\[
: \mathrm{Prop}
\]
\end{definition}

\begin{definition}[Lattice Fourier Series]\leanlink{Formalism.H.LatticeFourierSeries}
\[
: \mathrm{Prop}
\]
\end{definition}

\begin{definition}[Poisson Schwartz]\leanlink{Formalism.H.PoissonSchwartz}
\[
: \mathrm{Prop}
\]
\end{definition}

\subsection{Main Results}
\begin{theorem}[lattice fourier series]\leanlink{Formalism.H.lattice_fourier_series}
\[
(V : \mathrm{Type}) : \mathrm{LatticeFourierSeries} V
\]
\end{theorem}
\begin{proof}
The argument operates on the structural objects dual \leanlink{Formalism.H.dual}, dual Equiv \leanlink{Formalism.H.dualEquiv}, torus Fun C \leanlink{Formalism.H.torusFunC}, Lattice Fourier Series \leanlink{Formalism.H.LatticeFourierSeries}. Combining \emph{m Fourier Coeff torus Fun C} \leanlink{Formalism.H.mFourierCoeff_torusFunC}, \emph{torus Fun proj} \leanlink{Formalism.H.torusFun_proj}, \emph{integral fundamental Domain eq} \leanlink{Formalism.H.integral_fundamentalDomain_eq}, \emph{m Fourier torus Proj} \leanlink{Formalism.H.mFourier_torusProj}, \emph{coeff integrand periodic} \leanlink{Formalism.H.coeff_integrand_periodic} with \emph{dual Vec mem} \leanlink{Formalism.H.dualVec_mem} yields the stated result.
\end{proof}

\begin{theorem}[covolume Dual]\leanlink{Formalism.H.covolumeDual}
\[
\mathrm{CovolumeDual} V
\]
\end{theorem}
\begin{proof}
The argument operates on the structural objects dual \leanlink{Formalism.H.dual}, Covolume Dual \leanlink{Formalism.H.CovolumeDual}. Combining \emph{covol Dual det mul} \leanlink{Formalism.H.covolDual_det_mul}, \emph{exists basis span eq} \leanlink{Formalism.H.exists_basis_span_eq}, \emph{covol Dual vol fd} \leanlink{Formalism.H.covolDual_vol_fd} with \emph{covol Dual nondegenerate} \leanlink{Formalism.H.covolDual_nondegenerate} yields the stated result.
\end{proof}

\begin{proposition}[poisson periodization]\leanlink{Formalism.H.poisson_periodization}
\[
(V : \mathrm{Type}) : \mathrm{Periodization} V
\]
\end{proposition}
\begin{proof}
The argument operates on the structural objects Periodization \leanlink{Formalism.H.Periodization}. Combining \emph{exists global bound} \leanlink{Formalism.H.PoissonPeriodization.exists_global_bound}, \emph{summable one add norm} \leanlink{Formalism.H.PoissonPeriodization.summable_one_add_norm} with \emph{shift bound} \leanlink{Formalism.H.PoissonPeriodization.shift_bound} yields the stated result.
\end{proof}

\begin{proposition}[poisson Schwartz]\leanlink{Formalism.H.poissonSchwartz}
\[
\mathrm{PoissonSchwartz} V
\]
\end{proposition}
\begin{proof}
The argument operates on the structural objects Poisson Schwartz \leanlink{Formalism.H.PoissonSchwartz}. Combining \emph{poisson Decay} \leanlink{Formalism.H.poissonDecay}, \emph{schwartz decay gaussian} \leanlink{Formalism.H.schwartz_decay_gaussian} with \emph{poisson Schwartz of} \leanlink{Formalism.H.poissonSchwartz_of} yields the stated result.
\end{proof}

\begin{proposition}[poisson Decay]\leanlink{Formalism.H.poissonDecay}
\[
\mathrm{PoissonDecay} V
\]
\end{proposition}
\begin{proof}
The argument operates on the structural objects Poisson Decay \leanlink{Formalism.H.PoissonDecay}. Combining \emph{lattice fourier series} \leanlink{Formalism.H.lattice_fourier_series}, \emph{poisson periodization} \leanlink{Formalism.H.poisson_periodization}, \emph{dual Is Lattice} \leanlink{Formalism.H.dualIsLattice}, \emph{poisson Decay of} \leanlink{Formalism.H.poissonDecay_of}, \emph{lattice summable decay} \leanlink{Formalism.H.lattice_summable_decay} with \emph{unfolding fourier coefficient} \leanlink{Formalism.H.unfolding_fourier_coefficient} yields the stated result.
\end{proof}

\begin{proposition}[dual Is Lattice]\leanlink{Formalism.H.dualIsLattice}
\[
\mathrm{DualIsLattice} V
\]
\end{proposition}
\begin{proof}
The argument operates on the structural objects dual \leanlink{Formalism.H.dual}, Dual Is Lattice \leanlink{Formalism.H.DualIsLattice}. Combining \emph{dual Lattice innerl nondegenerate} \leanlink{Formalism.H.dualLattice_innerₗ_nondegenerate}, \emph{is Lattice of eq span} \leanlink{Formalism.H.isLattice_of_eq_span} with \emph{exists basis span eq} \leanlink{Formalism.H.exists_basis_span_eq} yields the stated result.
\end{proof}

\begin{lemma}[poisson Decay of]\leanlink{Formalism.H.poissonDecay_of}
\begin{quote}\small\raggedright$\displaystyle (hD : \mathrm{DualIsLattice} V) \allowbreak (h1 : \mathrm{LatticeSummableDecay} V) \allowbreak (h2 : \mathrm{Periodization} V) \allowbreak (h3 : \mathrm{UnfoldCoeff} V) \allowbreak (h4 : \mathrm{LatticeFourierSeries} V) : \mathrm{PoissonDecay} V$\end{quote}
\end{lemma}
\begin{proof}
The argument operates on the structural objects dual \leanlink{Formalism.H.dual}, Lattice Fourier Series \leanlink{Formalism.H.LatticeFourierSeries}, Poisson Decay \leanlink{Formalism.H.PoissonDecay}, Periodization \leanlink{Formalism.H.Periodization}. The claim follows directly from \emph{norm exp two pi I mul} \leanlink{Formalism.H.norm_exp_two_pi_I_mul}.
\end{proof}

\begin{lemma}[dual Dual]\leanlink{Formalism.H.dualDual}
\[
\mathrm{DualDual} V
\]
\end{lemma}
\begin{proof}
The argument operates on the structural objects Dual Dual \leanlink{Formalism.H.DualDual}. Combining \emph{innerl is Symm} \leanlink{Formalism.H.innerₗ_isSymm} with \emph{exists basis span eq} \leanlink{Formalism.H.exists_basis_span_eq} yields the stated result.
\end{proof}

\begin{lemma}[m Fourier Coeff torus Fun C]\leanlink{Formalism.H.mFourierCoeff_torusFunC}
\begin{quote}\small\raggedright$\displaystyle (G : V \to \allowbreak \mathbb{C} ) \allowbreak (hG : \mathrm{Continuous} G) \allowbreak (\mathrm{hper} : \forall l \in \mathrm{span} \mathbb{Z} (\mathcal{P}.\mathrm{range} e), \allowbreak \forall x, \allowbreak G (l + x) = G x) \allowbreak (m : \iota \to \allowbreak \mathbb{Z} ) : \mathrm{mFourierCoeff} (\mathrm{torusFunC} e G hG \mathrm{hper}) m = \int t in \mathrm{boxIco} \iota , \allowbreak G (\mathrm{symm} t) * \exp (-(2 * \pcReal .pi * i * (\mathrm{dualVec} e m, \allowbreak \mathrm{symm} t : \mathbb{R} )))$\end{quote}
\end{lemma}
\begin{proof}
The argument operates on the structural objects torus Fun C \leanlink{Formalism.H.torusFunC}, torus Fun \leanlink{Formalism.H.torusFun}, torus Proj \leanlink{Formalism.H.torusProj}, dual Vec \leanlink{Formalism.H.dualVec}. Combining \emph{torus Fun proj} \leanlink{Formalism.H.torusFun_proj}, \emph{m Fourier torus Proj} \leanlink{Formalism.H.mFourier_torusProj} with \emph{box Ioc ae eq} \leanlink{Formalism.H.boxIoc_ae_eq} yields the stated result.
\end{proof}

\begin{lemma}[torus Fun proj]\leanlink{Formalism.H.torusFun_proj}
\begin{quote}\small\raggedright$\displaystyle (G : V \to \allowbreak \mathbb{C} ) \allowbreak (\mathrm{hper} : \forall l \in \mathrm{span} \mathbb{Z} (\mathcal{P}.\mathrm{range} e), \allowbreak \forall x, \allowbreak G (l + x) = G x) \allowbreak (t : \iota \to \allowbreak \mathbb{R} ) : \mathrm{torusFun} e G (\mathrm{torusProj} \iota t) = G (\mathrm{symm} t)$\end{quote}
\end{lemma}
\begin{proof}
The argument operates on the structural objects torus Fun \leanlink{Formalism.H.torusFun}, torus Proj \leanlink{Formalism.H.torusProj}. Combining \emph{torus Proj eq} \leanlink{Formalism.H.torusProj_eq}, \emph{periodic comp} \leanlink{Formalism.H.periodic_comp} with \emph{torus Proj surjective} \leanlink{Formalism.H.torusProj_surjective} yields the stated result.
\end{proof}

\begin{lemma}[integral fundamental Domain eq]\leanlink{Formalism.H.integral_fundamentalDomain_eq}
\begin{quote}\small\raggedright$\displaystyle (H : V \to \allowbreak \mathbb{C} ) : \int y in \mathrm{fundamentalDomain} e, \allowbreak H y = (\mathrm{real} (\mathrm{fundamentalDomain} e) : \mathbb{C} ) * \int t in \mathrm{boxIco} \iota , \allowbreak H (\mathrm{symm} t)$\end{quote}
\end{lemma}
\begin{proof}
The argument operates on the structural objects box Ico \leanlink{Formalism.H.boxIco}. Combining \emph{volume box Ico} \leanlink{Formalism.H.volume_boxIco} with \emph{fundamental Domain eq preimage} \leanlink{Formalism.H.fundamentalDomain_eq_preimage} yields the stated result.
\end{proof}

\begin{lemma}[inner dual Vec basis]\leanlink{Formalism.H.inner_dualVec_basis}
\[
(m : \iota \to \mathbb{Z} ) (j : \iota ) : \mathrm{dualVec} e m, e j = m j
\]
\end{lemma}
\begin{proof}
The argument operates on the structural objects dual Vec \leanlink{Formalism.H.dualVec}. Combining \emph{inner dual B} \leanlink{Formalism.H.inner_dualB} with \emph{dual Vec eq} \leanlink{Formalism.H.dualVec_eq} yields the stated result.
\end{proof}

\section{Development / Abelian Variety Theory: se Good}

This chapter encompasses 337,931 formal declarations (4,621,552~LOC). Primary proof idioms detected: \emph{abelian varieties \& heights, Iwasawa tower, automorphic / spectral theory}.
\begin{quote}\small
\textbf{Formal Interface ($\mathrm{Inputs} \vdash \mathrm{Outputs}$):}
\par\noindent\emph{Inputs:} D7~\leanlink{Formalism.BasicNotions.D7}, Sp~\leanlink{ReciprocityTheory.ChartRestriction.G.Geo.Tor.Sp}, HSp~\leanlink{Formalism.BasicNotions.TH.HSp}, algebra Map~\leanlink{Formalism.GradedWeylDivisibility.CoeffHom.algebraMap}, Host~\leanlink{ReciprocityTheory.ChartRestriction.G.Geo.Tor.Host}, LB~\leanlink{ReciprocityTheory.ChartRestriction.G.Semi.LB}.
\par\noindent\emph{Outputs:} AE~\leanlink{GlobalTheory.IndefiniteRowSideFullHecke.WeilTensorModels.AE}, Spc~\leanlink{GlobalTheory.IndefiniteRowSideFullHecke.WeilTensorModels.Spc}, GA~\leanlink{GlobalTheory.CuspRep.GA}, Base~\leanlink{GlobalTheory.IndefiniteRowSideFullHecke.WeilTensorModels.Base}, Pl~\leanlink{GlobalTheory.BCT.Pl}, map mul~\leanlink{GlobalTheory.Serre.Lie.map_mul}.
\end{quote}

\subsection{Core Definitions and Structures}
\begin{definition}[cut N]\leanlink{GlobalTheory.HeightCounting.CutOff.cutN}
\[
\mathrm{TwNew} D
\]
\end{definition}

\begin{definition}[cols]\leanlink{GlobalTheory.Mem.Col2.TC.cols}
\[
\mathrm{Cols} B
\]
\end{definition}

\begin{definition}[bas]\leanlink{GlobalTheory.ThFr.Gauss.bas}
\[
(\mathrm{Fin} 2 \oplus \mathrm{Fin} 2 \to \mathbb{Z} ) \simeq _{l} \mathrm{flagData}.\Lambda
\]
\end{definition}

\begin{definition}[vec Lin]\leanlink{GlobalTheory.ThFr.Gauss.vecLin}
\[
(\mathrm{Fin} 2 \oplus \mathrm{Fin} 2 \to \mathbb{Z} ) \to _{l} \mathrm{flagData}.\Lambda
\]
\end{definition}

\begin{definition}[Moved]\leanlink{GlobalTheory.Mem.R3.Rw.Moved}
\[
(o : q.\pcOver ) : \mathrm{Type}1
\]
\end{definition}

\begin{definition}[Supp]\leanlink{GlobalTheory.Mem.Top.Bl.Supp}
\[
WP
\]
\end{definition}

\subsection{Main Results}
\begin{theorem}[se Good]\leanlink{GlobalTheory.HeightCounting.SerCopy.seGood}
\[
(h : \rho .\pcGood ) : (\mathrm{seT} o' \rho ).\pcGood
\]
\end{theorem}
\begin{proof}
The argument operates on the structural objects NData \leanlink{Formalism.BasicNotions.NData}, Rows Rest \leanlink{GlobalTheory.Mem.RowsRest}, NTw \leanlink{Formalism.BasicNotions.NTw}, Base \leanlink{GlobalTheory.Mem.Base}. Combining \emph{ba Good} \leanlink{GlobalTheory.HeightCounting.SerCopy.baGood}, \emph{tw Good} \leanlink{GlobalTheory.HeightCounting.SerCopy.twGood}, \emph{ba Weil} \leanlink{GlobalTheory.HeightCounting.SerCopy.baWeil}, \emph{rc Good} \leanlink{GlobalTheory.HeightCounting.SerCopy.rcGood} with \emph{weil} \leanlink{GlobalTheory.AbLaw.FromCube.weil} yields the stated result.
\end{proof}

\begin{theorem}[family Unif of laws]\leanlink{GlobalTheory.familyUnif_of_laws}
\[
\mathrm{Out}.\mathrm{Laws} \to \mathrm{N1c}
\]
\end{theorem}
\begin{proof}
The argument operates on the structural objects Laws \leanlink{ReciprocityTheory.ChartRestriction.G.Out.Laws}, Theta Unif AV \leanlink{ReciprocityTheory.ChartRestriction.G.Geo.Tor.ThetaUnifAV}, N1c \leanlink{ReciprocityTheory.ChartRestriction.G.Geo.Tor.N1c}. Combining \emph{family Uniformisation of abelian Variety uniformisation} \leanlink{GlobalTheory.T.familyUniformisation_of_abelianVariety_uniformisation} with \emph{theta Unif AV of laws} \leanlink{GlobalTheory.thetaUnifAV_of_laws} yields the stated result.
\end{proof}

\begin{theorem}[holds tr]\leanlink{GlobalTheory.Mem.R3.Tr.holds_tr}
\[
(h : \mathrm{Holds} R F) : \mathrm{letI}
\]
\end{theorem}
\begin{proof}
The argument operates on the structural objects real3 \leanlink{GlobalTheory.Mem.R3.Tr.real3}, Base \leanlink{GlobalTheory.Mem.Base}, prim \leanlink{Formalism.MainNotions.GL3Theta.prim}, Face3 \leanlink{Formalism.IndefR.Links.Face3}. Combining \emph{real Adelic Inert Transport further Rows} \leanlink{GlobalTheory.Mem.R3.Tr.realAdelicInertTransport_furtherRows}, \emph{prim Rows3 tr} \leanlink{GlobalTheory.Mem.R3.Tr.primRows3_tr}, \emph{ok} \leanlink{Formalism.L2.Col.GSD.Iso.ok}, \emph{ok} \leanlink{Formalism.L2.Col.CharT.WeilTied.ok} with \emph{sep} \leanlink{Formalism.DivisionValues.GluingInclusion.sep} yields the stated result.
\end{proof}

\begin{theorem}[detects]\leanlink{GlobalTheory.GeoTie.SideTie.detects}
\[
\mathrm{Detects} YS fw (T.\mathrm{iC} i) b j
\]
\end{theorem}
\begin{proof}
The argument operates on the structural objects D7 \leanlink{Formalism.BasicNotions.D7}, NTh \leanlink{Formalism.BasicNotions.NTh}, Rat Dat \leanlink{Formalism.IndefR.NX.RatDat}, Wt \leanlink{Formalism.MainNotions.WeightFamilies.Wt}. Combining \emph{ev inj} \leanlink{GlobalTheory.GeoTie.GeoLaws.ev_inj} with \emph{norm le iff} \leanlink{GlobalTheory.GeoTie.SideTie.norm_le_iff} yields the stated result.
\end{proof}

\begin{theorem}[hecke Std Plain]\leanlink{GlobalTheory.HeckeStdFacts.heckeStdPlain}
\[
\mathrm{HeckeStdPlain}
\]
\end{theorem}
\begin{proof}
The argument operates on the structural objects Hecke Std Plain \leanlink{Formalism.L2.HeckeStdPlain}, alg \leanlink{ReciprocityTheory.PGsh.TowerExample.alg}. Combining \emph{hecke D comm} \leanlink{GlobalTheory.HeckeStdFacts.heckeD_comm}, \emph{eigen integral} \leanlink{GlobalTheory.HeckeStdFacts.eigen_integral}, \emph{hecke D mem} \leanlink{GlobalTheory.HeckeStdFacts.heckeD_mem}, \emph{not dvd q} \leanlink{GlobalTheory.HeckeStdFacts.not_dvd_q}, \emph{exists data} \leanlink{GlobalTheory.HeckeStdFacts.exists_data}, \emph{q ne zero} \leanlink{GlobalTheory.HeckeStdFacts.q_ne_zero}, \emph{hecke D type} \leanlink{GlobalTheory.HeckeStdFacts.heckeD_type} with \emph{hecke D lin} \leanlink{GlobalTheory.HeckeStdFacts.heckeD_lin} yields the stated result.
\end{proof}

\begin{theorem}[vec Lin bijective]\leanlink{GlobalTheory.ThFr.Gauss.vecLin_bijective}
\[
\mathrm{Bijective} \mathrm{vecLin}
\]
\end{theorem}
\begin{proof}
The argument operates on the structural objects vec Lin \leanlink{GlobalTheory.ThFr.Gauss.vecLin}, Data \leanlink{ReciprocityTheory.ChartRestriction.G.FD.Data}, vec \leanlink{GlobalTheory.ThFr.Gauss.vec}, herm \leanlink{ReciprocityTheory.ChartRestriction.G.FD.herm}. Combining \emph{flag Data } \leanlink{GlobalTheory.HyperbolicCusp.flagData_Λ}, \emph{EF vec} \leanlink{GlobalTheory.ThFr.Gauss.EF_vec} with \emph{exists int of is Integral} \leanlink{GlobalTheory.HyperbolicCusp.exists_int_of_isIntegral} yields the stated result.
\end{proof}

\begin{theorem}[real Adelic Inert Transport further Rows]\leanlink{GlobalTheory.Mem.R3.Tr.realAdelicInertTransport_furtherRows}
\[
(h : \mathrm{UnusedChoiceStatements} R F) : \mathrm{letI}
\]
\end{theorem}
\begin{proof}
The argument operates on the structural objects real3 \leanlink{GlobalTheory.Mem.R3.Tr.real3}, Base \leanlink{GlobalTheory.Mem.Base}, Face3 \leanlink{Formalism.IndefR.Links.Face3}, Unused Choice Statements \leanlink{Formalism.IndefR.Links.UnusedChoiceStatements}. Combining \emph{ex At tr} \leanlink{GlobalTheory.Mem.R3.Tr.exAt_tr} with \emph{l Level W tr} \leanlink{GlobalTheory.Mem.R3.Tr.lLevelW_tr} yields the stated result.
\end{proof}

\begin{theorem}[mult eq]\leanlink{GlobalTheory.Mem.R3.Tr.mult_eq}
\[
@\mathrm{Mult} b.\mathrm{D} = @\mathrm{multB} B C a b
\]
\end{theorem}
\begin{proof}
The argument operates on the structural objects Base \leanlink{GlobalTheory.Mem.Base}, J \leanlink{Formalism.BasicNotions.TH.FockData.J}, Cols \leanlink{GlobalTheory.Mem.Cols}, Mult \leanlink{Formalism.IndefR.Mult.Mult}. This result is established directly from standard library foundations within its local module (5~lines of formal proof).
\end{proof}

\begin{theorem}[same Std]\leanlink{GlobalTheory.Mem.SRe.CarModel.sameStd}
\[
\mathrm{SameStd}
\]
\end{theorem}
\begin{proof}
The argument operates on the structural objects Same Std \leanlink{GlobalTheory.Mem.SRe.SameStd}, Col W \leanlink{GlobalTheory.Mem.Top.Bl.ColW}, car \leanlink{GlobalTheory.Mem.SRe.CarModel.car}, Base \leanlink{GlobalTheory.Mem.Base}. The claim follows directly from \emph{car eq} \leanlink{GlobalTheory.Mem.SRe.CarModel.car_eq}.
\end{proof}

\begin{theorem}[flag Data Gm]\leanlink{GlobalTheory.HyperbolicCusp.flagData_Gm}
\[
\mathrm{Gm} = gm
\]
\end{theorem}
\begin{proof}
The argument operates on the structural objects flag Data \leanlink{GlobalTheory.HyperbolicCusp.flagData}, gm \leanlink{GlobalTheory.HyperbolicCusp.gm}. The claim follows directly from \emph{exists flag Data} \leanlink{GlobalTheory.HyperbolicCusp.exists_flagData}.
\end{proof}

\begin{theorem}[harish Chandra Finiteness Totally Complex GL proof]\leanlink{GlobalTheory.T.harishChandraFinitenessTotallyComplexGL_proof}
\[
\mathrm{CoreFiniteGL}
\]
\end{theorem}
\begin{proof}
This result is established directly from standard library foundations within its local module (2~lines of formal proof).
\end{proof}

\begin{theorem}[in Mt state]\leanlink{GlobalTheory.MtId.inMt_state}
\begin{quote}\small\raggedright$\displaystyle (\mathrm{pin} : \mathrm{ColPin} w) : \exists m : \mathrm{InMt} \mathrm{Y} (\mathrm{cfOf} x) \mathrm{xi} (\mathrm{QOf} P), \allowbreak \forall i, \allowbreak \mathrm{IdFrame} \mathrm{hF} (\mathrm{PV} i)$\end{quote}
\end{theorem}
\begin{proof}
The argument operates on the structural objects Col Pin \leanlink{GlobalTheory.Mem.Top.Bl.ColPin}, Pre2 \leanlink{GlobalTheory.Mem.R3.Rw.Pre2}, Col W \leanlink{GlobalTheory.Mem.Top.Bl.ColW}, Dof \leanlink{GlobalTheory.HeightCounting.Up.Dof}. Combining \emph{real state} \leanlink{GlobalTheory.MtId.real_state}, \emph{exists unit frame} \leanlink{GlobalTheory.MtIn.exists_unit_frame}, \emph{gram V} \leanlink{GlobalTheory.Mem.R3.StdModel.gram_V} with \emph{gen state} \leanlink{GlobalTheory.MtId.gen_state} yields the stated result.
\end{proof}

\section{Reciprocity Theory: segment Representation}

This chapter encompasses 163,457 formal declarations (2,120,740~LOC). Primary proof idioms detected: \emph{abelian varieties \& heights, automorphic / spectral theory, Hermitian curvature / L\textasciicircum{}2 dbar}.
\begin{quote}\small
\textbf{Formal Interface ($\mathrm{Inputs} \vdash \mathrm{Outputs}$):}
\par\noindent\emph{Inputs:} AF~\leanlink{Formalism.MainNotions.GlobalTheoryTheory.GenusMassVolumes.AF}, D7~\leanlink{Formalism.BasicNotions.D7}, algebra Map~\leanlink{Formalism.GradedWeylDivisibility.CoeffHom.algebraMap}, Spc~\leanlink{GlobalTheory.IndefiniteRowSideFullHecke.WeilTensorModels.Spc}, AE~\leanlink{GlobalTheory.IndefiniteRowSideFullHecke.WeilTensorModels.AE}, GA~\leanlink{GlobalTheory.CuspRep.GA}.
\par\noindent\emph{Outputs:} Sp~\leanlink{ReciprocityTheory.ChartRestriction.G.Geo.Tor.Sp}, Host~\leanlink{ReciprocityTheory.ChartRestriction.G.Geo.Tor.Host}, Datum~\leanlink{ReciprocityTheory.ChartRestriction.G.Geo.Tor.Datum}, LB~\leanlink{ReciprocityTheory.ChartRestriction.G.Semi.LB}, sp Map~\leanlink{ReciprocityTheory.ChartRestriction.G.Geo.Tor.spMap}, Level~\leanlink{ReciprocityTheory.ChartRestriction.G.Geo.Tor.Level}.
\end{quote}

\subsection{Core Definitions and Structures}
\begin{definition}[Run Frame]\leanlink{PositivityTheory.Rg.BcB.F.RunFrame}
\[
(\mathrm{F\_0} : \mathrm{Frame0} hR \pi S) (T : \mathcal{F} (Pl E)) : \mathrm{Type}1
\]
\end{definition}

\begin{definition}[ug]\leanlink{PositivityTheory.Rg.BcB.F.RunFrame.ug}
\[
F.\Gamma \to U.\mathrm{UGen} E n
\]
\end{definition}

\begin{definition}[rad Int]\leanlink{GlobalTheory.CompanionTheory.Rodier.RadInt.radInt}
\begin{quote}\small\raggedright$\displaystyle (h\psi : \mathrm{Continuous} \psi ) \allowbreak (h\rho : \mathrm{IsSmooth} \rho ) \allowbreak (\mu : \mathrm{Measure} (\mathrm{unipU} k d hd)) : \mathrm{openFns} d hd \rho \to \allowbreak _{l} W$\end{quote}
\end{definition}

\begin{definition}[Conj Pl]\leanlink{PositivityTheory.Rg.BcB.F.ConjPl}
\[
(w u : \mathrm{HeightOneSpectrum} (\mathcal{O} E)) : \mathrm{Prop}
\]
\end{definition}

\begin{definition}[Frame0]\leanlink{PositivityTheory.Rg.BcB.F.Frame0}
\begin{quote}\small\raggedright$\displaystyle (hR : Rd.\mathrm{ReadLimit}) \allowbreak (\pi : \mathrm{Submodule} \mathbb{C} (UA E n \to \allowbreak \mathbb{C} )) \allowbreak (S : \mathcal{F} (Pl E)) : \mathrm{Type}1$\end{quote}
\end{definition}

\begin{definition}[]\leanlink{PositivityTheory.Rg.BcB.F.RunFrame.Γ}
\[
\mathrm{Type}
\]
\end{definition}

\subsection{Main Results}
\begin{theorem}[segment Representation]\leanlink{PositivityTheory.RgGu.segmentRepresentation}
\[
\mathrm{SegmentRepresentation}
\]
\end{theorem}
\begin{proof}
The argument operates on the structural objects Segment Representation \leanlink{GlobalTheory.Published.SegmentRepresentation}, Geometric Lemma \leanlink{GlobalTheory.Published.GeometricLemma}. Combining \emph{rho Nu Rho Reducible C} \leanlink{PositivityTheory.RgGu.NR.rhoNuRhoReducibleC}, \emph{stage P} \leanlink{PositivityTheory.RgGu.SRC.stageP}, \emph{seg Rep of} \leanlink{PositivityTheory.RgGu.SRC.segRep_of}, \emph{seg Two of} \leanlink{PositivityTheory.RgGu.SRC.segTwo_of}, \emph{geometric Lemma of filtration} \leanlink{GlobalTheory.geometricLemma_of_filtration}, \emph{bernstein Zelevinsky Finiteness} \leanlink{GlobalTheory.CompanionTheory.Published.bernsteinZelevinskyFiniteness} with \emph{geometric Lemma Filtration} \leanlink{PositivityTheory.Geo.geometricLemmaFiltration} yields the stated result.
\end{proof}

\begin{theorem}[red Cl C]\leanlink{PositivityTheory.RgGu.NR.redClC}
\[
\mathrm{RedCl} k
\]
\end{theorem}
\begin{proof}
Combining \emph{cusp Pair Irr} \leanlink{PositivityTheory.RgGu.NR.cuspPairIrr}, \emph{signature} \leanlink{PositivityTheory.RgGu.NR.signature}, \emph{complementary of} \leanlink{PositivityTheory.RgGu.NR.complementary_of}, \emph{conj Dual} \leanlink{PositivityTheory.RgGu.NR.conjDual}, \emph{some Red of} \leanlink{PositivityTheory.RgGu.NR.someRed_of}, \emph{herm of B} \leanlink{PositivityTheory.RgGu.NR.herm_ofB}, \emph{red Cl of C} \leanlink{PositivityTheory.RgGu.NR.redCl_ofC} with \emph{prod Model} \leanlink{PositivityTheory.RgGu.NR.prodModel} yields the stated result.
\end{proof}

\begin{theorem}[power Sum Hom injective]\leanlink{GroundTheory.Heights₂.Heights.RestTerms.Flag.powerSumHom_injective}
\[
(N : \mathbb{N} ) : \mathrm{Injective} (\mathrm{powerSumHom} N)
\]
\end{theorem}
\begin{proof}
The argument operates on the structural objects power Sum Hom \leanlink{GroundTheory.Heights₂.Heights.RestTerms.powerSumHom}. Combining \emph{injective of surjective endo} \leanlink{GroundTheory.Heights₂.Heights.RestTerms.Flag.injective_of_surjective_endo}, \emph{exists power Sum Hom eq} \leanlink{GroundTheory.Heights₂.Heights.RestTerms.Flag.exists_powerSumHom_eq} with \emph{power Sum Hom is Symmetric} \leanlink{GroundTheory.Heights₂.Heights.RestTerms.Flag.powerSumHom_isSymmetric} yields the stated result.
\end{proof}

\begin{theorem}[grothendieck Ring Polynomial of seg]\leanlink{PositivityTheory.Groth.grothendieckRingPolynomial_of_seg}
\[
(hZ : \mathrm{SegmentRepresentation}) : \mathrm{GrothendieckRingPolynomial}
\]
\end{theorem}
\begin{proof}
The argument operates on the structural objects Grothendieck Ring Polynomial \leanlink{GlobalTheory.Published.GrothendieckRingPolynomial}, Segment Representation \leanlink{GlobalTheory.Published.SegmentRepresentation}. Combining \emph{l Unlinked Irr of seg} \leanlink{PositivityTheory.Exh.lUnlinkedIrr_of_seg}, \emph{mid P'} \leanlink{PositivityTheory.Exh.midP'}, \emph{swap P'} \leanlink{PositivityTheory.Exh.swapP'}, \emph{groth of three} \leanlink{PositivityTheory.Groth.groth_of_three}, \emph{refine P'} \leanlink{PositivityTheory.Exh.refineP'}, \emph{irreducible Admissible GL} \leanlink{PositivityTheory.Exh.irreducibleAdmissibleGL} with \emph{comm of} \leanlink{PositivityTheory.Groth.ClsOfF.comm_of} yields the stated result.
\end{proof}

\begin{theorem}[cusp Pair Irr]\leanlink{PositivityTheory.RgGu.NR.cuspPairIrr}
\[
\mathrm{CuspPairIrr} k
\]
\end{theorem}
\begin{proof}
The argument operates on the structural objects Cusp Pair Irr \leanlink{GlobalTheory.BZSeg.CuspPairIrr}. Combining \emph{cusp Pair Irr of} \leanlink{PositivityTheory.RgGu.NR.cuspPairIrr_of} with \emph{nu Pair} \leanlink{PositivityTheory.RgGu.NR.nuPair} yields the stated result.
\end{proof}

\begin{theorem}[fin Rat F mem of]\leanlink{PositivityTheory.ThKW.AutLatA.finRatF_mem_of}
\[
(hy : y \in \Lambda ) : \mathrm{finRatF} F y \in \mathrm{latF} F E nS nT \Lambda
\]
\end{theorem}
\begin{proof}
The argument operates on the structural objects Base \leanlink{GlobalTheory.IndefiniteRowSideFullHecke.WeilTensorModels.Base}, theta Sub F \leanlink{GlobalTheory.WeilPartsF.thetaSubF}, fin Rat F \leanlink{PositivityTheory.ThKW.finRatF}, lat F \leanlink{GlobalTheory.FinLatF.latF}. The claim follows directly from \emph{one int} \leanlink{PositivityTheory.ThKW.AutLatA.one_int}.
\end{proof}

\begin{theorem}[whit Heredity]\leanlink{GlobalTheory.CompanionTheory.Rodier.whitHeredity}
\[
\forall (k : \mathrm{Type}) , \mathrm{WhitHeredity} k
\]
\end{theorem}
\begin{proof}
The argument operates on the structural objects open Fns \leanlink{GlobalTheory.CompanionTheory.Rodier.openFns}, Whit Heredity \leanlink{GlobalTheory.BZMir.Leib.WhitHeredity}, whit Char \leanlink{GlobalTheory.LocGL.whitChar}, multi Prod \leanlink{GlobalTheory.SmGL.multiProd}. Combining \emph{induced function open cell modulo whittaker} \leanlink{GlobalTheory.T.induced_function_open_cell_modulo_whittaker}, \emph{open cell jacquet map surjective whittaker} \leanlink{GlobalTheory.T.open_cell_jacquet_map_surjective_whittaker} with \emph{is Smooth pi Tensor} \leanlink{GlobalTheory.SmGL.isSmooth_piTensor} yields the stated result.
\end{proof}

\begin{theorem}[ug ok]\leanlink{PositivityTheory.Rg.BcB.F.B3.PTop.ug_ok}
\[
(x : F.\Gamma ) : (F.\mathrm{ug} x).Ok \mathrm{F\_0} .\mathrm{Sec} \mathrm{B}
\]
\end{theorem}
\begin{proof}
The argument operates on the structural objects Run Frame \leanlink{PositivityTheory.Rg.BcB.F.RunFrame}, ug \leanlink{PositivityTheory.Rg.BcB.F.RunFrame.ug}, Frame0 \leanlink{PositivityTheory.Rg.BcB.F.Frame0},  \leanlink{PositivityTheory.Rg.BcB.F.RunFrame.Γ}. Combining \emph{good notin B} \leanlink{PositivityTheory.Rg.BcB.F.B3.PTop.good_notin_B}, \emph{ok} \leanlink{Formalism.L2.Col.GSD.Iso.ok}, \emph{ok} \leanlink{Formalism.MainNotions.S.TieN.JtOrbOKN.ok} with \emph{is Hecke h U} \leanlink{GlobalTheory.CompanionTheory.MatchWord.isHecke_hU} yields the stated result.
\end{proof}

\begin{theorem}[rho Nu Rho Reducible C]\leanlink{PositivityTheory.RgGu.NR.rhoNuRhoReducibleC}
\[
\mathrm{RhoNuRhoReducible}
\]
\end{theorem}
\begin{proof}
The argument operates on the structural objects Rho Nu Rho Reducible \leanlink{PositivityTheory.Pf1.SegRepC.SegTwo.RhoNuRhoReducible}. Combining \emph{red Cl C} \leanlink{PositivityTheory.RgGu.NR.redClC} with \emph{rho Nu Rho of C} \leanlink{PositivityTheory.RgGu.NR.rhoNuRho_ofC} yields the stated result.
\end{proof}

\begin{lemma}[is Locally Constant phi]\leanlink{GlobalTheory.CompanionTheory.Rodier.RadInt.isLocallyConstant_phi}
\begin{quote}\small\raggedright$\displaystyle (h\psi : \mathrm{Continuous} \psi ) \allowbreak (f : \mathrm{IndSp} (\mathrm{stdBlocks} d hd) \rho ) : \mathrm{IsLocallyConstant} (\mathrm{phi} \psi f)$\end{quote}
\end{lemma}
\begin{proof}
The argument operates on the structural objects whit Char \leanlink{GlobalTheory.LocGL.whitChar}, par Ind \leanlink{GlobalTheory.SmGL.parInd}, theta U \leanlink{GlobalTheory.CompanionTheory.Rodier.RadInt.thetaU}, Ind Sp \leanlink{GlobalTheory.SmGL.IndSp}. Combining \emph{one mem} \leanlink{GlobalTheory.ThFr.Gauss.one_mem}, \emph{is Smooth Char whit Char} \leanlink{GlobalTheory.Whittaker.isSmoothChar_whitChar}, \emph{is Locally Constant of right} \leanlink{GlobalTheory.CompanionTheory.Rodier.RadInt.isLocallyConstant_of_right} with \emph{is Smooth par Ind} \leanlink{GlobalTheory.SmGL.isSmooth_parInd} yields the stated result.
\end{proof}

\begin{lemma}[param pf]\leanlink{PositivityTheory.Rg.BcB.G3.param_pf}
\begin{quote}\small\raggedright$\displaystyle (h1 : \mathrm{SkewEx}) \allowbreak (h2 : \mathrm{Struct}) \allowbreak (h3 : \mathrm{QuadField}) \allowbreak (h4 : \mathrm{LocDict}) \allowbreak (h5 : \mathrm{ClsNorm}) \allowbreak (h6 : \mathrm{ClsOnto}) \allowbreak (h7 : \mathrm{ClsEq}) \allowbreak (h8 : \mathrm{LocGood}) : \mathrm{Param}$\end{quote}
\end{lemma}
\begin{proof}
The argument operates on the structural objects Loc Good \leanlink{PositivityTheory.Rg.BcB.G3.LocGood}, Cls Norm \leanlink{PositivityTheory.Rg.BcB.G3.ClsNorm}, Cls Onto \leanlink{PositivityTheory.Rg.BcB.G3.ClsOnto}, Cls Eq \leanlink{PositivityTheory.Rg.BcB.G3.ClsEq}. The claim follows directly from \emph{param main} \leanlink{PositivityTheory.Rg.BcB.G3.ParamP.param_main}.
\end{proof}

\begin{lemma}[param main]\leanlink{PositivityTheory.Rg.BcB.G3.ParamP.param_main}
\begin{quote}\small\raggedright$\displaystyle (h1 : \mathrm{SkewEx}) \allowbreak (h2 : \mathrm{Struct}) \allowbreak (h3 : \mathrm{QuadField}) \allowbreak (h4 : \mathrm{LocDict}) \allowbreak (h5 : \mathrm{ClsNorm}) \allowbreak (h6 : \mathrm{ClsOnto}) \allowbreak (h7 : \mathrm{ClsEq}) \allowbreak (h8 : \mathrm{LocGood}) : \mathrm{Param}$\end{quote}
\end{lemma}
\begin{proof}
The argument operates on the structural objects under \leanlink{GlobalTheory.Ker.Ctx.under}. Combining \emph{loc El rel} \leanlink{PositivityTheory.Rg.BcB.G3.ParamP.locEl_rel}, \emph{charpoly nrm L} \leanlink{GlobalTheory.BCT.ThBasic.charpoly_nrmL}, \emph{exists over} \leanlink{PositivityTheory.Rg.BcB.G3.ParamP.exists_over}, \emph{exists norm} \leanlink{GlobalTheory.BCT.KapMul.exists_norm}, \emph{map inv} \leanlink{Formalism.L2.DIV.IsAbsIrrChar.map_inv}, \emph{tw R equiv} \leanlink{GlobalTheory.BCT.FinFib.twR_equiv}, \emph{ne zero} \leanlink{Formalism.MainNotions.G.HEv.ne_zero} with \emph{loc El eventually mem} \leanlink{GlobalTheory.BCT.FinFib.locEl_eventually_mem} yields the stated result.
\end{proof}

\section{Formalism: gluing mass holds}

This chapter encompasses 173,144 formal declarations (1,931,701~LOC). Primary proof idioms detected: \emph{Harish-Chandra / real reductive groups, automorphic / spectral theory, p-adic / local harmonic analysis}.
\begin{quote}\small
\textbf{Formal Interface ($\mathrm{Inputs} \vdash \mathrm{Outputs}$):}
\par\noindent\emph{Inputs:} Host~\leanlink{ReciprocityTheory.ChartRestriction.G.Geo.Tor.Host}, Core~\leanlink{ReciprocityTheory.ChartRestriction.T2.Core}, Datum~\leanlink{ReciprocityTheory.ChartRestriction.G.Geo.Tor.Datum}, Moduli~\leanlink{ReciprocityTheory.ChartRestriction.G.Geo.Tor.Moduli}, Coeff~\leanlink{ReciprocityTheory.ChartRestriction.G.Geo.Tor.Coeff}, G~\leanlink{Formalism.ConjugateSelfDual.G}.
\par\noindent\emph{Outputs:} D7~\leanlink{Formalism.BasicNotions.D7}, algebra Map~\leanlink{Formalism.GradedWeylDivisibility.CoeffHom.algebraMap}, HSp~\leanlink{Formalism.BasicNotions.TH.HSp}, AF~\leanlink{Formalism.MainNotions.GlobalTheoryTheory.GenusMassVolumes.AF}, X~\leanlink{Formalism.BasicNotions.LocWach.X}, GA~\leanlink{Formalism.L2.Theta.GA}.
\end{quote}

\subsection{Core Definitions and Structures}
\begin{definition}[Mram H]\leanlink{Formalism.L2.Col.MramH}
\[
(u : Pr.\mathrm{Word}) : \mathrm{HLat} X.\mathrm{K}
\]
\end{definition}

\begin{definition}[mem Q]\leanlink{Formalism.L2.DIV.memQ}
\[
\mathrm{KerChain} X.\mathrm{K}
\]
\end{definition}

\begin{definition}[ram Dim]\leanlink{Formalism.L2.DIV.ramDim}
\[
(u : Pr.\mathrm{Word}) : \mathbb{N}
\]
\end{definition}

\begin{definition}[First Cols]\leanlink{Formalism.L2.DIV.FirstCols}
\[
(T : \mathrm{ChainTie} ch (\mathrm{memQ} (A := A) (hF := hF) hE (x := x))) : \mathrm{Prop}
\]
\end{definition}

\begin{definition}[col Of Coord]\leanlink{Formalism.L2.DIV.colOfCoord}
\begin{quote}\small\raggedright$\displaystyle (u : Pr.\mathrm{Word}) : (\mathrm{Fin} (\mathrm{ramDim} W Ch hX u) \to \allowbreak X.\mathrm{K}) \to \allowbreak _{l} \mathrm{MfV} W Ch u$\end{quote}
\end{definition}

\begin{definition}[col Per]\leanlink{Formalism.L2.DIV.colPer}
\begin{quote}\small\raggedright$\displaystyle (hP : (\mathrm{mtildeGSD} W Ch \mathrm{sel} mu P).Ok) \allowbreak (Xg : \mathcal{F} (G0 X.\mathrm{E} Pr.\mathrm{vGs} Rk.\mathrm{r} \to \allowbreak X.\mathrm{K})) \allowbreak (A : \mathrm{ChKer} W Ch \mathrm{sel} mu \mathrm{cst} L P u m2 b hY) \allowbreak (hF : \forall j, \allowbreak (\mathrm{chF} W \mathrm{sel} mu L u j).Ok) \allowbreak (hL : L.\mathrm{Ok} Ch \mathrm{sel} mu \mathrm{cst} G m2) \allowbreak (ev : (\mathrm{Fin} n \to \allowbreak X.\mathrm{K}) \to \allowbreak _{l} \mathrm{V0X} W Ch \mathrm{sel} mu Xg P) : (\mathrm{Fin} (\mathrm{ramDim} W Ch hX u) \to \allowbreak X.\mathrm{K}) \to \allowbreak _{l} \mathrm{Dual} X.\mathrm{K} (\mathrm{Fin} n \to \allowbreak X.\mathrm{K})$\end{quote}
\end{definition}

\subsection{Main Results}
\begin{theorem}[gluing mass holds]\leanlink{Formalism.L2.SG.Pf.Cnt.gluing_mass_holds}
\[
\mathrm{GlueMass}
\]
\end{theorem}
\begin{proof}
The argument operates on the structural objects Glue Mass \leanlink{Formalism.L2.SG.Pf.Cnt.GlueMass}. Combining \emph{alpha discriminant holds} \leanlink{Formalism.L2.SG.Pf.Cnt.alpha_discriminant_holds}, \emph{glue Mass of'} \leanlink{Formalism.L2.SG.Pf.Cnt.glueMass_of'}, \emph{jordan invariants holds} \leanlink{Formalism.L2.SG.Pf.Cnt.jordan_invariants_holds}, \emph{jordan splitting exists} \leanlink{Formalism.L2.SG.Pf.Cnt.jordan_splitting_exists}, \emph{local gluing holds} \leanlink{Formalism.L2.SG.Pf.Cnt.local_gluing_holds} with \emph{gluing product holds} \leanlink{Formalism.L2.SG.Pf.Cnt.gluing_product_holds} yields the stated result.
\end{proof}

\begin{theorem}[alpha discriminant holds]\leanlink{Formalism.L2.SG.Pf.Cnt.alpha_discriminant_holds}
\[
\mathrm{AlphaDisc}
\]
\end{theorem}
\begin{proof}
The argument operates on the structural objects Alpha Disc \leanlink{Formalism.L2.SG.Pf.Cnt.AlphaDisc}. Combining \emph{alpha Disc of disc Onto} \leanlink{Formalism.L2.SG.Pf.G5.alphaDisc_of_discOnto}, \emph{unitary discriminant model} \leanlink{Formalism.L2.SG.Pf.Cnt.unitary_discriminant_model} with \emph{disc Onto of u DModel} \leanlink{Formalism.L2.SG.Pf.discriminantNormShapes.discOnto_of_uDModel} yields the stated result.
\end{proof}

\begin{theorem}[li Stat]\leanlink{Formalism.HermitianLattices.Ext.L3.liStat}
\[
\mathrm{LIStat}
\]
\end{theorem}
\begin{proof}
The argument operates on the structural objects LIStat \leanlink{Formalism.HermitianLattices.Ext.L3.LIStat}. Combining \emph{li Ratio} \leanlink{Formalism.HermitianLattices.Ext.L3.liRatio}, \emph{li Pin} \leanlink{Formalism.HermitianLattices.Ext.L3.liPin} with \emph{li Stat of} \leanlink{Formalism.HermitianLattices.Ext.L3.liStat_of} yields the stated result.
\end{proof}

\begin{theorem}[class Reps Definite]\leanlink{Formalism.ThetaLiftIndependence.classRepsDefinite}
\[
\mathrm{ClassRepsDefinite}
\]
\end{theorem}
\begin{proof}
The argument operates on the structural objects isom \leanlink{Formalism.HermitianLatticesTheory.isom}, Class Reps Definite \leanlink{Formalism.ThetaLiftIndependence.ClassRepsDefinite}, aut L \leanlink{Formalism.ThetaLiftIndependence.autL}, genus \leanlink{Formalism.HermitianLatticesTheory.genus}. Combining \emph{exists reps} \leanlink{Formalism.ThetaLiftIndependence.DefiniteClassRepresentatives.exists_reps}, \emph{lat Wt gmap} \leanlink{Formalism.ThetaLiftIndependence.LatticeLiftIndependence.latWt_gmap} with \emph{gmap one} \leanlink{Formalism.ThetaLiftIndependence.LatticeLiftIndependence.gmap_one} yields the stated result.
\end{proof}

\begin{theorem}[row lattice mass]\leanlink{PI.UnitaryGroupMassFormula.CntFormula.RowsMassW.row_lattice_mass}
\[
\mathrm{RowsMass}
\]
\end{theorem}
\begin{proof}
Combining \emph{row class holds} \leanlink{PI.UnitaryGroupMassFormula.CntFormula.RowsMassW.row_class_holds}, \emph{row weight formula} \leanlink{PI.UnitaryGroupMassFormula.CntFormula.RowsMassW.row_weight_formula}, \emph{row lattice mass holds} \leanlink{PI.UnitaryGroupMassFormula.CntFormula.RowsMassW.row_lattice_mass_holds}, \emph{glued lattice mass formula} \leanlink{PI.UnitaryGroupMassFormula.CntFormula.RowsMassW.GluedLatticeMass.glued_lattice_mass_formula}, \emph{Mof unimod} \leanlink{PI.UnitaryGroupMassFormula.CntFormula.RowsMassW.RowMPf.Mof_unimod}, \emph{mem} \leanlink{Formalism.ThetaLiftIndependence.ChainScalarLemmas.mem}, \emph{mem} \leanlink{Formalism.HermitianLattices.Ext.L3.U0.Dv.mem} with \emph{anisotropic complement definite} \leanlink{PI.UnitaryGroupMassFormula.CntFormula.RowsMassW.anisotropic_complement_definite} yields the stated result.
\end{proof}

\begin{theorem}[mass Basic]\leanlink{Formalism.HermitianLatticesTheory.massBasic}
\[
\mathrm{MassBasic}
\]
\end{theorem}
\begin{proof}
The argument operates on the structural objects Mass Basic \leanlink{Formalism.HermitianLatticesTheory.MassBasic}, aut L \leanlink{Formalism.ThetaLiftIndependence.autL}, lat Wt \leanlink{Formalism.ThetaLiftIndependence.latWt}. Combining \emph{class Reps Definite} \leanlink{Formalism.ThetaLiftIndependence.classRepsDefinite}, \emph{lat Lift Indep} \leanlink{Formalism.ThetaLiftIndependence.latLiftIndep}, \emph{mass Basic self mem genus} \leanlink{Formalism.HermitianLatticesTheory.massBasic_self_mem_genus} with \emph{mass Basic mass eq} \leanlink{Formalism.HermitianLatticesTheory.massBasic_mass_eq} yields the stated result.
\end{proof}

\begin{theorem}[normalizing arithmetic holds]\leanlink{Formalism.L2.SG.Pf.NcGlob.normalizing_arithmetic_holds}
\[
\mathrm{NcArith}
\]
\end{theorem}
\begin{proof}
The argument operates on the structural objects Nc Arith \leanlink{Formalism.L2.SG.Pf.NcVal.NcArith}, R \leanlink{Formalism.MainNotions.S.Loc.Model.R}, h P \leanlink{Formalism.L2.DIV.NumDatum.hP}. Combining \emph{nc Value C of global} \leanlink{Formalism.L2.SG.Pf.NcArithG.ncValueC_of_global}, \emph{mem} \leanlink{Formalism.ThetaLiftIndependence.ChainScalarLemmas.mem}, \emph{mem} \leanlink{Formalism.HermitianLattices.Ext.L3.U0.Dv.mem}, \emph{field dictionary holds} \leanlink{Formalism.L2.SG.Pf.Cnt.field_dictionary_holds}, \emph{mem} \leanlink{Formalism.HermitianLattices.Ext.L3.WCA.InL.mem}, \emph{mem} \leanlink{Formalism.MainNotions.G.KapIDef.mem} with \emph{algebra Map} \leanlink{Formalism.GradedWeylDivisibility.CoeffHom.algebraMap} yields the stated result.
\end{proof}

\begin{theorem}[phi counting holds]\leanlink{Formalism.L2.SG.Pf.LocalizedFibreSets.phi_counting_holds}
\[
\mathrm{PhiCount}
\]
\end{theorem}
\begin{proof}
The argument operates on the structural objects Phi Count \leanlink{Formalism.L2.SG.Pf.G5.PhiCount}, Is Fix F \leanlink{Formalism.HermitianLattices.Ext.L3.IsFixF}. Combining \emph{phi mod p holds} \leanlink{Formalism.L2.SG.Pf.LocalizedFibreSets.phi_mod_p_holds} with \emph{phi Count of} \leanlink{Formalism.L2.SG.Pf.LocalizedFibreSets.phiCount_of} yields the stated result.
\end{proof}

\begin{theorem}[phi mod p holds]\leanlink{Formalism.L2.SG.Pf.LocalizedFibreSets.phi_mod_p_holds}
\[
\mathrm{PhiModP}
\]
\end{theorem}
\begin{proof}
The argument operates on the structural objects Phi Mod P \leanlink{Formalism.L2.SG.Pf.LocalizedFibreSets.PhiModP}. Combining \emph{card U diag} \leanlink{Formalism.L2.SG.Pf.LocalizedFibreSets.card_U_diag}, \emph{card from Blocks cayley eq mul} \leanlink{Formalism.L2.SG.Pf.LocalizedFibreSets.card_fromBlocks_cayley_eq_mul}, \emph{card ker pi Q} \leanlink{Formalism.HermitianLattices.Ext.L3.idealPowerSpanBounds.card_ker_piQ}, \emph{v F ne zero iff} \leanlink{Formalism.L2.SG.Pf.LocalizedFibreSets.vF_ne_zero_iff}, \emph{star O } \leanlink{Formalism.L2.SG.Pf.LocalizedFibreSets.starO_ι}, \emph{exists UJ umat eq of in UJ} \leanlink{Formalism.L2.SG.Pf.LocalizedFibreSets.exists_UJ_umat_eq_of_inUJ}, \emph{card herm Sub} \leanlink{Formalism.HermitianLattices.Ext.L3.reductionMapsOfForms.card_hermSub} with \emph{exists unit quot} \leanlink{Formalism.HermitianLattices.Ext.L3.reductionMapsOfForms.exists_unit_quot} yields the stated result.
\end{proof}

\begin{theorem}[li Ratio]\leanlink{Formalism.HermitianLattices.Ext.L3.liRatio}
\[
\mathrm{LIRatio}
\]
\end{theorem}
\begin{proof}
The argument operates on the structural objects LIRatio \leanlink{Formalism.HermitianLattices.Ext.L3.LIRatio}, vd \leanlink{Formalism.HermitianLattices.Ext.L3.vd}, Is Integral Lat \leanlink{Formalism.HermitianLatticesTheory.IsIntegralLat}, alpha Std \leanlink{Formalism.HermitianLattices.Ext.L3.alphaStd}. Combining \emph{exists loc Reps} \leanlink{Formalism.HermitianLattices.Ext.L3.exists_locReps}, \emph{rep Ratio eq zero} \leanlink{Formalism.HermitianLattices.Ext.L3.LIOrb.repRatio_eq_zero}, \emph{main term} \leanlink{Formalism.HermitianLattices.Ext.L3.LIOrb.main_term}, \emph{psi P eq sum} \leanlink{Formalism.HermitianLattices.Ext.L3.psiP_eq_sum}, \emph{rep Ratio eq} \leanlink{Formalism.HermitianLattices.Ext.L3.LIOrb.repRatio_eq} with \emph{eigenvalue eq2} \leanlink{Formalism.HermitianLattices.Ext.L3.LIOrb.eigenvalue_eq₂} yields the stated result.
\end{proof}

\begin{theorem}[alpha Disc of disc Onto]\leanlink{Formalism.L2.SG.Pf.G5.alphaDisc_of_discOnto}
\[
(hD : \mathrm{DiscOnto}) : \mathrm{AlphaDisc}
\]
\end{theorem}
\begin{proof}
The argument operates on the structural objects Ker Count \leanlink{Formalism.L2.SG.Pf.G5.KerCount}, Disc Onto \leanlink{Formalism.L2.SG.Pf.G5.DiscOnto}, Alpha Disc \leanlink{Formalism.L2.SG.Pf.Cnt.AlphaDisc}. Combining \emph{phi counting holds} \leanlink{Formalism.L2.SG.Pf.LocalizedFibreSets.phi_counting_holds}, \emph{ker Count main} \leanlink{Formalism.L2.SG.Pf.discriminantGroupCounts.kerCount_main} with \emph{alpha Disc of} \leanlink{Formalism.L2.SG.Pf.G5.alphaDisc_of} yields the stated result.
\end{proof}

\begin{theorem}[row class holds]\leanlink{PI.UnitaryGroupMassFormula.CntFormula.RowsMassW.row_class_holds}
\[
\mathrm{RowClass}
\]
\end{theorem}
\begin{proof}
The argument operates on the structural objects Row Class \leanlink{PI.UnitaryGroupMassFormula.CntFormula.RowsMassW.RowClass}. Combining \emph{row unit condition} \leanlink{PI.UnitaryGroupMassFormula.CntFormula.RowsMassW.row_unit_condition} with \emph{row Class of zunit} \leanlink{PI.UnitaryGroupMassFormula.CntFormula.RowsMassW.rowClass_of_zunit} yields the stated result.
\end{proof}

\section{Ground Theory: zel}

This chapter encompasses 116,103 formal declarations (1,632,709~LOC). Primary proof idioms detected: \emph{automorphic / spectral theory, Vandiver / cohomology, p-adic / local harmonic analysis}.
\begin{quote}\small
\textbf{Formal Interface ($\mathrm{Inputs} \vdash \mathrm{Outputs}$):}
\par\noindent\emph{Inputs:} algebra Map~\leanlink{Formalism.GradedWeylDivisibility.CoeffHom.algebraMap}, D7~\leanlink{Formalism.BasicNotions.D7}, F~\leanlink{Formalism.HC.HostRaw.F}, map mul~\leanlink{GlobalTheory.Serre.Lie.map_mul}, Spc~\leanlink{GlobalTheory.EtC.Spc}, n~\leanlink{Formalism.BasicNotions.Ranks.n}.
\par\noindent\emph{Outputs:} G~\leanlink{Formalism.ConjugateSelfDual.G}, Is Place Over~\leanlink{Formalism.ConjugateSelfDual.IsPlaceOver}, Qbar~\leanlink{Formalism.ConjugateSelfDual.Qbar}, Ine~\leanlink{ReciprocityTheory.Gal.Ine}, Dec~\leanlink{ReciprocityTheory.Gal.Dec}, Ideles~\leanlink{Neutral.B1.Ideles}.
\end{quote}

\subsection{Core Definitions and Structures}
\begin{definition}[Good K]\leanlink{UpperTheory.TP.W.GoodK}
\[
(\sigma : E \to _{a} E) (B : \pcBase F E) : \mathrm{Prop}
\]
\end{definition}

\begin{definition}[Same Galois Level]\leanlink{GlobalTheory.HeightCounting.Up.SameGaloisLevel}
\[
(hL : \mathrm{RecTExists}) (x : St P) : \mathrm{Type}1
\]
\end{definition}

\begin{definition}[Intermediate Level]\leanlink{GlobalTheory.HeightCounting.Up.IntermediateLevel}
\begin{quote}\small\raggedright$\displaystyle (hL : \mathrm{RecTExists}) \allowbreak (x : St P) : \mathrm{Type}1 \mathrm{extends} \mathrm{toLA} : \mathrm{SameGaloisLevel} hL x$\end{quote}
\end{definition}

\begin{definition}[Second Intermediate Level]\leanlink{GlobalTheory.HeightCounting.Up.SecondIntermediateLevel}
\begin{quote}\small\raggedright$\displaystyle (hL : \mathrm{RecTExists}) \allowbreak (x : St P) : \mathrm{Type}1 \mathrm{extends} \mathrm{toLI} : \mathrm{IntermediateLevel} hL x$\end{quote}
\end{definition}

\begin{definition}[First Intermediate Level]\leanlink{GlobalTheory.HeightCounting.Up.FirstIntermediateLevel}
\begin{quote}\small\raggedright$\displaystyle (hL : \mathrm{RecTExists}) \allowbreak (x : St P) : \mathrm{Type}1 \mathrm{extends} \mathrm{toLI} : \mathrm{IntermediateLevel} hL x$\end{quote}
\end{definition}

\begin{definition}[Height Counting Base Data]\leanlink{GlobalTheory.HeightCounting.HeightCountingBaseData}
\begin{quote}\small\raggedright$\displaystyle (hL : \mathrm{RecTExists}) \allowbreak (D : D7) \allowbreak (cf : \mathrm{CharFrame} D) \allowbreak (\mathrm{rep} : D.\pcWhat .FT \to \allowbreak NG.\mathrm{FinTyp} D.\pcHostA .E \mathrm{r}) : \mathrm{Type}1$\end{quote}
\end{definition}

\subsection{Main Results}
\begin{theorem}[zel]\leanlink{UpperTheory.Zero.LLC.Pr.zel}
\[
\mathrm{ZelevinskyPolynomialRing}
\]
\end{theorem}
\begin{proof}
The argument operates on the structural objects Zelevinsky Polynomial Ring \leanlink{GlobalTheory.Published.ZelevinskyPolynomialRing}. Combining \emph{segment Representation} \leanlink{PositivityTheory.RgGu.segmentRepresentation}, \emph{grothendieck Ring Polynomial holds} \leanlink{PositivityTheory.Groth.grothendieckRingPolynomial_holds}, \emph{zelevinsky Polynomial Ring of} \leanlink{PositivityTheory.Prw.zelevinskyPolynomialRing_of} with \emph{bz} \leanlink{UpperTheory.Zero.LLC.Pr.bz} yields the stated result.
\end{proof}

\begin{theorem}[sb1l alt alt]\leanlink{GlobalTheory.HeightCounting.Up.SeriesStep.sb1l_alt_alt}
\begin{quote}\small\raggedright$\displaystyle \mathrm{PublishedInputsThirdLayerAlt} \to \allowbreak \mathrm{SerN} \to \allowbreak \forall hL : \mathrm{RecTExists}, \allowbreak \forall \mathrm{cDep} : \mathbb{R} , \allowbreak \mathrm{SeriesExistList}' hL \mathrm{cDep}$\end{quote}
\end{theorem}
\begin{proof}
The argument operates on the structural objects Series Exist List' \leanlink{GlobalTheory.HeightCounting.Up.SeriesExistList'}, Published Inputs Third Layer Alt \leanlink{GlobalTheory.HeightCounting.Up.PublishedInputsThirdLayerAlt}, Pre2 \leanlink{GlobalTheory.Mem.R3.Rw.Pre2}, Col W \leanlink{GlobalTheory.Mem.Top.Bl.ColW}. Combining \emph{twisted state transport heq levels'} \leanlink{GlobalTheory.T.twisted_state_transport_heq_levels'}, \emph{list St o Tw} \leanlink{GlobalTheory.HeightCounting.Up.SeriesStep.listSt_oTw}, \emph{and} \leanlink{GlobalTheory.Mem.Top.Bl.StepW.and}, \emph{of all} \leanlink{GlobalTheory.Mem.Top.Bl.StepW.of_all} with \emph{exists matrix input identity frame of inputs} \leanlink{GlobalTheory.MtId.exists_matrix_input_identity_frame_of_inputs} yields the stated result.
\end{proof}

\begin{proposition}[exists reduced]\leanlink{GroundTheory.Automorphic.Reduction.exists_reduced}
\begin{quote}\small\raggedright$\displaystyle \exists \mathrm{t\_0} : \mathbb{R} , \allowbreak 0 < \mathrm{t\_0} \land \allowbreak \forall (n : \mathbb{N} ) \allowbreak (g : \mathrm{AdelicGL} K n), \allowbreak \exists \gamma : GL (\mathrm{Fin} n) K, \allowbreak \exists b \in \mathrm{adelicBorel}, \allowbreak \exists k \in \mathrm{maxCompact}, \allowbreak \mathrm{diagonal} K n \gamma * g = b * k \land \allowbreak \forall i j : \mathrm{Fin} n, \allowbreak (j : \mathbb{N} ) = i + 1 \to \allowbreak \mathrm{t\_0} * \mathrm{adeleAbs} ((b : \mathrm{Matrix} (\mathrm{Fin} n) \allowbreak (\mathrm{Fin} n) A ) j j) \le \mathrm{adeleAbs} ((b : \mathrm{Matrix} (\mathrm{Fin} n) \allowbreak (\mathrm{Fin} n) A ) i i)$\end{quote}
\end{proposition}
\begin{proof}
The argument operates on the structural objects max Compact \leanlink{GroundTheory.Automorphic.Reduction.maxCompact}, adelic Borel \leanlink{GroundTheory.Automorphic.Reduction.adelicBorel}, adele Abs \leanlink{GroundTheory.Automorphic.Reduction.adeleAbs}, diagonal \leanlink{GroundTheory.Automorphic.AdelicGLn.diagonal}. Combining \emph{exists reduced of} \leanlink{GroundTheory.Automorphic.Reduction.exists_reduced_of} with \emph{exists adelic Norm vec Mul le} \leanlink{GroundTheory.Automorphic.Reduction.exists_adelicNorm_vecMul_le} yields the stated result.
\end{proof}

\begin{proposition}[pre at of side b tau true lam ok]\leanlink{GlobalTheory.HeightCounting.SerCopy.pre_at_of_side_b_tau_true_lam_ok}
\begin{quote}\small\raggedright$\displaystyle (hA : \mathrm{GlobalArtinReciprocity}) \allowbreak (hK : K.\pcOver B.\mathrm{toPre}) \allowbreak (hS : \mathrm{Start} K B f T w) \allowbreak (x : St P) \allowbreak (l : \mathrm{IntermediateLevel} hL x) \allowbreak (hB : \exists Bs : \mathrm{SideB} (\mathrm{Dof} x) \allowbreak (\mathrm{cfOf} x) \mathrm{row} (\mathrm{lowAll} (\mathrm{Dof} x)) l.\mathrm{Pq} \mathrm{S} \mathrm{hL} \mathrm{eig} \mathrm{In}\theta , \allowbreak \mathrm{TauTrue} Bs.\mathrm{exp}) \allowbreak (\mathrm{hSb} : \forall i, \allowbreak \mathrm{Sb} \subseteq (\mathrm{Dof} x).\mathrm{ntw}.\mathrm{S0} i) \allowbreak (hZ : \exists Z : \mathrm{LamDat} (\mathrm{Dof} x), \allowbreak Z.\mathrm{Ok}) : \mathrm{PreAt} (P := P) \mathrm{z} \mathrm{row} l.\mathrm{Pq} \mathrm{hL} \mathrm{eig} \mathrm{In}\theta$\end{quote}
\end{proposition}
\begin{proof}
The argument operates on the structural objects Start \leanlink{GlobalTheory.Mem.Top.Bl.Start}, Pre2 \leanlink{GlobalTheory.Mem.R3.Rw.Pre2}, Col W \leanlink{GlobalTheory.Mem.Top.Bl.ColW}, Dof \leanlink{GlobalTheory.HeightCounting.Up.Dof}. Combining \emph{is Choice of start of over} \leanlink{GlobalTheory.T.isChoice_of_start_of_over} with \emph{exp} \leanlink{Formalism.L2.SG.Pf.BunEx.Unip.exp} yields the stated result.
\end{proof}

\begin{proposition}[is Choice of start of over]\leanlink{GlobalTheory.T.isChoice_of_start_of_over}
\begin{quote}\small\raggedright$\displaystyle \forall (hL : \mathrm{RecTExists}) , \allowbreak K.\pcOver B.\mathrm{toPre} \to \allowbreak \mathrm{Start} K B f T w \to \allowbreak \forall (x : \mathrm{St} P) \allowbreak (l : \mathrm{SameGaloisLevel} hL x), \allowbreak \mathrm{IsChoice} (D := \mathrm{Dof} x) \allowbreak (\mathcal{P}.\mathrm{range} l.\mathrm{Pq}) \land \allowbreak (\forall t, \allowbreak \forall Q \in \mathrm{Q} t, \allowbreak Q \in \mathcal{P}.\mathrm{range} l.\mathrm{Pq}) \land \allowbreak (\forall j, \allowbreak \mathrm{FWT}.\mathrm{VStd} (\mathrm{Dof} x) j) \land \allowbreak \mathrm{C} = \mathrm{cfOf} x \land \allowbreak \forall (k : \mathrm{Fin} 2) \allowbreak (z : \mathbb{C} ), \allowbreak z ^ ((\mathrm{Dof})) \;\cdots$\end{quote}
\end{proposition}
\begin{proof}
The argument operates on the structural objects Start \leanlink{GlobalTheory.Mem.Top.Bl.Start}, St \leanlink{GlobalTheory.Mem.R3.Rw.St}, Pre2 \leanlink{GlobalTheory.Mem.R3.Rw.Pre2}, Col W \leanlink{GlobalTheory.Mem.Top.Bl.ColW}. The claim follows directly from \emph{facts alt} \leanlink{GlobalTheory.HeightCounting.FirstFacts.facts_alt}.
\end{proof}

\begin{proposition}[pre at twisted copy of side b tau true]\leanlink{GlobalTheory.HeightCounting.SerCopy.pre_at_twisted_copy_of_side_b_tau_true}
\begin{quote}\small\raggedright$\displaystyle (hA : \mathrm{GlobalArtinReciprocity}) \allowbreak (hK : K.\pcOver B.\mathrm{toPre}) \allowbreak (hS : \mathrm{Start} K B f T w) \allowbreak (x : St P) \allowbreak (l : \mathrm{IntermediateLevel} hL x) \allowbreak (N : \mathrm{TwNew} (\mathrm{Dof} x)) \allowbreak (hB : \exists Bs : \mathrm{SideB} N.\mathrm{toD7} (N.\mathrm{cf} (\mathrm{cfOf} x)) \allowbreak (N.\mathrm{row} \mathrm{row}) \allowbreak (\mathrm{lowAll} (\mathrm{Dof} x)) l.\mathrm{Pq} (N.\mathrm{rat} \mathrm{S}) \mathrm{hL} \mathrm{eig} \mathrm{In}\theta , \allowbreak \mathrm{TauTrue} Bs.\mathrm{exp}) \allowbreak (\mathrm{hSb} : \forall i, \allowbreak \mathrm{Sb} \subseteq \mathrm{S0} i) \allowbreak (hZ : \exists Z : \mathrm{LamDat} N.\mathrm{toD7}, \allowbreak Z.\mathrm{Ok}) : \mathrm{PreAt} (P := P) \allowbreak (\mathrm{datTw} x N) \allowbreak ((by \mathrm{with\_unfolding\_all} \mathrm{exact} ((N.\mathrm{row})))) \;\cdots$\end{quote}
\end{proposition}
\begin{proof}
The argument operates on the structural objects Start \leanlink{GlobalTheory.Mem.Top.Bl.Start}, Pre2 \leanlink{GlobalTheory.Mem.R3.Rw.Pre2}, Col W \leanlink{GlobalTheory.Mem.Top.Bl.ColW}, Dof \leanlink{GlobalTheory.HeightCounting.Up.Dof}. Combining \emph{pre at of side b tau true lam ok} \leanlink{GlobalTheory.HeightCounting.SerCopy.pre_at_of_side_b_tau_true_lam_ok}, \emph{moved eq dat Tw} \leanlink{GlobalTheory.HeightCounting.Up.moved_eq_datTw} with \emph{twisted state transport heq levels'} \leanlink{GlobalTheory.T.twisted_state_transport_heq_levels'} yields the stated result.
\end{proof}

\begin{lemma}[read package twist heq alt]\leanlink{GlobalTheory.T.read_package_twist_heq_alt}
\begin{quote}\small\raggedright$\displaystyle \forall (hL : \mathrm{RecTExists}) \allowbreak (x : St P) \allowbreak (l : \mathrm{IntermediateLevel} hL x) \allowbreak (N : \mathrm{TwNew} (\mathrm{Dof} x)) \allowbreak (m : \mathrm{Moved} (\mathrm{oTw} x N)) \allowbreak (hm : m.\mathrm{z} = \mathrm{datTw} x N) \allowbreak (hc : m.\mathrm{Copy}) \allowbreak (hh : m.\mathrm{Coh}) \allowbreak (r : \mathrm{ReadPk} m (\mathrm{sideTw} x N)) \allowbreak (hp : r.\mathrm{Pin}), \allowbreak \mathrm{HEq} (\mathrm{repOf} r.\mathrm{Y}) \allowbreak (\mathrm{repOf} \mathrm{Y}) \to \allowbreak (\forall \mu h i y, \allowbreak \mathrm{HEq} (\mathrm{jMt} \mu h i y).1 (\mathrm{jMt} \mu h i y).1) \to \allowbreak (\forall i y, \allowbreak \mathrm{HEq} (\mathrm{eW} i y) \allowbreak (\mathrm{eW} i y)) \to \allowbreak \exists l' : \mathrm{IntermediateLevel} hL (\mathrm{stTw} x N m hc hh r hp), \allowbreak \mathrm{HEq} l'.\mathrm{row} (N.\mathrm{row}) \;\cdots$\end{quote}
\end{lemma}
\begin{proof}
The argument operates on the structural objects Moved \leanlink{GlobalTheory.Mem.R3.Rw.Moved}, Copy \leanlink{GlobalTheory.Mem.R3.Rw.Moved.Copy}, Coh \leanlink{GlobalTheory.Mem.R3.Rw.Moved.Coh}, St \leanlink{GlobalTheory.Mem.R3.Rw.St}. This result is established directly from standard library foundations within its local module (2~lines of formal proof).
\end{proof}

\begin{lemma}[twisted state transport heq levels']\leanlink{GlobalTheory.T.twisted_state_transport_heq_levels'}
\begin{quote}\small\raggedright$\displaystyle \forall (hL : \mathrm{RecTExists}) \allowbreak (x : St P) \allowbreak (l : \mathrm{IntermediateLevel} hL x) \allowbreak (N : \mathrm{TwNew} (\mathrm{Dof} x)), \allowbreak \exists (x' : St P) \allowbreak (l' : \mathrm{IntermediateLevel} hL x'), \allowbreak x'.o = \mathrm{oTw} x N \land \allowbreak \mathrm{HEq} x'.\mathrm{toLevCh} (\mathrm{levCh} (a := \mathrm{Z}) \allowbreak (\mathrm{oTw} x N).Z \mathrm{toLevCh}) \land \allowbreak \mathrm{HEq} l'.\mathrm{row} (N.\mathrm{row} \mathrm{row}) \land \allowbreak \mathrm{HEq} l'.Pq l.\mathrm{Pq} \land \allowbreak \mathrm{HEq} l'.\mathrm{rat}.\mathrm{S} (N.\mathrm{rat} \mathrm{S}) \land \allowbreak \mathrm{HEq} l'.\mathrm{Sb} \mathrm{Sb} \land \allowbreak \mathrm{HEq} l'.\mathrm{hL} \mathrm{hL} \land \allowbreak \mathrm{HEq} \;\cdots$\end{quote}
\end{lemma}
\begin{proof}
The argument operates on the structural objects St \leanlink{GlobalTheory.Mem.R3.Rw.St}, Pre2 \leanlink{GlobalTheory.Mem.R3.Rw.Pre2}, real3 \leanlink{GlobalTheory.Mem.R3.Tr.real3}, Col W \leanlink{GlobalTheory.Mem.Top.Bl.ColW}. The claim follows directly from \emph{restage copy real alt} \leanlink{GlobalTheory.HeightCounting.Up.restage_copy_real_alt}.
\end{proof}

\begin{lemma}[exists reduced of]\leanlink{GroundTheory.Automorphic.Reduction.exists_reduced_of}
\begin{quote}\small\raggedright$\displaystyle (hC : 0 < C) \allowbreak (hD : \forall (\alpha \delta : \mathrm{IdeleGroup} K) \allowbreak (u : A ), \allowbreak \exists x : \mathrm{Fin} 2 \to \allowbreak K, \allowbreak x \neq 0 \land \allowbreak \mathrm{adelicNorm} (\mathrm{ratVec} x _{v} * !! ) \le C * (\mathrm{ideleModule} K (\alpha * \delta ) : \mathbb{R} )) : \forall (n : \mathbb{N} ) \allowbreak (g : \mathrm{AdelicGL} K n), \allowbreak \exists \gamma : GL (\mathrm{Fin} n) K, \allowbreak \exists b \in \mathrm{adelicBorel}, \allowbreak \exists k \in \mathrm{maxCompact}, \allowbreak \mathrm{diagonal} K n \gamma * g = b * k \land \allowbreak \forall i j : \mathrm{Fin} n, \allowbreak (j : \mathbb{N} ) = i + 1 \to \allowbreak \mathrm{adeleAbs} ((b : \mathrm{Matrix} (\mathrm{Fin} n) \allowbreak (\mathrm{Fin} n) A ) j j) \le C ^ 2 * \mathrm{adeleAbs} ((b : \mathrm{Matrix} (\mathrm{Fin} n) \allowbreak (\mathrm{Fin} n) A ) i i)$\end{quote}
\end{lemma}
\begin{proof}
The argument operates on the structural objects max Compact \leanlink{GroundTheory.Automorphic.Reduction.maxCompact}, idele Module \leanlink{GroundTheory.Global.IdeleGroup.ideleModule}. Combining \emph{exists adelic Borel mul max Compact} \leanlink{GroundTheory.Automorphic.Reduction.exists_adelicBorel_mul_maxCompact}, \emph{exists height min} \leanlink{GroundTheory.Automorphic.Reduction.exists_height_min}, \emph{adelic Norm single add single} \leanlink{GroundTheory.Automorphic.Reduction.adelicNorm_single_add_single}, \emph{exists last Row eq} \leanlink{GroundTheory.Automorphic.Reduction.exists_lastRow_eq}, \emph{block One Hom mem max Compact} \leanlink{GroundTheory.Automorphic.Reduction.blockOneHom_mem_maxCompact}, \emph{is Unit apply self of mem adelic Borel} \leanlink{GroundTheory.Automorphic.Reduction.isUnit_apply_self_of_mem_adelicBorel}, \emph{adelic Norm single units} \leanlink{GroundTheory.Automorphic.Reduction.adelicNorm_single_units} with \emph{height diagonal mul} \leanlink{GroundTheory.Automorphic.Reduction.height_diagonal_mul} yields the stated result.
\end{proof}

\begin{lemma}[facts alt]\leanlink{GlobalTheory.HeightCounting.FirstFacts.facts_alt}
\begin{quote}\small\raggedright$\displaystyle \forall (hL : \mathrm{RecTExists}) , \allowbreak K.\pcOver B.\mathrm{toPre} \to \allowbreak \mathrm{Start} K B f T w \to \allowbreak \forall (x : \mathrm{St} P) \allowbreak (l : \mathrm{SameGaloisLevel} hL x), \allowbreak \mathrm{IsChoice} (D := \mathrm{Dof} x) \allowbreak (\mathcal{P}.\mathrm{range} l.\mathrm{Pq}) \land \allowbreak (\forall t, \allowbreak \forall Q \in \mathrm{Q} t, \allowbreak Q \in \mathcal{P}.\mathrm{range} l.\mathrm{Pq}) \land \allowbreak (\forall j, \allowbreak \mathrm{FWT}.\mathrm{VStd} (\mathrm{Dof} x) j) \land \allowbreak \mathrm{C} = \mathrm{cfOf} x \land \allowbreak \forall (k : \mathrm{Fin} 2) \allowbreak (z : \mathbb{C} ), \allowbreak z ^ ((\mathrm{Dof})) \;\cdots$\end{quote}
\end{lemma}
\begin{proof}
The argument operates on the structural objects Start \leanlink{GlobalTheory.Mem.Top.Bl.Start}, St \leanlink{GlobalTheory.Mem.R3.Rw.St}, Pre2 \leanlink{GlobalTheory.Mem.R3.Rw.Pre2}, Col W \leanlink{GlobalTheory.Mem.Top.Bl.ColW}. Combining \emph{pair of over} \leanlink{GlobalTheory.HeightCounting.FirstFacts.pair_of_over} with \emph{ok} \leanlink{Formalism.L2.Col.GSD.Iso.ok} yields the stated result.
\end{proof}

\begin{lemma}[exists int add conj eq one of base good k]\leanlink{GlobalTheory.exists_int_add_conj_eq_one_of_base_good_k}
\begin{quote}\small\raggedright$\displaystyle \forall (B : \pcBase F E), \allowbreak \mathrm{GoodK} \sigma B \to \allowbreak \exists \theta : \mathcal{O} E, \allowbreak (\theta : E) + \sigma \theta = 1$\end{quote}
\end{lemma}
\begin{proof}
The argument operates on the structural objects Good K \leanlink{UpperTheory.TP.W.GoodK}, Base \leanlink{GlobalTheory.IndefiniteRowSideFullHecke.WeilTensorModels.Base}, cm \leanlink{Formalism.HermitianLattices.FiltrationBoundaryCriterion.PfN.cm}, cm \leanlink{Formalism.FourierDecayAlt.S3.cm}. Combining \emph{exists integer add conj eq one of formally Unramified} \leanlink{GlobalTheory.exists_integer_add_conj_eq_one_of_formallyUnramified}, \emph{cm} \leanlink{GlobalTheory.Published.HasseNorm.cm}, \emph{unram} \leanlink{GlobalTheory.Mem.FieldsLaws.unram} with \emph{unram} \leanlink{GlobalTheory.InertOfConj.unram} yields the stated result.
\end{proof}

\begin{lemma}[li copy alt]\leanlink{GlobalTheory.HeightCounting.Up.LiTw.li_copy_alt}
\begin{quote}\small\raggedright$\displaystyle \forall (hL : \mathrm{RecTExists}) \allowbreak (x : St P) \allowbreak (l : \mathrm{IntermediateLevel} hL x) \allowbreak (N : \mathrm{TwNew} (\mathrm{Dof} x)) \allowbreak (m : \mathrm{Moved} (\mathrm{oTw} x N)) \allowbreak (hm : m.\mathrm{z} = \mathrm{datTw} x N) \allowbreak (hc : m.\mathrm{Copy}) \allowbreak (hh : m.\mathrm{Coh}) \allowbreak (r : \mathrm{ReadPk} m (\mathrm{sideTw} x N)) \allowbreak (hp : r.\mathrm{Pin}), \allowbreak \mathrm{HEq} (\mathrm{repOf} r.\mathrm{Y}) \allowbreak (\mathrm{repOf} \mathrm{Y}) \to \allowbreak (\forall \mu h i y, \allowbreak \mathrm{HEq} (\mathrm{jMt} \mu h i y).1 (\mathrm{jMt} \mu h i y).1) \to \allowbreak (\forall i y, \allowbreak \mathrm{HEq} (\mathrm{eW} i y) \allowbreak (\mathrm{eW} i y)) \to \allowbreak \exists l' : \mathrm{IntermediateLevel} hL (\mathrm{stTw} x N m hc hh r hp), \allowbreak \mathrm{HEq} l'.\mathrm{row} (N.\mathrm{row}) \;\cdots$\end{quote}
\end{lemma}
\begin{proof}
The argument operates on the structural objects Moved \leanlink{GlobalTheory.Mem.R3.Rw.Moved}, Copy \leanlink{GlobalTheory.Mem.R3.Rw.Moved.Copy}, Coh \leanlink{GlobalTheory.Mem.R3.Rw.Moved.Coh}, St \leanlink{GlobalTheory.Mem.R3.Rw.St}. This result is established directly from standard library foundations within its local module (12~lines of formal proof).
\end{proof}

\section{Development / Resolution Of Singularities: functorial Resolution Char Zero}

This chapter encompasses 7,590 formal declarations (116,064~LOC). Primary proof idioms detected: \emph{abelian varieties \& heights, automorphic / spectral theory, determinant / Stepanov}.
\begin{quote}\small
\textbf{Formal Interface ($\mathrm{Inputs} \vdash \mathrm{Outputs}$):}
\par\noindent\emph{Inputs:} saturation By Powers~\leanlink{UpperTheory.Reformulated.B.SaturationByPowers.saturationByPowers}.
\par\noindent\emph{Outputs:} Blowup Sequence~\leanlink{FunctorialResolution.BlowupSequence}, Etale Chart~\leanlink{FunctorialResolution.EmbeddedResolution.EtaleChart}, coord~\leanlink{FunctorialResolution.EmbeddedResolution.EtaleChart.coord}, Has Simple Normal Crossings~\leanlink{FunctorialResolution.HasSimpleNormalCrossings}, Is Principalizing Sequence~\leanlink{FunctorialResolution.IsPrincipalizingSequence}, struct Map~\leanlink{FunctorialResolution.BlowupSequence.structMap}.
\end{quote}

\subsection{Core Definitions and Structures}
\begin{definition}[Blowup Chart]\leanlink{FunctorialResolution.EmbeddedResolution.BlowupChart}
\begin{quote}\small\raggedright$\displaystyle (\pi : \mathrm{X\_1} \longrightarrow \allowbreak X) \allowbreak (c : \mathrm{EtaleChart} p U m) \allowbreak (S : \mathcal{F} (\mathrm{Fin} m)) \allowbreak (s : \mathrm{Fin} m)$\end{quote}
\end{definition}

\begin{definition}[coord Relation]\leanlink{FunctorialResolution.EmbeddedResolution.BlowupChart.coordRelation}
\begin{quote}\small\raggedright$\displaystyle (_B : \mathrm{BlowupChart} \pi c S s) \allowbreak (i : \mathrm{Fin} m) : \mathrm{MvPolynomial} (\mathrm{ULift}. (\mathrm{Fin} m)) \Gamma (X, \allowbreak U)$\end{quote}
\end{definition}

\begin{definition}[Etale Chart]\leanlink{FunctorialResolution.EmbeddedResolution.EtaleChart}
\begin{quote}\small\raggedright$\displaystyle (p : X \longrightarrow \allowbreak \mathrm{Spec} (\mathrm{of} k)) \allowbreak (U : X.\mathrm{Opens}) \allowbreak (m : \mathbb{N} ) : \mathrm{Type}u$\end{quote}
\end{definition}

\begin{definition}[open Sections To Coord Relation Quotient]\leanlink{FunctorialResolution.EmbeddedResolution.BlowupChart.openSectionsToCoordRelationQuotient}
\begin{quote}\small\raggedright$\displaystyle (B : \mathrm{BlowupChart} \pi c S s) : \Gamma (U.\mathrm{toScheme}, \allowbreak \top ) \longrightarrow \allowbreak \mathrm{of} (\mathrm{MvPolynomial} (\mathrm{ULift}. (\mathrm{Fin} m)) \Gamma (X, \allowbreak U) / \mathrm{span} (\mathcal{P}.\mathrm{range} B.\mathrm{coordRelation}))$\end{quote}
\end{definition}

\begin{definition}[base Change Emb]\leanlink{FunctorialResolution.EmbeddedResolution.EmbeddingDatum.baseChangeEmb}
\begin{quote}\small\raggedright$\displaystyle (D : \mathrm{EmbeddingDatum} k Y p) \allowbreak (K : \mathrm{Type}u) : \mathrm{baseChange} p K \longrightarrow \allowbreak \mathrm{baseChange} D.\mathrm{ambientStruct} K$\end{quote}
\end{definition}

\begin{definition}[comparison]\leanlink{FunctorialResolution.EmbeddedResolution.LocalResolutionData.comparison}
\begin{quote}\small\raggedright$\displaystyle (d : \mathrm{LocalResolutionData} k f s) \allowbreak (i : \mathrm{I\_0} ) : \mathrm{BlowupSequencePullbackComparison} (\mathrm{f} i) \allowbreak (d.\mathrm{seq} i) s$\end{quote}
\end{definition}

\subsection{Main Results}
\begin{theorem}[functorial Resolution Char Zero]\leanlink{FunctorialResolution.functorialResolutionCharZero}
\[
\mathrm{FunctorialResolutionCharZero}.
\]
\end{theorem}
\begin{proof}
The argument operates on the structural objects tower \leanlink{FunctorialResolution.EmbeddedResolution.BlowupSequenceData.tower}, Functorial Resolution Char Zero \leanlink{FunctorialResolution.FunctorialResolutionCharZero}, tower \leanlink{FunctorialResolution.EmbeddedResolution.AffineSpacePrincipalizationData.tower}. The claim follows directly from \emph{strong Resolution Of Singularities Char Zero} \leanlink{FunctorialResolution.EmbeddedResolution.strongResolutionOfSingularitiesCharZero}.
\end{proof}

\begin{theorem}[canonical Algorithm Rules is Correct]\leanlink{FunctorialResolution.EmbeddedResolution.canonicalAlgorithmRules_isCorrect}
\begin{quote}\small\raggedright$\displaystyle (\mathrm{canonicalAlgorithmRules} \mathrm{strictTransformIsBlowupOfClosedSubscheme}. ).\mathrm{ReducesInductivelyOnDimension}$\end{quote}
\end{theorem}
\begin{proof}
The argument operates on the structural objects Reduces Inductively On Dimension \leanlink{FunctorialResolution.EmbeddedResolution.PrincipalizationRules.ReducesInductivelyOnDimension}, canonical Algorithm Rules \leanlink{FunctorialResolution.EmbeddedResolution.canonicalAlgorithmRules}. Combining \emph{is Correct canonical Algorithm Rules} \leanlink{FunctorialResolution.EmbeddedResolution.isCorrect_canonicalAlgorithmRules} with \emph{strict Transform Is Blowup Of Closed Subscheme} \leanlink{FunctorialResolution.EmbeddedResolution.strictTransformIsBlowupOfClosedSubscheme} yields the stated result.
\end{proof}

\begin{proposition}[strong Resolution Of Singularities Char Zero]\leanlink{FunctorialResolution.EmbeddedResolution.strongResolutionOfSingularitiesCharZero}
\[
\mathrm{StrongResolutionOfSingularitiesCharZero}.
\]
\end{proposition}
\begin{proof}
The argument operates on the structural objects Strong Resolution Of Singularities Char Zero \leanlink{FunctorialResolution.EmbeddedResolution.StrongResolutionOfSingularitiesCharZero}. Combining \emph{transitions Are Blowups For Embeddings principalization With Radical Centers} \leanlink{FunctorialResolution.EmbeddedResolution.transitionsAreBlowupsForEmbeddings_principalizationWithRadicalCenters}, \emph{embedded Principalization Yields Blowups principalization With Radical Centers} \leanlink{FunctorialResolution.EmbeddedResolution.embeddedPrincipalizationYieldsBlowups_principalizationWithRadicalCenters}, \emph{is Strong Resolution canonical Affine Resolution Rule} \leanlink{FunctorialResolution.EmbeddedResolution.isStrongResolution_canonicalAffineResolutionRule}, \emph{affine Embeds Into Affine Space} \leanlink{FunctorialResolution.EmbeddedResolution.affineEmbedsIntoAffineSpace}, \emph{closed Embeddings Related By Automorphism} \leanlink{FunctorialResolution.EmbeddedResolution.closedEmbeddingsRelatedByAutomorphism}, \emph{strong Resolution Of Singularities Char Zero of principalization Algorithm} \leanlink{FunctorialResolution.EmbeddedResolution.strongResolutionOfSingularitiesCharZero_of_principalizationAlgorithm}, \emph{is Functorial Principalization principalization With Radical Centers} \leanlink{FunctorialResolution.EmbeddedResolution.isFunctorialPrincipalization_principalizationWithRadicalCenters} with \emph{affine Space Is Smooth} \leanlink{FunctorialResolution.EmbeddedResolution.affineSpaceIsSmooth} yields the stated result.
\end{proof}

\begin{proposition}[is Correct canonical Algorithm Rules]\leanlink{FunctorialResolution.EmbeddedResolution.isCorrect_canonicalAlgorithmRules}
\begin{quote}\small\raggedright$\displaystyle (\mathrm{hres} : \mathrm{StrictTransformIsBlowupOfClosedSubscheme} ) : (\mathrm{canonicalAlgorithmRules} \mathrm{hres}).\mathrm{ReducesInductivelyOnDimension}$\end{quote}
\end{proposition}
\begin{proof}
The argument operates on the structural objects Reduces Inductively On Dimension \leanlink{FunctorialResolution.EmbeddedResolution.PrincipalizationRules.ReducesInductivelyOnDimension}, Strict Transform Is Blowup Of Closed Subscheme \leanlink{FunctorialResolution.Blowup.StrictTransformIsBlowupOfClosedSubscheme}, canonical Algorithm Rules \leanlink{FunctorialResolution.EmbeddedResolution.canonicalAlgorithmRules}. Combining \emph{is Correct canonical Algorithm Rules of hypotheses} \leanlink{FunctorialResolution.EmbeddedResolution.isCorrect_canonicalAlgorithmRules_of_hypotheses}, \emph{principalization Geometry Hypotheses} \leanlink{FunctorialResolution.EmbeddedResolution.principalizationGeometryHypotheses} with \emph{divisor Reduction Hypotheses canonical Algorithm Rules} \leanlink{FunctorialResolution.EmbeddedResolution.divisorReductionHypotheses_canonicalAlgorithmRules} yields the stated result.
\end{proof}

\begin{proposition}[functorial Principalization Char Zero]\leanlink{FunctorialResolution.functorialPrincipalizationCharZero}
\[
\mathrm{FunctorialPrincipalizationCharZero}.
\]
\end{proposition}
\begin{proof}
The argument operates on the structural objects Functorial Principalization Char Zero \leanlink{FunctorialResolution.FunctorialPrincipalizationCharZero}, principalization With Radical Centers \leanlink{FunctorialResolution.EmbeddedResolution.principalizationWithRadicalCenters}. Combining \emph{is Functorial Principalization principalization With Radical Centers} \leanlink{FunctorialResolution.EmbeddedResolution.isFunctorialPrincipalization_principalizationWithRadicalCenters} with \emph{functorial Principalization Char Zero of algorithm} \leanlink{FunctorialResolution.functorialPrincipalizationCharZero_of_algorithm} yields the stated result.
\end{proof}

\begin{proposition}[transitions Are Blowups For Embeddings principalization With Radical Centers]\leanlink{FunctorialResolution.EmbeddedResolution.transitionsAreBlowupsForEmbeddings_principalizationWithRadicalCenters}
\[
\mathrm{TransitionsAreBlowupsForEmbeddings} \mathrm{principalizationWithRadicalCenters}.
\]
\end{proposition}
\begin{proof}
The argument operates on the structural objects Transitions Are Blowups For Embeddings \leanlink{FunctorialResolution.EmbeddedResolution.TransitionsAreBlowupsForEmbeddings}, principalization With Radical Centers \leanlink{FunctorialResolution.EmbeddedResolution.principalizationWithRadicalCenters}. Combining \emph{detached Locus Is Open For Embeddings principalization With Radical Centers} \leanlink{FunctorialResolution.EmbeddedResolution.detachedLocusIsOpenForEmbeddings_principalizationWithRadicalCenters} with \emph{transitions Are Blowups For Embeddings of detached Locus Is Open} \leanlink{FunctorialResolution.EmbeddedResolution.transitionsAreBlowupsForEmbeddings_of_detachedLocusIsOpen} yields the stated result.
\end{proof}

\begin{lemma}[is Correct canonical Algorithm Rules of hypotheses]\leanlink{FunctorialResolution.EmbeddedResolution.isCorrect_canonicalAlgorithmRules_of_hypotheses}
\begin{quote}\small\raggedright$\displaystyle (\mathrm{hres} : \mathrm{StrictTransformIsBlowupOfClosedSubscheme} ) \allowbreak (\mathrm{hdiv} : \mathrm{DivisorReductionHypotheses} \mathrm{hres}) \allowbreak (hG : \mathrm{PrincipalizationGeometryHypotheses}. ) : (\mathrm{canonicalAlgorithmRules} \mathrm{hres}).\mathrm{ReducesInductivelyOnDimension}$\end{quote}
\end{lemma}
\begin{proof}
The argument operates on the structural objects Principalization Geometry Hypotheses \leanlink{FunctorialResolution.EmbeddedResolution.PrincipalizationGeometryHypotheses}, Reduces Inductively On Dimension \leanlink{FunctorialResolution.EmbeddedResolution.PrincipalizationRules.ReducesInductivelyOnDimension}, Divisor Reduction Hypotheses \leanlink{FunctorialResolution.EmbeddedResolution.DivisorReduction.DivisorReductionHypotheses}, Strict Transform Is Blowup Of Closed Subscheme \leanlink{FunctorialResolution.Blowup.StrictTransformIsBlowupOfClosedSubscheme}. Combining \emph{is Boundary Position Invariant divisor Family canonical Algorithm Rules} \leanlink{FunctorialResolution.EmbeddedResolution.isBoundaryPositionInvariant_divisorFamily_canonicalAlgorithmRules}, \emph{is Order Reduction order Family of strict Transform Is Blowup} \leanlink{FunctorialResolution.EmbeddedResolution.OrderReductionProperties.isOrderReduction_orderFamily_of_strictTransformIsBlowup}, \emph{is Max Order Reduction canonical Algorithm Rules} \leanlink{FunctorialResolution.EmbeddedResolution.MaxOrderReduction.isMaxOrderReduction_canonicalAlgorithmRules} with \emph{is Boundary Divisor Reduction divisor Family canonical Algorithm Rules} \leanlink{FunctorialResolution.EmbeddedResolution.isBoundaryDivisorReduction_divisorFamily_canonicalAlgorithmRules} yields the stated result.
\end{proof}

\begin{lemma}[order Reduction Hypotheses]\leanlink{FunctorialResolution.EmbeddedResolution.OrderReductionProperties.orderReductionHypotheses}
\[
\mathrm{OrderReductionHypotheses}
\]
\end{lemma}
\begin{proof}
The argument operates on the structural objects Order Reduction Hypotheses \leanlink{FunctorialResolution.EmbeddedResolution.OrderReduction.OrderReductionHypotheses}. Combining \emph{max Contact Hypersurfaces Etale Equivalent} \leanlink{FunctorialResolution.EmbeddedResolution.OrderReductionProperties.maxContactHypersurfacesEtaleEquivalent}, \emph{blowup Lift Eq On Controlled Transform} \leanlink{FunctorialResolution.EmbeddedResolution.OrderReductionCorrectness.blowupLiftEqOnControlledTransform}, \emph{exists Max Contact Controlled Transforms} \leanlink{FunctorialResolution.EmbeddedResolution.OrderReductionCorrectness.existsMaxContactControlledTransforms}, \emph{krull Dim Succ le of etale} \leanlink{FunctorialResolution.EmbeddedResolution.krullDimSucc_le_of_etale}, \emph{projection To Base mem order At Least Locus of mem center} \leanlink{FunctorialResolution.EmbeddedResolution.OrderReductionCorrectness.projectionToBase_mem_orderAtLeastLocus_of_mem_center}, \emph{is Admissible pullback Along of smooth} \leanlink{FunctorialResolution.EmbeddedResolution.MarkedIdealWithBoundary.isAdmissible_pullbackAlong_of_smooth}, \emph{is Balanced comap of smooth} \leanlink{FunctorialResolution.EmbeddedResolution.OrderReductionCorrectness.isBalanced_comap_of_smooth} with \emph{is Locally Noetherian scheme At} \leanlink{FunctorialResolution.EmbeddedResolution.OrderReductionCorrectness.isLocallyNoetherian_schemeAt} yields the stated result.
\end{proof}

\begin{lemma}[embedded Principalization Yields Blowups principalization With Radical Centers]\leanlink{FunctorialResolution.EmbeddedResolution.embeddedPrincipalizationYieldsBlowups_principalizationWithRadicalCenters}
\begin{quote}\small\raggedright$\displaystyle \mathrm{EmbeddedPrincipalizationYieldsBlowups} \mathrm{principalizationWithRadicalCenters}. \mathrm{affineEmbedsIntoAffineSpace} \mathrm{nilIsStrictNormalCrossingsOnAffineSpace} \mathrm{affineSpaceIsSmooth}$\end{quote}
\end{lemma}
\begin{proof}
The argument operates on the structural objects Embedded Principalization Yields Blowups \leanlink{FunctorialResolution.EmbeddedResolution.EmbeddedPrincipalizationYieldsBlowups}, principalization With Radical Centers \leanlink{FunctorialResolution.EmbeddedResolution.principalizationWithRadicalCenters}. Combining \emph{detached Locus Is Open For Embeddings principalization With Radical Centers} \leanlink{FunctorialResolution.EmbeddedResolution.detachedLocusIsOpenForEmbeddings_principalizationWithRadicalCenters}, \emph{affine Embeds Into Affine Space} \leanlink{FunctorialResolution.EmbeddedResolution.affineEmbedsIntoAffineSpace}, \emph{affine Space Is Smooth} \leanlink{FunctorialResolution.EmbeddedResolution.affineSpaceIsSmooth}, \emph{embedded Principalization Yields Blowups of discarded Locus Is Open} \leanlink{FunctorialResolution.EmbeddedResolution.embeddedPrincipalizationYieldsBlowups_of_discardedLocusIsOpen}, \emph{nil Is Strict Normal Crossings On Affine Space} \leanlink{FunctorialResolution.EmbeddedResolution.nilIsStrictNormalCrossingsOnAffineSpace}, \emph{strict Transform Is Blowup Of Closed Subscheme} \leanlink{FunctorialResolution.EmbeddedResolution.strictTransformIsBlowupOfClosedSubscheme}, \emph{embedded Principalization Discarded Locus Is Open of detached Locus Is Open} \leanlink{FunctorialResolution.EmbeddedResolution.embeddedPrincipalizationDiscardedLocusIsOpen_of_detachedLocusIsOpen} with \emph{embedded Principalization Ambient Is Noetherian} \leanlink{FunctorialResolution.EmbeddedResolution.embeddedPrincipalizationAmbientIsNoetherian} yields the stated result.
\end{proof}

\begin{lemma}[smooth Pullback Has Graph Chart]\leanlink{FunctorialResolution.EmbeddedResolution.smoothPullbackHasGraphChart}
\[
\mathrm{SmoothPullbackHasGraphChart}.
\]
\end{lemma}
\begin{proof}
The argument operates on the structural objects Smooth Pullback Has Graph Chart \leanlink{FunctorialResolution.EmbeddedResolution.SmoothPullbackHasGraphChart}, spec To Affine Space Extend Coords \leanlink{FunctorialResolution.EmbeddedResolution.specToAffineSpaceExtendCoords}. Combining \emph{smooth graph Morphism To Affine Space basic Open} \leanlink{FunctorialResolution.EmbeddedResolution.smooth_graphMorphismToAffineSpace_basicOpen}, \emph{exists affine Open is Pullback spec Map} \leanlink{FunctorialResolution.EmbeddedResolution.PresentedSmoothChart.exists_affineOpen_isPullback_specMap}, \emph{ker spec To Affine Space Extend Coords app Top} \leanlink{FunctorialResolution.EmbeddedResolution.ker_specToAffineSpaceExtendCoords_appTop}, \emph{app Top spec To Affine Space Extend Coords jacobian Det Section} \leanlink{FunctorialResolution.EmbeddedResolution.appTop_specToAffineSpaceExtendCoords_jacobianDetSection}, \emph{is Closed Immersion spec To Affine Space Extend Coords} \leanlink{FunctorialResolution.EmbeddedResolution.isClosedImmersion_specToAffineSpaceExtendCoords}, \emph{spec To Affine Space Extend Coords comp proj Left} \leanlink{FunctorialResolution.EmbeddedResolution.specToAffineSpaceExtendCoords_comp_projLeft}, \emph{smooth spec Map comp iso Spec inv} \leanlink{FunctorialResolution.EmbeddedResolution.PresentedSmoothChart.smooth_specMap_comp_isoSpec_inv} with \emph{preimage basic Open spec To Affine Space Extend Coords} \leanlink{FunctorialResolution.EmbeddedResolution.preimage_basicOpen_specToAffineSpaceExtendCoords} yields the stated result.
\end{proof}

\begin{lemma}[is Strong Resolution canonical Affine Resolution Rule]\leanlink{FunctorialResolution.EmbeddedResolution.isStrongResolution_canonicalAffineResolutionRule}
\begin{quote}\small\raggedright$\displaystyle (k : \mathrm{Type}u) \allowbreak (Y : \mathrm{Scheme}. ) \allowbreak (p : Y \longrightarrow \allowbreak \mathrm{Spec} (\mathrm{of} k)) : \mathrm{IsStrongResolution} k Y p (\mathrm{canonicalAffineResolutionRule} k Y p)$\end{quote}
\end{lemma}
\begin{proof}
Combining \emph{embedded Principalization Yields Blowups principalization With Radical Centers} \leanlink{FunctorialResolution.EmbeddedResolution.embeddedPrincipalizationYieldsBlowups_principalizationWithRadicalCenters}, \emph{affine Centers Lie Over Singular Locus principalization With Radical Centers} \leanlink{FunctorialResolution.EmbeddedResolution.affineCentersLieOverSingularLocus_principalizationWithRadicalCenters}, \emph{detached Locus Is Open For Embeddings principalization With Radical Centers} \leanlink{FunctorialResolution.EmbeddedResolution.detachedLocusIsOpenForEmbeddings_principalizationWithRadicalCenters}, \emph{affine Embeds Into Affine Space} \leanlink{FunctorialResolution.EmbeddedResolution.affineEmbedsIntoAffineSpace}, \emph{strict Normal Crossings Restrict To Clopen} \leanlink{FunctorialResolution.EmbeddedResolution.strictNormalCrossingsRestrictToClopen}, \emph{blowup Is Proper} \leanlink{FunctorialResolution.Blowup.blowupIsProper}, \emph{is Functorial Principalization principalization With Radical Centers} \leanlink{FunctorialResolution.EmbeddedResolution.isFunctorialPrincipalization_principalizationWithRadicalCenters} with \emph{affine Space Is Smooth} \leanlink{FunctorialResolution.EmbeddedResolution.affineSpaceIsSmooth} yields the stated result.
\end{proof}

\begin{lemma}[affine Centers Lie Over Singular Locus principalization With Radical Centers]\leanlink{FunctorialResolution.EmbeddedResolution.affineCentersLieOverSingularLocus_principalizationWithRadicalCenters}
\begin{quote}\small\raggedright$\displaystyle \mathrm{AffineCentersLieOverSingularLocus}. \mathrm{principalizationWithRadicalCenters} \mathrm{affineEmbedsIntoAffineSpace} \mathrm{nilIsStrictNormalCrossingsOnAffineSpace} \mathrm{affineSpaceIsSmooth} \mathrm{embeddedPrincipalizationYieldsBlowups\_principalizationWithRadicalCenters}$\end{quote}
\end{lemma}
\begin{proof}
Combining \emph{canonical Algorithm Rules is Correct} \leanlink{FunctorialResolution.EmbeddedResolution.canonicalAlgorithmRules_isCorrect}, \emph{transitions Are Blowups For Embeddings principalization With Radical Centers} \leanlink{FunctorialResolution.EmbeddedResolution.transitionsAreBlowupsForEmbeddings_principalizationWithRadicalCenters}, \emph{embedded Principalization Yields Blowups principalization With Radical Centers} \leanlink{FunctorialResolution.EmbeddedResolution.embeddedPrincipalizationYieldsBlowups_principalizationWithRadicalCenters}, \emph{affine Embeds Into Affine Space} \leanlink{FunctorialResolution.EmbeddedResolution.affineEmbedsIntoAffineSpace}, \emph{to Embedded Blowup Sequence centers Lie Over smooth Locus Complement} \leanlink{FunctorialResolution.EmbeddedResolution.toEmbeddedBlowupSequence_centersLieOver_smoothLocusComplement}, \emph{commutes With Closed Immersion of is Correct} \leanlink{FunctorialResolution.EmbeddedResolution.commutesWithClosedImmersion_of_isCorrect}, \emph{smooth Subscheme Ideal Is Radical} \leanlink{FunctorialResolution.EmbeddedResolution.smoothSubschemeIdealIsRadical} with \emph{affine Space Is Smooth} \leanlink{FunctorialResolution.EmbeddedResolution.affineSpaceIsSmooth} yields the stated result.
\end{proof}

\section{Reciprocity Theory / Geometry Of Numbers: exists chi P succ sub le}

This chapter encompasses 227 formal declarations (2,762~LOC). Primary proof idioms detected: \emph{Bost slope inequality}.
\begin{quote}\small
\textbf{Formal Interface ($\mathrm{Inputs} \vdash \mathrm{Outputs}$):}
\par\noindent\emph{Inputs:} measure Preserving swap Pi~\leanlink{UpperTheory.Reformulated.B.Qe3b.measurePreserving_swapPi}.
\par\noindent\emph{Outputs:} Herm Fam~\leanlink{ReciprocityTheory.GeometryOfNumbersTheory.HermFam}, deg~\leanlink{ReciprocityTheory.GeometryOfNumbersTheory.deg}, gram~\leanlink{ReciprocityTheory.GeometryOfNumbersTheory.HermFam.gram}, lat Rank~\leanlink{ReciprocityTheory.GeometryOfNumbersTheory.latRank}, nrm~\leanlink{ReciprocityTheory.GeometryOfNumbersTheory.HermFam.nrm}, Is Frame~\leanlink{ReciprocityTheory.GeometryOfNumbersTheory.IsFrame}.
\end{quote}

\subsection{Core Definitions and Structures}
\begin{definition}[Herm Fam]\leanlink{ReciprocityTheory.GeometryOfNumbersTheory.HermFam}
\[
(K : \mathrm{Type}) (V : \mathrm{Type})
\]
\end{definition}

\begin{definition}[Is Conj Inv]\leanlink{ReciprocityTheory.GeometryOfNumbersTheory.HermFam.IsConjInv}
\[
(h : \mathrm{HermFam} K V) : \mathrm{Prop}
\]
\end{definition}

\begin{definition}[ell T]\leanlink{ReciprocityTheory.GeometryOfNumbersTheory.ellT}
\begin{quote}\small\raggedright$\displaystyle (h : \mathrm{HermFam} K V) \allowbreak (v : \mathrm{Fin} r \to \allowbreak V) \allowbreak (R : \mathbb{R} ) : \mathcal{P} (( \to \allowbreak \mathrm{Fin} r \to \allowbreak \mathbb{R} ) \times ( \to \allowbreak \mathrm{Fin} r \to \allowbreak \mathbb{C} ))$\end{quote}
\end{definition}

\begin{definition}[ell S]\leanlink{ReciprocityTheory.GeometryOfNumbersTheory.ellS}
\begin{quote}\small\raggedright$\displaystyle (h : \mathrm{HermFam} K V) \allowbreak (v : \mathrm{Fin} r \to \allowbreak V) \allowbreak (R : \mathbb{R} ) : \mathcal{P} (\mathrm{Fin} r \to \allowbreak \mathrm{mixedSpace} K)$\end{quote}
\end{definition}

\begin{definition}[crd]\leanlink{ReciprocityTheory.GeometryOfNumbersTheory.crd}
\[
(\mathrm{hli} : \mathrm{LinearIndependent} K v) : V \to _{l} (\mathrm{Fin} r \to K)
\]
\end{definition}

\begin{definition}[degree Of Family]\leanlink{ReciprocityTheory.GeometryOfNumbersTheory.degreeOfFamily}
\begin{quote}\small\raggedright$\displaystyle (h : \mathrm{HermFam} K V) \allowbreak (E : \mathrm{Submodule} (\mathcal{O} K) V) \allowbreak (v : \mathrm{Fin} r \to \allowbreak V) : \mathbb{R}$\end{quote}
\end{definition}

\subsection{Main Results}
\begin{theorem}[exists chi P succ sub le]\leanlink{ReciprocityTheory.GeometryOfNumbersTheory.exists_chiP_succ_sub_le}
\begin{quote}\small\raggedright$\displaystyle (D : \mathbb{N} ) : \exists C : \mathbb{R} , \allowbreak \forall n : \mathbb{N} , \allowbreak 1 \le n \to \allowbreak \mathrm{chiP} D (n + 1) - \mathrm{chiP} D n \le C * (n : \mathbb{R} ) ^ D * \pcReal .\mathrm{log} (n + 1)$\end{quote}
\end{theorem}
\begin{proof}
The argument operates on the structural objects chi P \leanlink{ReciprocityTheory.GeometryOfNumbersTheory.chiP}, mono Norm Sq \leanlink{ReciprocityTheory.GeometryOfNumbersTheory.monoNormSq}. Combining \emph{chi PAux neg log mono Norm Sq le} \leanlink{ReciprocityTheory.GeometryOfNumbersTheory.chiPAux_neg_log_monoNormSq_le}, \emph{chi P succ sub eq} \leanlink{ReciprocityTheory.GeometryOfNumbersTheory.chiP_succ_sub_eq}, \emph{chi PAux sum filter zero le} \leanlink{ReciprocityTheory.GeometryOfNumbersTheory.chiPAux_sum_filter_zero_le}, \emph{chi PAux card adt le} \leanlink{ReciprocityTheory.GeometryOfNumbersTheory.chiPAux_card_adt_le}, \emph{chi PAux neg log step mem} \leanlink{ReciprocityTheory.GeometryOfNumbersTheory.chiPAux_neg_log_step_mem} with \emph{log add add one le log add log} \leanlink{ReciprocityTheory.GeometryOfNumbersTheory.log_add_add_one_le_log_add_log} yields the stated result.
\end{proof}

\begin{theorem}[integral form sq]\leanlink{ReciprocityTheory.GeometryOfNumbersTheory.integral_form_sq}
\begin{quote}\small\raggedright$\displaystyle (D n : \mathbb{N} ) \allowbreak (c : (\mathrm{Fin} (D + 1) \to \allowbreak \mathbb{N} ) \to \allowbreak \mathbb{C} ) : \int z, \allowbreak \| \sum a \in \mathcal{F}.\mathbb{N}.\mathrm{antidiagonalTuple} (D + 1) n, \allowbreak c a * \mathrm{monoC} a z\| ^ 2 \partial (\mathrm{sphProb} (D + 1)) = \sum a \in \mathcal{F}.\mathbb{N}.\mathrm{antidiagonalTuple} (D + 1) n, \allowbreak \| c a\| ^ 2 * \mathrm{monoNormSq} a$\end{quote}
\end{theorem}
\begin{proof}
The argument operates on the structural objects sph Prob \leanlink{ReciprocityTheory.GeometryOfNumbersTheory.sphProb}, mono Norm Sq \leanlink{ReciprocityTheory.GeometryOfNumbersTheory.monoNormSq}, mono C \leanlink{ReciprocityTheory.GeometryOfNumbersTheory.monoC}. The claim follows directly from \emph{sup St gram} \leanlink{ReciprocityTheory.GeometryOfNumbersTheory.supSt_gram}.
\end{proof}

\begin{proposition}[arakelov Degree eq degree Of Family]\leanlink{ReciprocityTheory.GeometryOfNumbersTheory.arakelovDegree_eq_degreeOfFamily}
\begin{quote}\small\raggedright$\displaystyle (h : \mathrm{HermFam} K V) \allowbreak (hE : E.\mathrm{FG}) \allowbreak (hv : \mathrm{IsFrame} K E v) : \mathrm{deg} h E = \mathrm{degreeOfFamily} h E v$\end{quote}
\end{proposition}
\begin{proof}
The argument operates on the structural objects Herm Fam \leanlink{ReciprocityTheory.GeometryOfNumbersTheory.HermFam}, degree Of Family \leanlink{ReciprocityTheory.GeometryOfNumbersTheory.degreeOfFamily}, deg \leanlink{ReciprocityTheory.GeometryOfNumbersTheory.deg}, Is Frame \leanlink{ReciprocityTheory.GeometryOfNumbersTheory.IsFrame}. The claim follows directly from \emph{degree Of Family eq degree Of Family} \leanlink{ReciprocityTheory.GeometryOfNumbersTheory.degreeOfFamily_eq_degreeOfFamily}.
\end{proof}

\begin{proposition}[sup St gram]\leanlink{ReciprocityTheory.GeometryOfNumbersTheory.supSt_gram}
\begin{quote}\small\raggedright$\displaystyle (a b : \mathrm{Fin} (D + 1) \to \allowbreak \mathbb{N} ) : \int z, \allowbreak \mathrm{starRingEnd} \mathbb{C} (\mathrm{monoC} a z) * \mathrm{monoC} b z \partial (\mathrm{sphProb} (D + 1)) = if a = b \mathrm{then} (\mathrm{monoNormSq} a : \mathbb{C} ) \mathrm{else} 0$\end{quote}
\end{proposition}
\begin{proof}
The argument operates on the structural objects sph Prob \leanlink{ReciprocityTheory.GeometryOfNumbersTheory.sphProb}, mono Norm Sq \leanlink{ReciprocityTheory.GeometryOfNumbersTheory.monoNormSq}, mono C \leanlink{ReciprocityTheory.GeometryOfNumbersTheory.monoC}. Combining \emph{integral mono sq sph Prob} \leanlink{ReciprocityTheory.GeometryOfNumbersTheory.integral_mono_sq_sphProb} with \emph{integral mono sph Prob of ne} \leanlink{ReciprocityTheory.GeometryOfNumbersTheory.integral_mono_sphProb_of_ne} yields the stated result.
\end{proof}

\begin{proposition}[degree Of Family eq of is Full Rank Family]\leanlink{ReciprocityTheory.GeometryOfNumbersTheory.degreeOfFamily_eq_of_isFullRankFamily}
\begin{quote}\small\raggedright$\displaystyle (h : \mathrm{HermFam} K V) \allowbreak (hE : E.\mathrm{FG}) \allowbreak (hv : \mathrm{IsFrame} K E v) \allowbreak (hv' : \mathrm{IsFrame} K E v') : \mathrm{degreeOfFamily} h E v = \mathrm{degreeOfFamily} h E v'$\end{quote}
\end{proposition}
\begin{proof}
The argument operates on the structural objects Herm Fam \leanlink{ReciprocityTheory.GeometryOfNumbersTheory.HermFam}, degree Of Family \leanlink{ReciprocityTheory.GeometryOfNumbersTheory.degreeOfFamily}, Is Frame \leanlink{ReciprocityTheory.GeometryOfNumbersTheory.IsFrame}. Combining \emph{degree Of Family linear Combination} \leanlink{ReciprocityTheory.GeometryOfNumbersTheory.degreeOfFamily_linearCombination} with \emph{rel Index span frame ne zero} \leanlink{ReciprocityTheory.GeometryOfNumbersTheory.IsFrame.relIndex_span_frame_ne_zero} yields the stated result.
\end{proof}

\begin{proposition}[sph Monom key]\leanlink{ReciprocityTheory.GeometryOfNumbersTheory.sphMonom_key}
\begin{quote}\small\raggedright$\displaystyle (f : (\mathrm{Fin} N \to \allowbreak \mathbb{C} ) \to \allowbreak \mathbb{C} ) \allowbreak (hf : \forall (r : \mathbb{R} ) \allowbreak (z : \mathrm{Fin} N \to \allowbreak \mathbb{C} ), \allowbreak f (r \cdot z) = (r : \mathbb{C} ) ^ k * f z) : (\int z, \allowbreak f z * (\pcReal .\mathrm{exp} (-\mathrm{sqN} z) : \mathbb{C} )) * ((\pcReal .\mathrm{Gamma} ((((2 * N - 1 + 0 : \mathbb{N} ) : \mathbb{R} ) + 1) / 2) / 2 : \mathbb{R} ) : \mathbb{C} ) = (\int z, \allowbreak f ((\pcReal .\mathrm{sqrt} (\mathrm{sqN} z))^{-1} \cdot z) * (\pcReal .\mathrm{exp} (-\mathrm{sqN} z) : \mathbb{C} )) * ((\pcReal .\mathrm{Gamma} ((((2 * N - 1 + k : \mathbb{N} ) : \mathbb{R} ) + 1) / 2) / 2 : \mathbb{R} ) : \mathbb{C} )$\end{quote}
\end{proposition}
\begin{proof}
The argument operates on the structural objects sq N \leanlink{ReciprocityTheory.GeometryOfNumbersTheory.sqN}, sph Monom Mu \leanlink{ReciprocityTheory.GeometryOfNumbersTheory.sphMonomMu}, sph Monom L \leanlink{ReciprocityTheory.GeometryOfNumbersTheory.sphMonomL}. Combining \emph{sph Monom radial} \leanlink{ReciprocityTheory.GeometryOfNumbersTheory.sphMonom_radial}, \emph{sph Monom polar mul} \leanlink{ReciprocityTheory.GeometryOfNumbersTheory.sphMonom_polar_mul}, \emph{norm sph Monom L} \leanlink{ReciprocityTheory.GeometryOfNumbersTheory.norm_sphMonomL}, \emph{finrank sph Monom E} \leanlink{ReciprocityTheory.GeometryOfNumbersTheory.finrank_sphMonomE} with \emph{integral sph Monom Mu} \leanlink{ReciprocityTheory.GeometryOfNumbersTheory.integral_sphMonomMu} yields the stated result.
\end{proof}

\begin{lemma}[chi PAux neg log mono Norm Sq le]\leanlink{ReciprocityTheory.GeometryOfNumbersTheory.chiPAux_neg_log_monoNormSq_le}
\begin{quote}\small\raggedright$\displaystyle (b : \mathrm{Fin} (D + 1) \to \allowbreak \mathbb{N} ) : -\pcReal .\mathrm{log} (\mathrm{monoNormSq} b) \le (\sum i, \allowbreak b i : \mathbb{N} ) * \pcReal .\mathrm{log} (D + 1) + D * \pcReal .\mathrm{log} ((\sum i, \allowbreak b i : \mathbb{N} ) + D)$\end{quote}
\end{lemma}
\begin{proof}
The argument operates on the structural objects mono Norm Sq \leanlink{ReciprocityTheory.GeometryOfNumbersTheory.monoNormSq}. Combining \emph{le mono Norm Sq} \leanlink{ReciprocityTheory.GeometryOfNumbersTheory.le_monoNormSq} with \emph{mono Norm Sq pos} \leanlink{ReciprocityTheory.GeometryOfNumbersTheory.monoNormSq_pos} yields the stated result.
\end{proof}

\begin{lemma}[integral bernstein]\leanlink{ReciprocityTheory.GeometryOfNumbersTheory.integral_bernstein}
\begin{quote}\small\raggedright$\displaystyle (D i : \mathbb{N} ) \allowbreak (hi : i \le D) : \int t in (0 : \mathbb{R} )..1, \allowbreak (D.\mathrm{choose} i : \mathbb{R} ) * t ^ i * (1 - t) ^ (D - i) = 1 / (D + 1 : \mathbb{R} )$\end{quote}
\end{lemma}
\begin{proof}
The claim follows directly from \emph{b I closed} \leanlink{ReciprocityTheory.GeometryOfNumbersTheory.bI_closed}.
\end{proof}

\begin{lemma}[integral mono sq sph Prob]\leanlink{ReciprocityTheory.GeometryOfNumbersTheory.integral_mono_sq_sphProb}
\begin{quote}\small\raggedright$\displaystyle (a : \mathrm{Fin} (D + 1) \to \allowbreak \mathbb{N} ) : \int z, \allowbreak \| \mathrm{monoC} a z\| ^ 2 \partial (\mathrm{sphProb} (D + 1)) = \mathrm{monoNormSq} a$\end{quote}
\end{lemma}
\begin{proof}
The argument operates on the structural objects sph Prob \leanlink{ReciprocityTheory.GeometryOfNumbersTheory.sphProb}, sq N \leanlink{ReciprocityTheory.GeometryOfNumbersTheory.sqN}, mono Norm Sq \leanlink{ReciprocityTheory.GeometryOfNumbersTheory.monoNormSq}, mono C \leanlink{ReciprocityTheory.GeometryOfNumbersTheory.monoC}. Combining \emph{sph Monom key} \leanlink{ReciprocityTheory.GeometryOfNumbersTheory.sphMonom_key}, \emph{integral mono sq mul exp} \leanlink{ReciprocityTheory.GeometryOfNumbersTheory.integral_mono_sq_mul_exp}, \emph{sph Monom integral sph Prob} \leanlink{ReciprocityTheory.GeometryOfNumbersTheory.sphMonom_integral_sphProb}, \emph{sph Monom f smul} \leanlink{ReciprocityTheory.GeometryOfNumbersTheory.sphMonom_f_smul} with \emph{sph Monom gamma} \leanlink{ReciprocityTheory.GeometryOfNumbersTheory.sphMonom_gamma} yields the stated result.
\end{proof}

\begin{lemma}[sph Prob compl sphere]\leanlink{ReciprocityTheory.GeometryOfNumbersTheory.sphProb_compl_sphere}
\[
(D : \mathbb{N} ) : \mathrm{sphProb} (D + 1) c = 0
\]
\end{lemma}
\begin{proof}
The argument operates on the structural objects sph Prob \leanlink{ReciprocityTheory.GeometryOfNumbersTheory.sphProb}, sq N \leanlink{ReciprocityTheory.GeometryOfNumbersTheory.sqN}, gauss C \leanlink{ReciprocityTheory.GeometryOfNumbersTheory.gaussC}. The claim follows directly from \emph{sph Monom sq N normalize} \leanlink{ReciprocityTheory.GeometryOfNumbersTheory.sphMonom_sqN_normalize}.
\end{proof}

\begin{lemma}[neg log det le]\leanlink{ReciprocityTheory.GeometryOfNumbersTheory.neg_log_det_le}
\begin{quote}\small\raggedright$\displaystyle (hr : r' \le r) \allowbreak (hG' : G'.\mathrm{PosDef}) \allowbreak (hG : G.\mathrm{PosDef}) \allowbreak (M : \mathrm{Matrix} (\mathrm{Fin} r) \allowbreak (\mathrm{Fin} r') \mathbb{Z} ) \allowbreak (hM : \mathrm{Injective} (M.\mathrm{map} ( : \mathbb{Z} \to \allowbreak \mathbb{Q} )).\mathrm{mulVec}) \allowbreak (N : \mathrm{Matrix} (\mathrm{Fin} r) \allowbreak (\mathrm{Fin} r) \mathbb{Z} ) \allowbreak (hN : N.\mathrm{det} \neq 0) \allowbreak (hA : 0 < A) \allowbreak (\mathrm{hlam} : 0 < \mathrm{lam}) \allowbreak (\mathrm{hnorm} : (((A : \mathbb{C} ) ^ 2) \cdot G' - (M.\mathrm{map} ( : \mathbb{Z} \to \allowbreak \mathbb{C} ))^{\mathrm{H}} * G * M.\mathrm{map} ( : \mathbb{Z} \to \allowbreak \mathbb{C} )).\mathrm{PosSemidef}) \allowbreak (\mathrm{hsmall} : \forall i, \allowbreak (((N.\mathrm{map} ( : \mathbb{Z} \to \allowbreak \mathbb{C} ))^{\mathrm{H}} * G * N.\mathrm{map} ( : \mathbb{Z} \to \allowbreak \mathbb{C} )) i i).re \le \mathrm{lam} ^ 2) : -(1 / 2) * \pcReal .\mathrm{log} G'.\mathrm{det}.\mathrm{re} \le -(1 / 2) * \pcReal .\mathrm{log} \mathrm{re} + r' * \pcReal .\mathrm{log} A + \;\cdots$\end{quote}
\end{lemma}
\begin{proof}
Combining \emph{exists append li} \leanlink{ReciprocityTheory.GeometryOfNumbersTheory.LinearIndependentExtension.exists_append_li}, \emph{det re le mul prod diag} \leanlink{ReciprocityTheory.GeometryOfNumbersTheory.det_re_le_mul_prod_diag} with \emph{pos Def det re pos} \leanlink{ReciprocityTheory.GeometryOfNumbersTheory.posDef_det_re_pos} yields the stated result.
\end{proof}

\begin{lemma}[gram R pos Def]\leanlink{ReciprocityTheory.GeometryOfNumbersTheory.gramR_posDef}
\begin{quote}\small\raggedright$\displaystyle (h : \mathrm{HermFam} K V) \allowbreak (\mathrm{hli} : \mathrm{LinearIndependent} K v) \allowbreak (w : ) : (\mathrm{gramR} h v w).\mathrm{PosDef}$\end{quote}
\end{lemma}
\begin{proof}
The argument operates on the structural objects Herm Fam \leanlink{ReciprocityTheory.GeometryOfNumbersTheory.HermFam}, gram R \leanlink{ReciprocityTheory.GeometryOfNumbersTheory.gramR}. Combining \emph{finsupp sum sum star mul re eq q C} \leanlink{ReciprocityTheory.GeometryOfNumbersTheory.finsupp_sum_sum_star_mul_re_eq_qC}, \emph{finsupp sum sum star mul eq q R} \leanlink{ReciprocityTheory.GeometryOfNumbersTheory.finsupp_sum_sum_star_mul_eq_qR} with \emph{gram pos Def} \leanlink{ReciprocityTheory.GeometryOfNumbersTheory.gram_posDef} yields the stated result.
\end{proof}

\section{Companion Theory / Thickening Estimates: nrm Up}

This chapter encompasses 216 formal declarations (1,829~LOC). Primary proof idioms detected: \emph{iterated theta lift, Cauchy-Schwarz / Parseval, automorphic / spectral theory}.
\begin{quote}\small
\textbf{Formal Interface ($\mathrm{Inputs} \vdash \mathrm{Outputs}$):}
\par\noindent\emph{Inputs:} Sch~\leanlink{GlobalTheory.RealWeil.Sch}, schwartz To L2~\leanlink{GlobalTheory.RealWeil.schwartzToL2}, L2~\leanlink{GlobalTheory.RealWeil.L2}, comp~\leanlink{GlobalTheory.WeilM.IsUnitaryOn.comp}, comp~\leanlink{GlobalTheory.WeilM.IsContOn.comp}, d Rho~\leanlink{GlobalTheory.RealWeil.dRho}.
\par\noindent\emph{Outputs:} ker~\leanlink{CompanionTheory.ThKR.Mp.ker}, nrm Up~\leanlink{CompanionTheory.ThKR.Mp.nrmUp}, sec Der~\leanlink{CompanionTheory.ThKR.Mp.secDer}, sec Ex of~\leanlink{CompanionTheory.ThKR.Mp.secEx_of}, dd Ne~\leanlink{CompanionTheory.ThKR.Mp.ddNe}, top G~\leanlink{CompanionTheory.ThKR.Mp.topG}.
\end{quote}

\subsection{Core Definitions and Structures}
\begin{definition}[qd Lin]\leanlink{CompanionTheory.ThKR.Mp.qdLin}
\[
(f : \mathrm{Sch} l) : \mathrm{Mat} l \to _{l} \mathrm{Sch} l
\]
\end{definition}

\begin{definition}[Good Position]\leanlink{CompanionTheory.ThKR.GoodPosition}
\begin{quote}\small\raggedright$\displaystyle (X Y : \mathrm{Matrix} (\mathrm{Fin} n) \allowbreak (\mathrm{Fin} r) E) \allowbreak (K : \mathrm{Subgroup} (\mathrm{GAd} \sigma J)) : \mathrm{Prop}$\end{quote}
\end{definition}

\begin{definition}[Mem]\leanlink{CompanionTheory.ThKR.Mp.Mem}
\[
(x : \mathrm{Amb} l) : \mathrm{Prop}
\]
\end{definition}

\begin{definition}[eps]\leanlink{CompanionTheory.ThKR.Mp.MpA.eps}
\[
G l
\]
\end{definition}

\begin{definition}[unip Min]\leanlink{CompanionTheory.ThKR.unipMin}
\[
(X : \mathrm{Matrix} (\mathrm{Fin} n) (\mathrm{Fin} r) E) : \mathcal{P} (\mathrm{GAd} \sigma J)
\]
\end{definition}

\begin{definition}[Qd]\leanlink{CompanionTheory.ThKR.Mp.Qd}
\[
(X : \mathrm{Mat} l) (f : \mathrm{Sch} l) : \mathrm{Sch} l
\]
\end{definition}

\subsection{Main Results}
\begin{theorem}[nrm Up]\leanlink{CompanionTheory.ThKR.Mp.nrmUp}
\[
\mathrm{NrmUp} l
\]
\end{theorem}
\begin{proof}
The argument operates on the structural objects Nrm Up \leanlink{CompanionTheory.ThKR.Mp.NrmUp}. Combining \emph{dd Ne} \leanlink{CompanionTheory.ThKR.Mp.ddNe}, \emph{has Deriv At UU} \leanlink{CompanionTheory.ThKR.Mp.hasDerivAt_UU}, \emph{UU d Rho} \leanlink{CompanionTheory.ThKR.Mp.UU_dRho}, \emph{hess} \leanlink{CompanionTheory.ThKR.Mp.hess}, \emph{fou mq fou Inv} \leanlink{CompanionTheory.ThKR.LL.fou_mq_fouInv}, \emph{has Deriv At det} \leanlink{CompanionTheory.ThKR.Mp.hasDerivAt_det}, \emph{trace eq} \leanlink{CompanionTheory.ThKR.Mp.trace_eq} with \emph{upper Part unipotent mul} \leanlink{CompanionTheory.ThKR.Mp.upperPart_unipotent_mul} yields the stated result.
\end{proof}

\begin{theorem}[sec Der]\leanlink{CompanionTheory.ThKR.Mp.secDer}
\[
\mathrm{SecDer} l
\]
\end{theorem}
\begin{proof}
The argument operates on the structural objects qd Lin \leanlink{CompanionTheory.ThKR.Mp.qdLin}, Sec Der \leanlink{CompanionTheory.ThKR.Mp.SecDer}, Qd \leanlink{CompanionTheory.ThKR.Mp.Qd}, sec \leanlink{CompanionTheory.ThKR.Mp.sec}. Combining \emph{comp} \leanlink{GlobalTheory.WeilM.IsUnitaryOn.comp}, \emph{comp} \leanlink{GlobalTheory.WeilM.IsContOn.comp}, \emph{curve} \leanlink{CompanionTheory.ThKR.Mp.MpG.curve}, \emph{has FDeriv At of curves} \leanlink{CompanionTheory.ThKR.Mp.MpG.hasFDerivAt_of_curves} with \emph{qd sum} \leanlink{CompanionTheory.ThKR.Mp.qd_sum} yields the stated result.
\end{proof}

\begin{proposition}[ker]\leanlink{CompanionTheory.ThKR.Mp.ker}
\[
\mathrm{Ker} l
\]
\end{proposition}
\begin{proof}
The argument operates on the structural objects eps \leanlink{CompanionTheory.ThKR.Mp.MpA.eps}, rho \leanlink{GlobalTheory.RealWeil.rho}. Combining \emph{ann gs} \leanlink{CompanionTheory.ThKR.Mp.MpA.ann_gs}, \emph{scalar} \leanlink{CompanionTheory.ThKF.E.Impl.scalar}, \emph{exists rho} \leanlink{CompanionTheory.ThKF.E.exists_rho}, \emph{eps ne} \leanlink{CompanionTheory.ThKR.Mp.MpA.eps_ne}, \emph{dd one} \leanlink{CompanionTheory.ThKR.Mp.MpA.dd_one}, \emph{om S eps} \leanlink{CompanionTheory.ThKR.Mp.MpA.omS_eps}, \emph{to Lp rho} \leanlink{CompanionTheory.ThKF.E.toLp_rho} with \emph{om of om S} \leanlink{CompanionTheory.ThKR.Mp.MpA.om_of_omS} yields the stated result.
\end{proof}

\begin{proposition}[sec Ex of]\leanlink{CompanionTheory.ThKR.Mp.secEx_of}
\[
(h1 : \mathrm{MkElt} l) (h2 : \mathrm{NrmUp} l) : \mathrm{SecEx} l
\]
\end{proposition}
\begin{proof}
The argument operates on the structural objects Nrm Up \leanlink{CompanionTheory.ThKR.Mp.NrmUp}, Impl \leanlink{CompanionTheory.ThKF.E.Impl}. Combining \emph{is Symm inv mul} \leanlink{CompanionTheory.ThKR.Mp.C.isSymm_inv_mul}, \emph{intw of sympl} \leanlink{CompanionTheory.ThKR.LL.intw_of_sympl}, \emph{dd mul eq} \leanlink{CompanionTheory.ThKR.Mp.C.dd_mul_eq}, \emph{is Unit big} \leanlink{CompanionTheory.ThKR.LL.isUnit_big}, \emph{sympl sp} \leanlink{CompanionTheory.ThKR.Mp.C.sympl_sp}, \emph{smul} \leanlink{CompanionTheory.ThKR.LL.Intw.smul}, \emph{sec apply zero} \leanlink{CompanionTheory.ThKR.Mp.C.sec_apply_zero} with \emph{smul} \leanlink{CompanionTheory.ThKR.LL.LD.smul} yields the stated result.
\end{proof}

\begin{proposition}[dd Ne]\leanlink{CompanionTheory.ThKR.Mp.ddNe}
\[
\mathrm{DdNe} l
\]
\end{proposition}
\begin{proof}
The argument operates on the structural objects dd \leanlink{CompanionTheory.ThKR.Mp.dd}, Dd Ne \leanlink{CompanionTheory.ThKR.Mp.DdNe}, upper Parabolic \leanlink{CompanionTheory.ThKR.Mp.upperParabolic}. The claim follows directly from \emph{sp rel} \leanlink{CompanionTheory.ThKR.Mp.sp_rel}.
\end{proof}

\begin{proposition}[top G]\leanlink{CompanionTheory.ThKR.Mp.topG}
\[
\mathrm{TopG} l
\]
\end{proposition}
\begin{proof}
The argument operates on the structural objects om \leanlink{CompanionTheory.ThKR.Mp.om}, G \leanlink{CompanionTheory.ThKR.Mp.G}, Top G \leanlink{CompanionTheory.ThKR.Mp.TopG}. Combining \emph{comp} \leanlink{GlobalTheory.WeilM.IsUnitaryOn.comp} with \emph{comp} \leanlink{GlobalTheory.WeilM.IsContOn.comp} yields the stated result.
\end{proof}

\begin{lemma}[cover of]\leanlink{CompanionTheory.ThKR.Mp.cover_of}
\begin{quote}\small\raggedright$\displaystyle (hT : \mathrm{TopG} l) \allowbreak (hO : \mathrm{Surj} l) \allowbreak (hK : \mathrm{Ker} l) \allowbreak (hS : \mathrm{SecEx} l) \allowbreak (hD : \mathrm{SecDer} l) : \mathrm{Cover} l$\end{quote}
\end{lemma}
\begin{proof}
The argument operates on the structural objects Sec Der \leanlink{CompanionTheory.ThKR.Mp.SecDer}, Sec Ex \leanlink{CompanionTheory.ThKR.Mp.SecEx}, om S \leanlink{CompanionTheory.ThKR.Mp.omS}, sec \leanlink{CompanionTheory.ThKR.Mp.sec}. Combining \emph{comp} \leanlink{GlobalTheory.WeilM.IsUnitaryOn.comp}, \emph{comp} \leanlink{GlobalTheory.WeilM.IsContOn.comp}, \emph{tendsto sec} \leanlink{CompanionTheory.ThKR.Mp.MpE.tendsto_sec}, \emph{metaplectic group t2 Space} \leanlink{CompanionTheory.ThKR.Mp.MpE.metaplectic_group_t2Space}, \emph{eventually det pos} \leanlink{CompanionTheory.ThKR.Mp.MpE.eventually_det_pos} with \emph{is Inducing F} \leanlink{CompanionTheory.ThKR.Mp.MpE.isInducing_F} yields the stated result.
\end{proof}

\begin{lemma}[surj of]\leanlink{CompanionTheory.ThKR.Mp.surj_of}
\[
(h1 : \mathrm{DdNe} l) (h2 : \mathrm{SecEx} l) : \mathrm{Surj} l
\]
\end{lemma}
\begin{proof}
The argument operates on the structural objects Sec Ex \leanlink{CompanionTheory.ThKR.Mp.SecEx}, Dd Ne \leanlink{CompanionTheory.ThKR.Mp.DdNe}, Surj \leanlink{CompanionTheory.ThKR.Mp.Surj}, lo \leanlink{CompanionTheory.ThKR.Mp.C.lo}. Combining \emph{exists det ne} \leanlink{CompanionTheory.ThKR.Mp.C.exists_det_ne} with \emph{lo mem} \leanlink{CompanionTheory.ThKR.Mp.C.lo_mem} yields the stated result.
\end{proof}

\begin{lemma}[mk Elt]\leanlink{CompanionTheory.ThKR.Mp.mkElt}
\[
\mathrm{MkElt} l
\]
\end{lemma}
\begin{proof}
The argument operates on the structural objects Mem \leanlink{CompanionTheory.ThKR.Mp.Mem}, rho \leanlink{GlobalTheory.RealWeil.rho}, d Rho \leanlink{GlobalTheory.RealWeil.dRho}, Impl \leanlink{CompanionTheory.ThKF.E.Impl}. Combining \emph{d Rho} \leanlink{CompanionTheory.ThKF.E.Impl.dRho}, \emph{exists rho} \leanlink{CompanionTheory.ThKF.E.exists_rho}, \emph{to Lp rho} \leanlink{CompanionTheory.ThKF.E.toLp_rho}, \emph{dense sch} \leanlink{CompanionTheory.ThKR.Mp.MpA.dense_sch} with \emph{ext sch} \leanlink{CompanionTheory.ThKR.Mp.MpA.ext_sch} yields the stated result.
\end{proof}

\begin{lemma}[real Metaplectic Representation holds]\leanlink{CompanionTheory.ThKR.realMetaplecticRepresentation_holds}
\[
\mathrm{RealMetaplecticRepresentation}
\]
\end{lemma}
\begin{proof}
Combining \emph{nrm Up} \leanlink{CompanionTheory.ThKR.Mp.nrmUp}, \emph{sec Der} \leanlink{CompanionTheory.ThKR.Mp.secDer}, \emph{ker} \leanlink{CompanionTheory.ThKR.Mp.ker}, \emph{sec Ex of} \leanlink{CompanionTheory.ThKR.Mp.secEx_of}, \emph{dd Ne} \leanlink{CompanionTheory.ThKR.Mp.ddNe}, \emph{top G} \leanlink{CompanionTheory.ThKR.Mp.topG}, \emph{cover of} \leanlink{CompanionTheory.ThKR.Mp.cover_of} with \emph{surj of} \leanlink{CompanionTheory.ThKR.Mp.surj_of} yields the stated result.
\end{proof}

\begin{lemma}[one Par of]\leanlink{CompanionTheory.ThKR.Mp.onePar_of}
\begin{quote}\small\raggedright$\displaystyle (h1 : \mathrm{TopG} l) \allowbreak (h2 : \mathrm{Cover} l) \allowbreak (h3 : \mathrm{NearOne} l) \allowbreak (h4 : \mathrm{SecDer} l) : \mathrm{OnePar} l$\end{quote}
\end{lemma}
\begin{proof}
The argument operates on the structural objects Sec Der \leanlink{CompanionTheory.ThKR.Mp.SecDer}, One Par \leanlink{CompanionTheory.ThKR.Mp.OnePar}, Qd \leanlink{CompanionTheory.ThKR.Mp.Qd}, quad \leanlink{GlobalTheory.RealWeil.quad}. Combining \emph{exp mem sp} \leanlink{CompanionTheory.ThKR.Mp.exp_mem_sp}, \emph{has Deriv At} \leanlink{CompanionTheory.ThKR.LL.LD.hasDerivAt}, \emph{qd sum} \leanlink{CompanionTheory.ThKR.Mp.qd_sum}, \emph{quad eq} \leanlink{CompanionTheory.ThKR.LL.quad_eq} with \emph{has Deriv At lin exp} \leanlink{CompanionTheory.ThKR.LL.hasDerivAt_lin_exp} yields the stated result.
\end{proof}

\begin{lemma}[smooth of]\leanlink{CompanionTheory.ThKR.Mp.smooth_of}
\[
(hT : \mathrm{TopG} l) (hN : \mathrm{NearOne} l) (hD : \mathrm{SecDer} l) : \mathrm{Smooth} l
\]
\end{lemma}
\begin{proof}
The argument operates on the structural objects Sec Der \leanlink{CompanionTheory.ThKR.Mp.SecDer}, Qd \leanlink{CompanionTheory.ThKR.Mp.Qd}, Smooth \leanlink{CompanionTheory.ThKR.Mp.Smooth}, dm \leanlink{CompanionTheory.ThKR.Mp.dm}. Combining \emph{comp} \leanlink{GlobalTheory.WeilM.IsUnitaryOn.comp}, \emph{comp} \leanlink{GlobalTheory.WeilM.IsContOn.comp}, \emph{has FDeriv At orb} \leanlink{CompanionTheory.ThKR.Mp.hasFDerivAt_orb} with \emph{cont Diff dm} \leanlink{CompanionTheory.ThKR.Mp.contDiff_dm} yields the stated result.
\end{proof}

\section{Formalism: host Conj of}

This chapter encompasses 72 formal declarations (1,017~LOC). Primary proof idioms detected: \emph{Vandiver / cohomology, L-function / Rankin-Selberg}.
\begin{quote}\small
\textbf{Formal Interface ($\mathrm{Inputs} \vdash \mathrm{Outputs}$):}
\par\noindent\emph{Inputs:} S~\leanlink{Formalism.S}.
\par\noindent\emph{Outputs:} e3~\leanlink{Formalism.Q3.e3}, stab G~\leanlink{Formalism.Q3.stabG}, not dvd one~\leanlink{Formalism.Q3.CMRootsOfUnity.not_dvd_one}, chi Q~\leanlink{Formalism.Q3.chiQ}, quad Discr~\leanlink{Formalism.HF.quadDiscr}, quad Discr ne one~\leanlink{Formalism.HF.P.quadDiscr_ne_one}.
\end{quote}

\subsection{Core Definitions and Structures}
\begin{definition}[gal Sign]\leanlink{Formalism.Q3.galSign}
\begin{quote}\small\raggedright$\displaystyle (\alpha : \mathrm{Fin} k \to \allowbreak E) \allowbreak (\sigma : \mathrm{Gal}(E/\mathbb{Q} )) : \mathrm{Fin} k \to \allowbreak \mathbb{Z} ^{\times}$\end{quote}
\end{definition}

\begin{definition}[Is CMType Emb]\leanlink{Formalism.Q3.IsCMTypeEmb}
\[
(S : \mathcal{F} (K \to +* \mathbb{C} )) : \mathrm{Prop}
\]
\end{definition}

\begin{definition}[Is Host]\leanlink{Formalism.Q3.IsHost}
\[
(q : \mathrm{Fin} 3 \to \mathbb{N} ) (\alpha : \mathrm{Fin} 3 \to E) : \mathrm{Prop}
\]
\end{definition}

\begin{definition}[Is MQ]\leanlink{Formalism.Q3.IsMQ}
\[
(a : \mathrm{Fin} k \to \mathbb{Z} ) (\alpha : \mathrm{Fin} k \to E) : \mathrm{Prop}
\]
\end{definition}

\begin{definition}[mqk Field]\leanlink{Formalism.Q3.mqkField}
\begin{quote}\small\raggedright$\displaystyle (F : \mathrm{Type}) \allowbreak (\alpha : \mathrm{Fin} k \to \allowbreak E) \allowbreak (T : \mathcal{F} (\mathrm{Fin} k)) : \mathrm{IntermediateField} F E$\end{quote}
\end{definition}

\begin{definition}[cc]\leanlink{Formalism.Q3.cc}
\[
\mathrm{Gal3}
\]
\end{definition}

\subsection{Main Results}
\begin{theorem}[host Conj of]\leanlink{Formalism.Q3.hostConj_of}
\[
(hH : \mathrm{HostGalois}) (hK : \mathrm{MQKummer}) : \mathrm{HostConj}
\]
\end{theorem}
\begin{proof}
The argument operates on the structural objects Host Conj \leanlink{Formalism.Q3.HostConj}, MQKummer \leanlink{Formalism.Q3.MQKummer}, Host Galois \leanlink{Formalism.Q3.HostGalois}, gal Sign \leanlink{Formalism.Q3.galSign}. Combining \emph{gal Sign eq cc iff} \leanlink{Formalism.Q3.CMConjugation.galSign_eq_cc_iff}, \emph{part1} \leanlink{Formalism.Q3.CMConjugation.part1}, \emph{phi conj} \leanlink{Formalism.Q3.CMConjugation.phi_conj}, \emph{not sq mul} \leanlink{Formalism.Q3.CMConjugation.not_sq_mul} with \emph{alpha sq} \leanlink{Formalism.Q3.CMConjugation.alpha_sq} yields the stated result.
\end{proof}

\begin{theorem}[mq Kummer of]\leanlink{Formalism.Q3.mqKummer_of}
\[
(hG : \mathrm{MQGalois}) : \mathrm{MQKummer}
\]
\end{theorem}
\begin{proof}
The argument operates on the structural objects MQGalois \leanlink{Formalism.Q3.MQGalois}, MQKummer \leanlink{Formalism.Q3.MQKummer}, mqk Field \leanlink{Formalism.Q3.mqkField}. The claim follows directly from \emph{mqk induction} \leanlink{Formalism.Q3.mqk_induction}.
\end{proof}

\begin{proposition}[gal Sign eq cc iff]\leanlink{Formalism.Q3.CMConjugation.galSign_eq_cc_iff}
\begin{quote}\small\raggedright$\displaystyle (hq : \mathrm{Hyp3} q) \allowbreak (h\alpha : \mathrm{IsHost} q \alpha ) \allowbreak (\sigma : \mathrm{Gal}(E/\mathbb{Q} )) : \mathrm{galSign} \alpha \sigma = cc \iff \forall i, \allowbreak \sigma (\alpha i) = -\alpha i$\end{quote}
\end{proposition}
\begin{proof}
The argument operates on the structural objects gal Sign \leanlink{Formalism.Q3.galSign}, Is Host \leanlink{Formalism.Q3.IsHost}, cc \leanlink{Formalism.Q3.cc}, Hyp3 \leanlink{Formalism.Q3.Hyp3}. Combining \emph{alpha ne zero} \leanlink{Formalism.Q3.CMConjugation.alpha_ne_zero} with \emph{sigma alpha} \leanlink{Formalism.Q3.CMConjugation.sigma_alpha} yields the stated result.
\end{proof}

\begin{proposition}[mqk induction]\leanlink{Formalism.Q3.mqk_induction}
\begin{quote}\small\raggedright$\displaystyle (F : \mathrm{Type}) \allowbreak (a : \mathrm{Fin} k \to \allowbreak \mathbb{Z} ) \allowbreak (\alpha : \mathrm{Fin} k \to \allowbreak E) \allowbreak (h\alpha : \forall i, \allowbreak \alpha i ^ 2 = (a i : E)) : \forall T : \mathcal{F} (\mathrm{Fin} k), \allowbreak \mathrm{finrank} F (\mathrm{mqkField} F \alpha T) = 2 ^ T.\mathrm{card} \land \allowbreak \forall d : F, \allowbreak \forall y \in \mathrm{mqkField} F \alpha T, \allowbreak y ^ 2 = \mathrm{algebraMap} F E d \to \allowbreak \exists S \subseteq T, \allowbreak \mathrm{IsSquare} (d * \prod i \in S, \allowbreak (a i : F))$\end{quote}
\end{proposition}
\begin{proof}
The argument operates on the structural objects mqk Field \leanlink{Formalism.Q3.mqkField}. Combining \emph{mqk sq step} \leanlink{Formalism.Q3.mqk_sq_step}, \emph{mqk finrank two} \leanlink{Formalism.Q3.mqk_finrank_two} with \emph{mqk insert} \leanlink{Formalism.Q3.mqk_insert} yields the stated result.
\end{proof}

\begin{proposition}[host Galois of]\leanlink{Formalism.Q3.hostGalois_of}
\[
(hG : \mathrm{MQGalois}) (hK : \mathrm{MQKummer}) : \mathrm{HostGalois}
\]
\end{proposition}
\begin{proof}
The argument operates on the structural objects MQGalois \leanlink{Formalism.Q3.MQGalois}, MQKummer \leanlink{Formalism.Q3.MQKummer}, Host Galois \leanlink{Formalism.Q3.HostGalois}, Is MQ \leanlink{Formalism.Q3.IsMQ}. The claim follows directly from \emph{host not sq} \leanlink{Formalism.Q3.host_not_sq}.
\end{proof}

\begin{proposition}[host Emb of]\leanlink{Formalism.Q3.hostEmb_of}
\[
(hC : \mathrm{HostConj}) (hH : \mathrm{HostGalois}) : \mathrm{HostEmb}
\]
\end{proposition}
\begin{proof}
The argument operates on the structural objects Host Emb \leanlink{Formalism.Q3.HostEmb}, Host Conj \leanlink{Formalism.Q3.HostConj}, Host Galois \leanlink{Formalism.Q3.HostGalois}, gal Sign \leanlink{Formalism.Q3.galSign}. This result is established directly from standard library foundations within its local module (45~lines of formal proof).
\end{proof}

\begin{lemma}[mqk sq step]\leanlink{Formalism.Q3.mqk_sq_step}
\begin{quote}\small\raggedright$\displaystyle (K : \mathrm{IntermediateField} F E) \allowbreak (\beta b d : E) \allowbreak (hb : b \in K) \allowbreak (hd : d \in K) \allowbreak (h\beta : \beta ^ 2 = b) \allowbreak (h\beta K : \beta \notin K) \allowbreak (hy : y \in \mathrm{adjoin} K ) \allowbreak (\mathrm{hyd} : y ^ 2 = d) : (\exists z \in K, \allowbreak z ^ 2 = d) \lor \allowbreak (\exists z \in K, \allowbreak z ^ 2 = d * b)$\end{quote}
\end{lemma}
\begin{proof}
The claim follows directly from \emph{mqk mem adjoin form} \leanlink{Formalism.Q3.mqk_mem_adjoin_form}.
\end{proof}

\begin{lemma}[mq is Galois]\leanlink{Formalism.Q3.mq_isGalois}
\[
(h : \mathrm{IsMQ} a \alpha ) : \mathrm{IsGalois} \mathbb{Q} E
\]
\end{lemma}
\begin{proof}
The argument operates on the structural objects Is MQ \leanlink{Formalism.Q3.IsMQ}. The claim follows directly from \emph{mq alg Adjoin} \leanlink{Formalism.Q3.mq_algAdjoin}.
\end{proof}

\begin{lemma}[part1]\leanlink{Formalism.Q3.CMConjugation.part1}
\begin{quote}\small\raggedright$\displaystyle (hq : \mathrm{Hyp3} q) \allowbreak (h\alpha : \mathrm{IsHost} q \alpha ) \allowbreak (\varphi : E \to \allowbreak +* \mathbb{C} ) \allowbreak (\sigma : \mathrm{Gal}(E/\mathbb{Q} )) : (\forall x, \allowbreak \varphi (\sigma x) = \mathrm{starRingEnd} \mathbb{C} (\varphi x)) \iff \mathrm{galSign} \alpha \sigma = cc$\end{quote}
\end{lemma}
\begin{proof}
The argument operates on the structural objects gal Sign \leanlink{Formalism.Q3.galSign}, Is Host \leanlink{Formalism.Q3.IsHost}, cc \leanlink{Formalism.Q3.cc}, Hyp3 \leanlink{Formalism.Q3.Hyp3}. Combining \emph{gal Sign eq cc iff} \leanlink{Formalism.Q3.CMConjugation.galSign_eq_cc_iff}, \emph{phi conj} \leanlink{Formalism.Q3.CMConjugation.phi_conj} with \emph{alg Adjoin eq top} \leanlink{Formalism.Q3.CMConjugation.algAdjoin_eq_top} yields the stated result.
\end{proof}

\begin{lemma}[quad Discr]\leanlink{Formalism.HF.quadDiscr}
\[
\mathrm{QuadDiscr}
\]
\end{lemma}
\begin{proof}
The argument operates on the structural objects Quad Discr \leanlink{Formalism.HF.QuadDiscr}. The claim follows directly from \emph{quad Discr ne one} \leanlink{Formalism.HF.P.quadDiscr_ne_one}.
\end{proof}

\begin{lemma}[mq Galois]\leanlink{Formalism.Q3.mqGalois}
\[
\mathrm{MQGalois}
\]
\end{lemma}
\begin{proof}
The argument operates on the structural objects MQGalois \leanlink{Formalism.Q3.MQGalois}. Combining \emph{mq is Galois} \leanlink{Formalism.Q3.mq_isGalois}, \emph{mq gal Sign mul} \leanlink{Formalism.Q3.mq_galSign_mul}, \emph{mq sigma} \leanlink{Formalism.Q3.mq_sigma} with \emph{mq injective} \leanlink{Formalism.Q3.mq_injective} yields the stated result.
\end{proof}

\begin{lemma}[mq gal Sign mul]\leanlink{Formalism.Q3.mq_galSign_mul}
\begin{quote}\small\raggedright$\displaystyle (h : \mathrm{IsMQ} a \alpha ) \allowbreak (ha : \forall i, \allowbreak a i \neq 0) \allowbreak (\sigma \tau : \mathrm{Gal}(E/\mathbb{Q} )) : \mathrm{galSign} \alpha (\sigma * \tau ) = \mathrm{galSign} \alpha \sigma * \mathrm{galSign} \alpha \tau$\end{quote}
\end{lemma}
\begin{proof}
The argument operates on the structural objects gal Sign \leanlink{Formalism.Q3.galSign}, Is MQ \leanlink{Formalism.Q3.IsMQ}. Combining \emph{mq sigma} \leanlink{Formalism.Q3.mq_sigma} with \emph{mq alpha ne zero} \leanlink{Formalism.Q3.mq_alpha_ne_zero} yields the stated result.
\end{proof}

\section{Formalism: flag Deg Formula statement}

This chapter encompasses 67 formal declarations (462~LOC). Primary proof idioms detected: \emph{determinant / Stepanov}.
\begin{quote}\small
\textbf{Formal Interface ($\mathrm{Inputs} \vdash \mathrm{Outputs}$):}
\par\noindent\emph{Inputs:} Upq~\leanlink{Formalism.UG.Upq}.
\par\noindent\emph{Outputs:} h P~\leanlink{Formalism.DegreeFormula.hP}, h Q~\leanlink{Formalism.DegreeFormula.hQ}, Dom Flat~\leanlink{Formalism.DegreeFormula.DomFlat}, unflat~\leanlink{Formalism.DegreeFormula.unflat}, flat~\leanlink{Formalism.DegreeFormula.flat}, moeb~\leanlink{Formalism.MJ.moeb}.
\end{quote}

\subsection{Core Definitions and Structures}
\begin{definition}[det Char]\leanlink{Formalism.DegreeFormula.detChar}
\begin{quote}\small\raggedright$\displaystyle (\iota : \mathrm{Type}) \allowbreak (p q : \mathbb{N} ) \allowbreak (\mathrm{k\_1} \mathrm{k\_2} : \iota \to \allowbreak \mathbb{Z} ) : (\iota \to \allowbreak GL (\mathrm{Fin} p) \mathbb{C} \times GL (\mathrm{Fin} q) \mathbb{C} ) \to \allowbreak * GL (\mathrm{Fin} 1) \mathbb{C}$\end{quote}
\end{definition}

\begin{definition}[kap Q]\leanlink{Formalism.DegreeFormula.kapQ}
\begin{quote}\small\raggedright$\displaystyle (g : GL (\mathrm{Fin} p \oplus \mathrm{Fin} q) \mathbb{C} ) \allowbreak (z : \mathrm{Matrix} (\mathrm{Fin} p) \allowbreak (\mathrm{Fin} q) \mathbb{C} ) : \mathrm{Matrix} (\mathrm{Fin} q) \allowbreak (\mathrm{Fin} q) \mathbb{C}$\end{quote}
\end{definition}

\begin{definition}[moeb I]\leanlink{Formalism.DegreeFormula.moebI}
\begin{quote}\small\raggedright$\displaystyle (g : \iota \to \allowbreak GL (\mathrm{Fin} p \oplus \mathrm{Fin} q) \mathbb{C} ) \allowbreak (z : \iota \to \allowbreak \mathrm{Matrix} (\mathrm{Fin} p) \allowbreak (\mathrm{Fin} q) \mathbb{C} ) : \iota \to \allowbreak \mathrm{Matrix} (\mathrm{Fin} p) \allowbreak (\mathrm{Fin} q) \mathbb{C}$\end{quote}
\end{definition}

\begin{definition}[moeb]\leanlink{Formalism.MJ.moeb}
\begin{quote}\small\raggedright$\displaystyle (g : \mathrm{Matrix} (\mathrm{Fin} p \oplus \mathrm{Fin} q) \allowbreak (\mathrm{Fin} p \oplus \mathrm{Fin} q) \mathbb{C} ) \allowbreak (z : \mathrm{Matrix} (\mathrm{Fin} p) \allowbreak (\mathrm{Fin} q) \mathbb{C} ) : \mathrm{Matrix} (\mathrm{Fin} p) \allowbreak (\mathrm{Fin} q) \mathbb{C}$\end{quote}
\end{definition}

\begin{definition}[hc Dom]\leanlink{Formalism.MJ.hcDom}
\begin{quote}\small\raggedright$\displaystyle (p q : \mathbb{N} ) : \mathcal{P} (\mathrm{Matrix} (\mathrm{Fin} p) \allowbreak (\mathrm{Fin} q) \mathbb{C} )$\end{quote}
\end{definition}

\begin{definition}[unflat]\leanlink{Formalism.DegreeFormula.unflat}
\begin{quote}\small\raggedright$\displaystyle (v : \iota \times \mathrm{Fin} p \times \mathrm{Fin} q \to \allowbreak \mathbb{C} ) : \iota \to \allowbreak \mathrm{Matrix} (\mathrm{Fin} p) \allowbreak (\mathrm{Fin} q) \mathbb{C}$\end{quote}
\end{definition}

\subsection{Main Results}
\begin{theorem}[flag Deg Formula statement]\leanlink{Formalism.DegreeFormula.flagDegFormula_statement}
\[
\mathrm{FlagDegFormula}
\]
\end{theorem}
\begin{proof}
The argument operates on the structural objects dim Fl \leanlink{Formalism.DegreeFormula.dimFl}, Flag Deg Formula \leanlink{Formalism.DegreeFormula.FlagDegFormula}, deg Fl \leanlink{Formalism.DegreeFormula.degFl}, flag Deg \leanlink{Formalism.DegreeFormula.flagDeg}. Combining \emph{coeff pow} \leanlink{Formalism.DegreeFormula.PFlagDeg.coeff_pow}, \emph{vandermonde sum} \leanlink{Formalism.DegreeFormula.PFlagDeg.vandermonde_sum} with \emph{pairs nat} \leanlink{Formalism.DegreeFormula.PFlagDeg.pairs_nat} yields the stated result.
\end{proof}

\begin{theorem}[deg Fl Esymm]\leanlink{Formalism.DegreeFormula.PPowSum.degFlEsymm}
\[
\mathrm{DegFlEsymm}
\]
\end{theorem}
\begin{proof}
The argument operates on the structural objects Deg Fl Esymm \leanlink{Formalism.DegreeFormula.DegFlEsymm}, deg Fl \leanlink{Formalism.DegreeFormula.degFl}. Combining \emph{deg Fl Power Sum} \leanlink{Formalism.DegreeFormula.PPowSum.degFlPowerSum}, \emph{Z sum} \leanlink{Formalism.DegreeFormula.PPowSum.Z_sum} with \emph{Z mul left} \leanlink{Formalism.DegreeFormula.PPowSum.Z_mul_left} yields the stated result.
\end{proof}

\begin{proposition}[slice]\leanlink{Formalism.DegreeFormula.PPowSum.slice}
\begin{quote}\small\raggedright$\displaystyle (\tau : \iota ) \allowbreak (k : \mathbb{N} ) \allowbreak (hk : 1 \le k) \allowbreak (P : \mathrm{MvPolynomial} (\iota \times \mathrm{Fin} N) \mathbb{R} ) \allowbreak (j : \mathrm{Fin} N) : \sum \sigma : \iota \to \allowbreak \pcEquiv .\mathrm{Perm} (\mathrm{Fin} N), \allowbreak \mathrm{sgnP} \sigma * (X (\tau , \allowbreak (\sigma \tau ).\mathrm{symm} j) ^ k * P).\mathrm{coeff} (\mathrm{permExp} \sigma ) = 0$\end{quote}
\end{proposition}
\begin{proof}
The argument operates on the structural objects sgn P \leanlink{Formalism.DegreeFormula.PPowSum.sgnP}, perm Exp \leanlink{Formalism.DegreeFormula.permExp}. Combining \emph{perm Sign mul Single swap mul} \leanlink{Formalism.DegreeFormula.PPowSum.permSign_mulSingle_swap_mul}, \emph{rem eq} \leanlink{Formalism.DegreeFormula.PPowSum.rem_eq}, \emph{mul Single mul apply} \leanlink{Formalism.DegreeFormula.PPowSum.mulSingle_mul_apply}, \emph{coeff Xpow mul} \leanlink{Formalism.DegreeFormula.PPowSum.coeff_Xpow_mul}, \emph{sum eq zero of anti} \leanlink{Formalism.DegreeFormula.PPowSum.sum_eq_zero_of_anti} with \emph{perm Exp apply} \leanlink{Formalism.DegreeFormula.PPowSum.permExp_apply} yields the stated result.
\end{proof}

\begin{proposition}[degree Esymm Mul Eq Zero]\leanlink{Formalism.DegreeFormula.degreeEsymmMulEqZero}
\[
\mathrm{DegFlEsymm}
\]
\end{proposition}
\begin{proof}
The argument operates on the structural objects Deg Fl Esymm \leanlink{Formalism.DegreeFormula.DegFlEsymm}. The claim follows directly from \emph{deg Fl Esymm} \leanlink{Formalism.DegreeFormula.PPowSum.degFlEsymm}.
\end{proof}

\begin{proposition}[coeff pow]\leanlink{Formalism.DegreeFormula.PFlagDeg.coeff_pow}
\begin{quote}\small\raggedright$\displaystyle (\mathrm{lam} : \iota \to \allowbreak \mathrm{Fin} N \to \allowbreak \mathbb{R} ) \allowbreak (\sigma : \iota \to \allowbreak \pcEquiv .\mathrm{Perm} (\mathrm{Fin} N)) : (\mathrm{linRoot} \mathrm{lam} ^ \mathrm{dimFl} \iota N).\mathrm{coeff} (\mathrm{permExp} \sigma ) = (-1) ^ \mathrm{dimFl} \iota N * * \prod \tau , \allowbreak \prod i : \mathrm{Fin} N, \allowbreak \mathrm{lam} \tau i ^ (N - 1 - ((\sigma \tau i : \mathrm{Fin} N) : \mathbb{N} ))$\end{quote}
\end{proposition}
\begin{proof}
The argument operates on the structural objects dim Fl \leanlink{Formalism.DegreeFormula.dimFl}, perm Exp \leanlink{Formalism.DegreeFormula.permExp}, lin Root \leanlink{Formalism.DegreeFormula.linRoot}, sfac \leanlink{Formalism.DegreeFormula.PFlagDeg.sfac}. Combining \emph{perm Exp sum'} \leanlink{Formalism.DegreeFormula.PFlagDeg.permExp_sum'}, \emph{prod perm rev} \leanlink{Formalism.DegreeFormula.PFlagDeg.prod_perm_rev}, \emph{lin Root eq} \leanlink{Formalism.DegreeFormula.PFlagDeg.linRoot_eq} with \emph{perm Exp apply} \leanlink{Formalism.DegreeFormula.PFlagDeg.permExp_apply} yields the stated result.
\end{proof}

\begin{proposition}[deg Fl Homog]\leanlink{Formalism.DegreeFormula.PPowSum.degFlHomog}
\[
\mathrm{DegFlHomog}
\]
\end{proposition}
\begin{proof}
The argument operates on the structural objects Deg Fl Homog \leanlink{Formalism.DegreeFormula.DegFlHomog}, deg Fl \leanlink{Formalism.DegreeFormula.degFl}. The claim follows directly from \emph{degree perm Exp} \leanlink{Formalism.DegreeFormula.PPowSum.degree_permExp}.
\end{proof}

\begin{lemma}[deg Fl Power Sum]\leanlink{Formalism.DegreeFormula.PPowSum.degFlPowerSum}
\[
\mathrm{DegFlPowerSum}
\]
\end{lemma}
\begin{proof}
The argument operates on the structural objects Deg Fl Power Sum \leanlink{Formalism.DegreeFormula.DegFlPowerSum}, deg Fl \leanlink{Formalism.DegreeFormula.degFl}, sgn P \leanlink{Formalism.DegreeFormula.PPowSum.sgnP}, perm Exp \leanlink{Formalism.DegreeFormula.permExp}. The claim follows directly from \emph{slice} \leanlink{Formalism.DegreeFormula.PPowSum.slice}.
\end{proof}

\begin{lemma}[vandermonde sum]\leanlink{Formalism.DegreeFormula.PFlagDeg.vandermonde_sum}
\begin{quote}\small\raggedright$\displaystyle (v : \mathrm{Fin} N \to \allowbreak \mathbb{R} ) : \sum s : \pcEquiv .\mathrm{Perm} (\mathrm{Fin} N), \allowbreak ((\pcEquiv .\mathrm{sign} s : \mathbb{Z} ) : \mathbb{R} ) * \prod i, \allowbreak v i ^ (N - 1 - ((s i : \mathrm{Fin} N) : \mathbb{N} )) = \prod i : \mathrm{Fin} N, \allowbreak \prod k : \mathrm{Fin} N, \allowbreak if i < k \mathrm{then} v i - v k \mathrm{else} 1$\end{quote}
\end{lemma}
\begin{proof}
The claim follows directly from \emph{pairs rev} \leanlink{Formalism.DegreeFormula.PFlagDeg.pairs_rev}.
\end{proof}

\begin{lemma}[perm Exp sum']\leanlink{Formalism.DegreeFormula.PFlagDeg.permExp_sum'}
\begin{quote}\small\raggedright$\displaystyle (\sigma : \iota \to \allowbreak \pcEquiv .\mathrm{Perm} (\mathrm{Fin} N)) : \sum x : \iota \times \mathrm{Fin} N, \allowbreak \mathrm{permExp} \sigma x = \mathrm{dimFl} \iota N$\end{quote}
\end{lemma}
\begin{proof}
The argument operates on the structural objects dim Fl \leanlink{Formalism.DegreeFormula.dimFl}, perm Exp \leanlink{Formalism.DegreeFormula.permExp}. Combining \emph{sum perm rev} \leanlink{Formalism.DegreeFormula.PFlagDeg.sum_perm_rev} with \emph{perm Exp apply} \leanlink{Formalism.DegreeFormula.PFlagDeg.permExp_apply} yields the stated result.
\end{proof}

\begin{lemma}[pairs nat]\leanlink{Formalism.DegreeFormula.PFlagDeg.pairs_nat}
\begin{quote}\small\raggedright$\displaystyle (N : \mathbb{N} ) : \prod i : \mathrm{Fin} N, \allowbreak \prod k : \mathrm{Fin} N, \allowbreak (if i < k \mathrm{then} ((k : \mathbb{R} ) - (i : \mathbb{R} )) \mathrm{else} 1) = (\mathrm{sfac} N : \mathbb{R} )$\end{quote}
\end{lemma}
\begin{proof}
The argument operates on the structural objects sfac \leanlink{Formalism.DegreeFormula.PFlagDeg.sfac}. This result is established directly from standard library foundations within its local module (24~lines of formal proof).
\end{proof}

\begin{lemma}[Z sum]\leanlink{Formalism.DegreeFormula.PPowSum.Z_sum}
\begin{quote}\small\raggedright$\displaystyle (s : \mathcal{F} \alpha ) \allowbreak (g : \alpha \to \allowbreak \mathrm{MvPolynomial} (\iota \times \mathrm{Fin} N) \mathbb{R} ) \allowbreak (h : \forall a \in s, \allowbreak \forall P, \allowbreak \mathrm{degFl} \iota N (g a * P) = 0) : \forall P, \allowbreak \mathrm{degFl} \iota N ((\sum a \in s, \allowbreak g a) * P) = 0$\end{quote}
\end{lemma}
\begin{proof}
The argument operates on the structural objects deg Fl \leanlink{Formalism.DegreeFormula.degFl}. The claim follows directly from \emph{deg Fl sum} \leanlink{Formalism.DegreeFormula.PPowSum.degFl_sum}.
\end{proof}

\begin{lemma}[degree perm Exp]\leanlink{Formalism.DegreeFormula.PPowSum.degree_permExp}
\begin{quote}\small\raggedright$\displaystyle (\sigma : \iota \to \allowbreak \pcEquiv .\mathrm{Perm} (\mathrm{Fin} N)) : (\mathrm{permExp} \sigma ).\mathrm{degree} = \mathrm{dimFl} \iota N$\end{quote}
\end{lemma}
\begin{proof}
The argument operates on the structural objects dim Fl \leanlink{Formalism.DegreeFormula.dimFl}, perm Exp \leanlink{Formalism.DegreeFormula.permExp}. The claim follows directly from \emph{perm Exp apply} \leanlink{Formalism.DegreeFormula.PPowSum.permExp_apply}.
\end{proof}

\appendix
\section{Symbol Glossary}
Table~\ref{tab:glossary} anchors each global LaTeX notation macro to its underlying Lean~4 definition and source module.

\begin{table}[ht]
\centering
\caption{Global Symbol Registry}\label{tab:glossary}
\small
\begin{tabular}{llp{0.34\textwidth}p{0.30\textwidth}}
\toprule
Symbol & Macro & Lean Constant & Module \\
\midrule
$\pcHSp$ & \texttt{\detokenize{\pcHSp}} & \raggedright \leanlink{Formalism.BasicNotions.TH.HSp} & \raggedright\arraybackslash \leanmod{Formalism.LocalRingCardinality.DataTheta} \\
$\pcHost$ & \texttt{\detokenize{\pcHost}} & \raggedright \leanlink{ReciprocityTheory.ChartRestriction.G.Geo.Tor.Host} & \raggedright\arraybackslash \leanmod{ReciprocityTheory.ReductiveGroups.ArithmeticGeometry.BirationalGeometry.ToroidalCompactification.TwoWeightEqOneAddRoot} \\
$\pcDatum$ & \texttt{\detokenize{\pcDatum}} & \raggedright \leanlink{ReciprocityTheory.ChartRestriction.G.Geo.Tor.Datum} & \raggedright\arraybackslash \leanmod{ReciprocityTheory.ReductiveGroups.ArithmeticGeometry.BirationalGeometry.ToroidalCompactification.TwoWeightEqOneAddRoot} \\
$\pcSpMap$ & \texttt{\detokenize{\pcSpMap}} & \raggedright \leanlink{ReciprocityTheory.ChartRestriction.G.Geo.Tor.spMap} & \raggedright\arraybackslash \leanmod{ReciprocityTheory.ReductiveGroups.ArithmeticGeometry.BirationalGeometry.ToroidalCompactification.IsCompletedStrictHens} \\
$\pcLevel$ & \texttt{\detokenize{\pcLevel}} & \raggedright \leanlink{ReciprocityTheory.ChartRestriction.G.Geo.Tor.Level} & \raggedright\arraybackslash \leanmod{ReciprocityTheory.ReductiveGroups.ArithmeticGeometry.BirationalGeometry.ToroidalCompactification.TwoWeightEqOneAddRoot} \\
$\pcModuli$ & \texttt{\detokenize{\pcModuli}} & \raggedright \leanlink{ReciprocityTheory.ChartRestriction.G.Geo.Tor.Moduli} & \raggedright\arraybackslash \leanmod{ReciprocityTheory.ReductiveGroups.ArithmeticGeometry.BirationalGeometry.ToroidalCompactification.PushInvertible} \\
$\pcHostA$ & \texttt{\detokenize{\pcHostA}} & \raggedright \leanlink{GlobalTheory.HostModel.host} & \raggedright\arraybackslash \leanmod{Development.AbelianVarietyTheory.Definitions.ArchimedeanThetaValues.OneAddZeta21PowSevenAddPowFourteenEqZero} \\
$\pcCoeff$ & \texttt{\detokenize{\pcCoeff}} & \raggedright \leanlink{ReciprocityTheory.ChartRestriction.G.Geo.Tor.Coeff} & \raggedright\arraybackslash \leanmod{ReciprocityTheory.ReductiveGroups.ArithmeticGeometry.BirationalGeometry.ToroidalCompactification.PushInvertible} \\
$\pcOver$ & \texttt{\detokenize{\pcOver}} & \raggedright \leanlink{GlobalTheory.Ex.ArcF.Over} & \raggedright\arraybackslash \leanmod{Development.AbelianVarietyTheory.Definitions.Degenerations.CharZeroOfIsPrimitiveRoot} \\
$\pcHostB$ & \texttt{\detokenize{\pcHostB}} & \raggedright \leanlink{Formalism.L2.GlobalTheoryTheory.Arith.host} & \raggedright\arraybackslash \leanmod{Formalism.LocalToGlobal.GlobalTheory.IsResidueClassZero} \\
$\pcRank$ & \texttt{\detokenize{\pcRank}} & \raggedright \leanlink{Formalism.BasicNotions.HLat.rank} & \raggedright\arraybackslash \leanmod{Formalism.LocalRingCardinality.DataArith} \\
$\pcRankA$ & \texttt{\detokenize{\pcRankA}} & \raggedright \leanlink{ReciprocityTheory.ChartRestriction.G.Out.DegPre.rank} & \raggedright\arraybackslash \leanmod{ReciprocityTheory.ReductiveGroups.ArithmeticGeometry.ExtensionDescent.ThickTransComp} \\
$\pcWhat$ & \texttt{\detokenize{\pcWhat}} & \raggedright \leanlink{Formalism.L2.Indef.Examples₂.one.what₂} & \raggedright\arraybackslash \leanmod{Formalism.LocalToGlobal.IndefiniteForms.IndefExample} \\
$\pcWhatA$ & \texttt{\detokenize{\pcWhatA}} & \raggedright \leanlink{Formalism.L2.Indef.Examples₂.two.what₂} & \raggedright\arraybackslash \leanmod{Formalism.LocalToGlobal.IndefiniteForms.IndefExample} \\
$\pcHom$ & \texttt{\detokenize{\pcHom}} & \raggedright \leanlink{GlobalTheory.Serre.Pts.hom} & \raggedright\arraybackslash \leanmod{Development.AbelianVarietyTheory.Definitions.ShimuraVarieties.PowObjGeometricallyConnected} \\
$\pcComp$ & \texttt{\detokenize{\pcComp}} & \raggedright \leanlink{Formalism.CoboundaryTransfers.Hom.comp} & \raggedright\arraybackslash \leanmod{Formalism.CoboundaryTransfers.Statement} \\
$\pcHidaSetting$ & \texttt{\detokenize{\pcHidaSetting}} & \raggedright \leanlink{Formalism.L2.SG.HidaSetting.HidaSetting} & \raggedright\arraybackslash \leanmod{Formalism.LocalToGlobal.SiegelLatticeTheory.HaarTransversals.HidaSetting} \\
$\pcHomA$ & \texttt{\detokenize{\pcHomA}} & \raggedright \leanlink{GlobalTheory.Serre.Lie.hom} & \raggedright\arraybackslash \leanmod{Development.AbelianVarietyTheory.Proofs.HomogeneousIdeals.HomLieEquivScalarTmul} \\
$\pcPeriodData$ & \texttt{\detokenize{\pcPeriodData}} & \raggedright \leanlink{Formalism.WeaklyAdmissibleFiltrations.PeriodData} & \raggedright\arraybackslash \leanmod{Formalism.WeaklyAdmissibleFiltrations.Statement} \\
$\pcIsLB$ & \texttt{\detokenize{\pcIsLB}} & \raggedright \leanlink{ReciprocityTheory.ChartRestriction.G.Semi.LB.IsLB} & \raggedright\arraybackslash \leanmod{ReciprocityTheory.ReductiveGroups.ArithmeticGeometry.SemiabelianSchemes.CubicalStructures.CubeSymmetricSwapFirst} \\
$\pcSizeData$ & \texttt{\detokenize{\pcSizeData}} & \raggedright \leanlink{Formalism.BasicNotions.SizeData₂} & \raggedright\arraybackslash \leanmod{Formalism.LocalRingCardinality.DataParams₂} \\
$\FourierTransform$ & \texttt{\detokenize{\FourierTransform}} & \raggedright \leanlink{Formalism.FourierDecayAlt.B.Mat} & \raggedright\arraybackslash \leanmod{Formalism.HermitianCurvature.StatementB} \\
$\pcHostC$ & \texttt{\detokenize{\pcHostC}} & \raggedright \leanlink{GlobalTheory.Mem.LevCh.host} & \raggedright\arraybackslash \leanmod{Development.AbelianVarietyTheory.Definitions.AutomorphicMembership.DuStarPos} \\
$\pcData$ & \texttt{\detokenize{\pcData}} & \raggedright \leanlink{ReciprocityTheory.ChartRestriction.G.FD.Data} & \raggedright\arraybackslash \leanmod{ReciprocityTheory.ReductiveGroups.ArithmeticGeometry.FundamentalDomains.Signature} \\
$\pcRanks$ & \texttt{\detokenize{\pcRanks}} & \raggedright \leanlink{Formalism.BasicNotions.Ranks} & \raggedright\arraybackslash \leanmod{Formalism.LocalRingCardinality.Params} \\
$\pcCompA$ & \texttt{\detokenize{\pcCompA}} & \raggedright \leanlink{Formalism.MainNotions.G.Levels.comp} & \raggedright\arraybackslash \leanmod{Formalism.SharpRestriction.GrowthBounds.DataGChain} \\
$\pcFinExt$ & \texttt{\detokenize{\pcFinExt}} & \raggedright \leanlink{Formalism.BasicNotions.Loc.FinExt} & \raggedright\arraybackslash \leanmod{Formalism.SharpRestriction.LevelLetterCocycles.DataNLoc} \\
$\pcFamIn$ & \texttt{\detokenize{\pcFamIn}} & \raggedright \leanlink{Formalism.L2.DIV.FamIn} & \raggedright\arraybackslash \leanmod{Formalism.LocalToGlobal.Divisibility.LiftLat} \\
$\pcInfFam$ & \texttt{\detokenize{\pcInfFam}} & \raggedright \leanlink{Formalism.IndefR.Datum.InfFam} & \raggedright\arraybackslash \leanmod{Formalism.LocalToGlobal.IndefiniteReal.LocalNewformFunction} \\
$\pcDegPre$ & \texttt{\detokenize{\pcDegPre}} & \raggedright \leanlink{ReciprocityTheory.ChartRestriction.G.Out.DegPre} & \raggedright\arraybackslash \leanmod{ReciprocityTheory.ReductiveGroups.ArithmeticGeometry.ExtensionDescent.ThickTransComp} \\
$\pcSec$ & \texttt{\detokenize{\pcSec}} & \raggedright \leanlink{ReciprocityTheory.ChartRestriction.G.Semi.LB.Sec} & \raggedright\arraybackslash \leanmod{ReciprocityTheory.ReductiveGroups.ArithmeticGeometry.SemiabelianSchemes.CubicalStructures.CubeSymmetricSwapFirst} \\
$\pcLocWach$ & \texttt{\detokenize{\pcLocWach}} & \raggedright \leanlink{Formalism.BasicNotions.LocWach} & \raggedright\arraybackslash \leanmod{Formalism.SharpRestriction.LevelLetterCocycles.DataNLoc} \\
$\pcRealDat$ & \texttt{\detokenize{\pcRealDat}} & \raggedright \leanlink{Formalism.IndefR.Links.RealDat} & \raggedright\arraybackslash \leanmod{Formalism.LocalToGlobal.IndefiniteReal.LinksDefs} \\
$\pcLinComb$ & \texttt{\detokenize{\pcLinComb}} & \raggedright \leanlink{ReciprocityTheory.ChartRestriction.G.Semi.LB.LinComb} & \raggedright\arraybackslash \leanmod{ReciprocityTheory.ReductiveGroups.ArithmeticGeometry.SemiabelianSchemes.CubicalStructures.CubeSymmetricSwapFirst} \\
$\pcPolAV$ & \texttt{\detokenize{\pcPolAV}} & \raggedright \leanlink{ReciprocityTheory.ChartRestriction.G.Out.PolAV} & \raggedright\arraybackslash \leanmod{ReciprocityTheory.ReductiveGroups.ArithmeticGeometry.ExtensionDescent.NonzeroLocus} \\
$\pcFrame$ & \texttt{\detokenize{\pcFrame}} & \raggedright \leanlink{ReciprocityTheory.ChartRestriction.G.FD.Frame} & \raggedright\arraybackslash \leanmod{ReciprocityTheory.ReductiveGroups.ArithmeticGeometry.FundamentalDomains.Signature} \\
$\pcGen$ & \texttt{\detokenize{\pcGen}} & \raggedright \leanlink{ReciprocityTheory.ChartRestriction.G.Semi.LB.Gen} & \raggedright\arraybackslash \leanmod{ReciprocityTheory.ReductiveGroups.ArithmeticGeometry.SemiabelianSchemes.CubicalStructures.CubeSymmetricSwapFirst} \\
$\pcIsGen$ & \texttt{\detokenize{\pcIsGen}} & \raggedright \leanlink{ReciprocityTheory.ChartRestriction.G.Semi.LB.Gen.IsGen} & \raggedright\arraybackslash \leanmod{ReciprocityTheory.ReductiveGroups.ArithmeticGeometry.SemiabelianSchemes.CubicalStructures.CubeSymmetricSwapFirst} \\
$\pcIsGrpPullback$ & \texttt{\detokenize{\pcIsGrpPullback}} & \raggedright \leanlink{ReciprocityTheory.ChartRestriction.G.Geo.Tor.IsGrpPullback} & \raggedright\arraybackslash \leanmod{ReciprocityTheory.ReductiveGroups.ArithmeticGeometry.BirationalGeometry.ToroidalCompactification.IsGrpPullback} \\
$\pcSpc$ & \texttt{\detokenize{\pcSpc}} & \raggedright \leanlink{GlobalTheory.IndefiniteRowSideFullHecke.WeilTensorModels.Spc} & \raggedright\arraybackslash \leanmod{Development.AbelianVarietyTheory.BlockTheory.DivisorBase.WeilTensorModels} \\
\bottomrule
\end{tabular}
\end{table}

\section{Compression and Deduplication Breakdown}
Reverse-dependency reachability pruning removes 262,467 unreachable declarations (2,104,562~LOC). Within the reachable core, canonical type-signature hashing identifies 72,152 duplicate signature groups accounting for 130,161 redundant declarations (1,958,459~LOC of candidate consolidation).

\section{Formal Verification and Axiom Summary}
Root target declarations audited: \leanlink{LFunction_eq_zero_re}, \leanlink{LFunction_eq_zero_re_of_isPrimitive}, \leanlink{LFunction_ne_zero_of_half_lt_re}, \leanlink{riemannHypothesis}, \leanlink{riemannZeta_eq_zero_re}.
Kernel axiom metadata was not populated in this database snapshot; run the Lean~4 \texttt{.olean} extractor or refactoring gate to record exact \texttt{\#print axioms} output.

\end{document}
